% Elementary Probability — Pure Mathematics with Mechanics Upper Sixth — Cours signé OnBuch+
% Official MINESEC syllabus. Compiler avec : tectonic PMM19-elementary-probability.tex
\newcommand{\DOCMATIERE}{Pure Mathematics with Mechanics}
\newcommand{\DOCNIVEAU}{Upper Sixth}
\newcommand{\DOCMODULE}{Upper Sixth — Pure mathematics with mechanics}
\newcommand{\DOCLECON}{19}
\newcommand{\DOCTITRE}{Elementary Probability}
% OnBuch+ — préambule « cours signés » ENRICHI (compilé avec tectonic / XeLaTeX)
% Système visuel « Pop » d'OnBuch+ : fond crème (jamais blanc), bordures encre
% épaisses, ombres « dures » décalées, titres Archivo Black en capitales.
% Version MATHÉMATIQUES : graphes (pgfplots/tikz), boîtes pédagogiques
% (définition, propriété, méthode, attention, le savais-tu, exercice résolu,
% exercice, TP, bilan), page de garde et signature OnBuch+ ancrée en bas.
%
% Macros attendues dans main.tex AVANT \input{preamble} :
%   \DOCMATIERE  \DOCNIVEAU  \DOCMODULE  \DOCLECON (numéro)  \DOCTITRE
\documentclass[11pt]{article}

\usepackage{fontspec}
\usepackage[english]{babel}
\usepackage[a4paper,margin=2cm,top=3.4cm,bottom=2.8cm,headheight=1.4cm,footskip=1.3cm]{geometry}
\usepackage{amsmath,amssymb,amsfonts}
\usepackage{mathtools} % \xmapsto, \coloneqq... souvent utilisés par les modèles
\usepackage{cancel} % simplifications barrées (bilans de forces)
% \abs{z} (module, valeur absolue) et \norm{u} (norme), utilisés par les modèles
\providecommand{\abs}{}\let\abs\relax\DeclarePairedDelimiter{\abs}{\lvert}{\rvert}
\providecommand{\norm}{}\let\norm\relax\DeclarePairedDelimiter{\norm}{\lVert}{\rVert}
\usepackage[table]{xcolor} % [table] : fournit aussi \rowcolors (alternance de lignes)
\usepackage{pagecolor}
\usepackage{tikz}
\usetikzlibrary{arrows.meta,positioning,calc,shapes.geometric,shapes.misc,shapes.arrows,decorations.pathmorphing,decorations.pathreplacing,decorations.markings,patterns,shadows,mindmap,trees,fit,backgrounds,matrix,intersections,angles,quotes}
\usepackage{pgfplots}
\usepackage[version=4]{mhchem} % équations nucléaires \ce{^{235}_{92}U}
\usepackage[european,siunitx]{circuitikz} % schémas électriques
\pgfplotsset{compat=1.18}
\usepgfplotslibrary{fillbetween,groupplots} % aires entre courbes (calcul intégral)
\usepackage{enumitem}
\usepackage{fancyhdr}
% [most] charge toutes les bibliothèques tcolorbox (listings, minted, raster,
% documentation...) et gonfle inutilement la mémoire TeX sur un document long
% (~300 boîtes) : on ne charge que ce qui est réellement utilisé.
\usepackage[skins,breakable]{tcolorbox}
\usepackage{graphicx}
\usepackage{colortbl}
\usepackage{booktabs}
\usepackage{tabularx}
\usepackage{diagbox} % case d'en-tête diagonale des tableaux à double entrée
% Colonnes extensibles centrée / gauche / droite (C, L, R), souvent utilisées par les modèles
\newcolumntype{C}{>{\centering\arraybackslash}X}
\newcolumntype{L}{>{\raggedright\arraybackslash}X}
\newcolumntype{R}{>{\raggedleft\arraybackslash}X}
\usepackage{array}
\usepackage{multicol}
\usepackage{siunitx}
% parse-numbers=false : accepte aussi \SI{2,5 \times 10^{-3}}{...}
\sisetup{locale=UK,output-decimal-marker={.},group-minimum-digits=4,parse-numbers=false}
\DeclareSIUnit\ohm{\ensuremath{\symup{\Omega}}} % U+2126 absent de la police
\DeclareSIUnit\division{div} % réglages d'oscilloscope (ms/div, V/div)
\DeclareSIUnit\tr{tr} % tours (vitesse de rotation, ex. \SI{1500}{\tr\per\minute})
\usepackage{qrcode}
% \degree : commande fréquemment utilisée par les modèles pour un angle ou une
% température en dehors de \SI{}{} (non fournie par les paquets chargés).
\providecommand{\degree}{^{\circ}}
% \qed (fin de démonstration), utilisé par les modèles sans amsthm
\providecommand{\qed}{\hfill$\square$}
% \barre : double barre des valeurs interdites dans un tableau de variations (tkz-tab)
\providecommand{\barre}{\ensuremath{\|}}
% environnement proof (amsthm non chargé) utilisé par les modèles
\ifdefined\proof\else\newenvironment{proof}[1][Proof]{\par\noindent\textit{#1.}\ }{\qed\par}\fi
\usepackage{lastpage}
\usepackage{hyperref}
\hypersetup{hidelinks}

% ============================================================
% Palette « Pop » OnBuch+ — mêmes hex que la classe P (plus_ui.dart)
% ============================================================
\definecolor{poporange}{HTML}{F59321}
\definecolor{popdark}{HTML}{E07A0C}
\definecolor{popcream}{HTML}{FFF6E8}
\definecolor{popink}{HTML}{1C1714}
\definecolor{poppurple}{HTML}{7A5AE0}
\definecolor{popgreen}{HTML}{2E9E6A}
\definecolor{poppink}{HTML}{E0457B}
\definecolor{popblue}{HTML}{2F7DE1}
\definecolor{popgold}{HTML}{C98A12}
\definecolor{popmuted}{HTML}{8A7F76}
\definecolor{popline}{HTML}{EADFD2}
\definecolor{popgray}{HTML}{8A7F76}
\colorlet{popgrey}{popgray}
% Alias tolérants (le modèle nomme parfois les couleurs différemment)
\colorlet{popwhite}{white}
\colorlet{popblack}{popink}
\colorlet{popred}{poppink}
\colorlet{popyellow}{popgold}
% Teintes claires pour fonds de tableaux / zones
\colorlet{popblueL}{popblue!10!white}
\colorlet{poporangeL}{poporange!14!white}
\colorlet{popgreenL}{popgreen!12!white}
\colorlet{poppurpleL}{poppurple!12!white}
\colorlet{poppinkL}{poppink!12!white}
\colorlet{popgoldL}{popgold!14!white}

\pagecolor{popcream}

% ============================================================
% Typographie
% ============================================================
\defaultfontfeatures{Ligatures=TeX}
\setmainfont{PlusJakartaSans-Medium.ttf}[Path=../../../fonts/,BoldFont=PlusJakartaSans-Bold.ttf]
% Maths en sans-serif (Fira Math) avec chiffres et lettres droites de Plus
% Jakarta, pour que les formules se fondent dans le texte.
\usepackage[math-style=ISO,bold-style=ISO]{unicode-math}
\setmathfont{FiraMath-Regular.otf}[Path=../../../fonts/]
\setmathfont{PlusJakartaSans-Medium.ttf}[Path=../../../fonts/,range={up/num,up/{Latin,latin},\mathup}]
% (icomma retiré : décimales avec point en anglais)
% Caractères mathématiques Unicode absents des polices de texte (Plus Jakarta,
% Archivo Black) : rendus par leur équivalent mathématique, en texte comme en
% formule — sinon ils s'affichent en carré vide (ex. « Arithmétique dans ℤ »).
\usepackage{newunicodechar}
\newunicodechar{ℝ}{\ensuremath{\mathbb{R}}}
\newunicodechar{ℤ}{\ensuremath{\mathbb{Z}}}
\newunicodechar{ℂ}{\ensuremath{\mathbb{C}}}
\newunicodechar{ℕ}{\ensuremath{\mathbb{N}}}
\newunicodechar{ℚ}{\ensuremath{\mathbb{Q}}}
\newunicodechar{𝔻}{\ensuremath{\mathbb{D}}}
\newunicodechar{α}{\ensuremath{\alpha}}
\newunicodechar{β}{\ensuremath{\beta}}
\newunicodechar{γ}{\ensuremath{\gamma}}
\newunicodechar{δ}{\ensuremath{\delta}}
\newunicodechar{ε}{\ensuremath{\varepsilon}}
\newunicodechar{θ}{\ensuremath{\theta}}
\newunicodechar{λ}{\ensuremath{\lambda}}
\newunicodechar{μ}{\ensuremath{\mu}}
\newunicodechar{σ}{\ensuremath{\sigma}}
\newunicodechar{τ}{\ensuremath{\tau}}
\newunicodechar{φ}{\ensuremath{\varphi}}
\newunicodechar{ψ}{\ensuremath{\psi}}
\newunicodechar{ω}{\ensuremath{\omega}}
\newunicodechar{Σ}{\ensuremath{\Sigma}}
% majuscules produites par \MakeUppercase dans les titres (α-aminés -> Α)
\newunicodechar{Α}{\ensuremath{\alpha}}
\newunicodechar{Β}{\ensuremath{\beta}}
\newunicodechar{Γ}{\ensuremath{\Gamma}}
\newunicodechar{Δ}{\ensuremath{\Delta}}
\newunicodechar{Ε}{\ensuremath{\varepsilon}}
\newunicodechar{Θ}{\ensuremath{\Theta}}
\newunicodechar{Λ}{\ensuremath{\Lambda}}
\newunicodechar{Μ}{\ensuremath{\mu}}
\newunicodechar{Τ}{\ensuremath{\tau}}
\newunicodechar{Φ}{\ensuremath{\Phi}}
\newunicodechar{Ψ}{\ensuremath{\Psi}}
\newunicodechar{Ω}{\ensuremath{\Omega}}
\newunicodechar{↦}{\ensuremath{\mapsto}}
\newunicodechar{∈}{\ensuremath{\in}}
\newunicodechar{∉}{\ensuremath{\notin}}
\newunicodechar{∧}{\ensuremath{\wedge}}
\newunicodechar{⇔}{\ensuremath{\Leftrightarrow}}
\newunicodechar{⟺}{\ensuremath{\Longleftrightarrow}}
\newunicodechar{⇒}{\ensuremath{\Rightarrow}}
\newunicodechar{⟹}{\ensuremath{\Longrightarrow}}
\newunicodechar{∘}{\ensuremath{\circ}}
\newunicodechar{∪}{\ensuremath{\cup}}
\newunicodechar{∩}{\ensuremath{\cap}}
\newunicodechar{≡}{\ensuremath{\equiv}}
\newunicodechar{⊂}{\ensuremath{\subset}}
\newunicodechar{⊥}{\ensuremath{\perp}}
\newunicodechar{∃}{\ensuremath{\exists}}
\newunicodechar{∀}{\ensuremath{\forall}}
\newunicodechar{∛}{\ensuremath{\sqrt[3]{\,}}}
\newunicodechar{∜}{\ensuremath{\sqrt[4]{\,}}}
\newunicodechar{⃗}{\ensuremath{{}^{\to}}}
\newunicodechar{ⁿ}{\ifmmode^{n}\else\textsuperscript{\ensuremath{n}}\fi}
\newunicodechar{ˣ}{\ifmmode^{x}\else\textsuperscript{\ensuremath{x}}\fi}
\newunicodechar{ᵇ}{\ifmmode^{b}\else\textsuperscript{\ensuremath{b}}\fi}
\newunicodechar{ʳ}{\ifmmode^{r}\else\textsuperscript{\ensuremath{r}}\fi}
\newunicodechar{ᵘ}{\ifmmode^{u}\else\textsuperscript{\ensuremath{u}}\fi}
\newunicodechar{ʸ}{\ifmmode^{y}\else\textsuperscript{\ensuremath{y}}\fi}
\newunicodechar{ᵃ}{\ifmmode^{a}\else\textsuperscript{\ensuremath{a}}\fi}
\newunicodechar{ᵉ}{\ifmmode^{e}\else\textsuperscript{\ensuremath{e}}\fi}
\newunicodechar{ᵏ}{\ifmmode^{k}\else\textsuperscript{\ensuremath{k}}\fi}
\newunicodechar{ᵖ}{\ifmmode^{p}\else\textsuperscript{\ensuremath{p}}\fi}
\newunicodechar{ᴺ}{\ifmmode^{N}\else\textsuperscript{\ensuremath{N}}\fi}
\newunicodechar{ᶜ}{\ifmmode^{c}\else\textsuperscript{\ensuremath{c}}\fi}
\newunicodechar{ʰ}{\ifmmode^{h}\else\textsuperscript{\ensuremath{h}}\fi}
\newunicodechar{ᵠ}{\ifmmode^{q}\else\textsuperscript{\ensuremath{q}}\fi}
\newunicodechar{ₙ}{\ifmmode_{n}\else\textsubscript{\ensuremath{n}}\fi}
\newunicodechar{ₐ}{\ifmmode_{a}\else\textsubscript{\ensuremath{a}}\fi}
\newunicodechar{ᵢ}{\ifmmode_{i}\else\textsubscript{\ensuremath{i}}\fi}
\newunicodechar{ᵣ}{\ifmmode_{r}\else\textsubscript{\ensuremath{r}}\fi}
\newunicodechar{ₖ}{\ifmmode_{k}\else\textsubscript{\ensuremath{k}}\fi}
\newunicodechar{ᵦ}{\ifmmode_{\beta}\else\textsubscript{\ensuremath{\beta}}\fi}
\newunicodechar{ₓ}{\ifmmode_{x}\else\textsubscript{\ensuremath{x}}\fi}
\newfontfamily\popdisplay{ArchivoBlack-Regular.ttf}[Path=../../../fonts/]
\newfontfamily\poplogo{SpaceGrotesk-Bold.ttf}[Path=../../../fonts/]
\newfontfamily\popsemi{PlusJakartaSans-SemiBold.ttf}[Path=../../../fonts/]
\newfontfamily\popxbold{PlusJakartaSans-ExtraBold.ttf}[Path=../../../fonts/]
\newfontfamily\popxboldls{PlusJakartaSans-ExtraBold.ttf}[Path=../../../fonts/,LetterSpace=3.0]

\setlist[enumerate]{leftmargin=1.7em,itemsep=5pt,topsep=4pt}
\setlist[itemize]{leftmargin=1.7em,itemsep=4pt,topsep=4pt,label={\color{poporange}\textbullet}}
\setlength{\parindent}{0pt}
\setlength{\parskip}{5pt}
\renewcommand{\arraystretch}{1.25}

% pgfplots : style Pop par défaut
\pgfplotsset{
  every axis/.append style={axis line style={popink,line width=1pt},tick style={popink},
    grid style={popline,line width=0.5pt},label style={font=\small},tick label style={font=\footnotesize}},
  popcurve/.style={poporange,line width=1.8pt,no markers,smooth},
}
% Même style disponible dans un \draw TikZ simple (hors pgfplots)
\tikzset{popcurve/.style={poporange,line width=1.8pt}}
\tikzset{popgrid/.style={popline,very thin}}
\colorlet{popcurve}{poporange} % le nom du style est parfois utilisé comme couleur

% ============================================================
% Pill / squiggle
% ============================================================
\newcommand{\poppill}[3]{%
  \tikz[baseline=-0.55ex]\node[draw=popink,line width=1.2pt,rounded corners=2.6pt,%
    fill=#2,inner xsep=6.5pt,inner ysep=3pt]{%
    \popxboldls\fontsize{7}{7}\selectfont\color{#3}\MakeUppercase{#1}};%
}
\newcommand{\popsquiggle}[1]{%
  \tikz[baseline=-0.4ex]{
    \draw[#1,line width=2.6pt,line cap=round]
      plot [smooth] coordinates {(0,0.09) (0.34,0.01) (0.68,0.21) (1.02,0.01) (1.36,0.10)};
  }%
}
% Mot-clé surligné (marqueur orange sous le texte)
\newcommand{\cle}[1]{{\popxbold\color{popdark}#1}}

% ============================================================
% En-tête / pied
% ============================================================
\pagestyle{fancy}
\fancyhf{}
\renewcommand{\headrulewidth}{0pt}
\renewcommand{\footrulewidth}{0pt}
\lhead{\raisebox{-0.15ex}{\poplogo\fontsize{13}{13}\selectfont\color{popink}OnBuch\color{poporange}+}}
\rhead{\poppill{\DOCMATIERE\ \textbullet\ \DOCNIVEAU\ \textbullet\ Chapter \DOCLECON}{white}{popink}}
\lfoot{\popxbold\fontsize{7.6}{7.6}\selectfont\color{popmuted}\MakeUppercase{Signed course OnBuch+}\ \textbullet\ {\popsemi\DOCTITRE}}
\rfoot{\tikz[baseline=-0.6ex]\node[rounded corners=8.5pt,fill=popink,minimum height=17pt,minimum width=17pt,inner xsep=5pt,inner sep=0]{\popxbold\fontsize{8}{8}\selectfont\color{popcream}\thepage\,/\,\pageref*{LastPage}};}

% ============================================================
% Titres
% ============================================================
\newcommand{\chaptitre}[2]{%
  \begin{tcolorbox}[enhanced,colback=poporange,colframe=popink,boxrule=2.6pt,arc=11pt,
      drop shadow=popink,left=15pt,right=15pt,top=12pt,bottom=11pt,before skip=3pt,after skip=18pt]
    \poppill{#2}{popink}{popcream}\par\vspace{9pt}
    {\raggedright\hyphenpenalty=10000\exhyphenpenalty=10000\popdisplay\fontsize{22}{24}\selectfont\color{popcream}\MakeUppercase{#1}\par}
  \end{tcolorbox}
}
% Section numérotée : pastille orange avec le numéro + titre Archivo
\newcounter{popsec}
\newcommand{\coursec}[1]{%
  \stepcounter{popsec}\setcounter{popsub}{0}%
  \par\vspace{16pt}\noindent
  \begin{minipage}{\linewidth}
  \tikz[baseline=-0.7ex]\node[draw=popink,line width=1.6pt,fill=poporange,rounded corners=4pt,minimum size=22pt,inner sep=0]{\popdisplay\fontsize{12}{12}\selectfont\color{popcream}\thepopsec};%
  \hspace{8pt}{\popdisplay\fontsize{15}{17}\selectfont\color{popink}\MakeUppercase{#1}}\par
  \vspace{4pt}\noindent{\color{poporange}\rule{36pt}{3.6pt}}
  \end{minipage}\par\nopagebreak\vspace{8pt}
}
\newcounter{popsub}
\newcommand{\courssub}[1]{%
  \stepcounter{popsub}%
  \par\vspace{9pt}\noindent{\popxbold\fontsize{11.5}{13}\selectfont\color{popink}{\color{poporange}\thepopsec.\thepopsub}\hspace{6pt}#1}\par\nopagebreak\vspace{4pt}
}

% ============================================================
% Boîtes pédagogiques — même recette : fond blanc, bordure épaisse colorée,
% ombre dure encre, badge plein. Titre optionnel [#1].
% ============================================================
\tcbset{popbase/.style={breakable,enhanced,colback=white,boxrule=2.3pt,arc=9pt,
  shadow={3pt}{-3pt}{0pt}{popink},left=14pt,right=14pt,top=11pt,bottom=11pt,before skip=11pt,after skip=15pt,
  fontupper=\popsemi\fontsize{10.6}{14.5}\selectfont\color{popink},
  fontlower=\popsemi\fontsize{10.6}{14.5}\selectfont\color{popink}}}
% Coche dessinée (le glyphe ✓ n'existe pas dans les polices du document)
\newcommand{\popcheck}[1]{\tikz[baseline=-0.5ex]\draw[#1,line width=1.6pt,line cap=round,line join=round](-0.13,0.0)--(-0.04,-0.09)--(0.13,0.1);}
\providecommand{\coche}{\popcheck{popgreen}}
\providecommand{\resultat}[1]{\par\smallskip\noindent\fcolorbox{popgreen}{popgreenL}{\parbox{\dimexpr\linewidth-2\fboxsep-2\fboxrule\relax}{\textbf{Result.} #1}}\par\smallskip}
\newcommand{\popboxhead}[3]{\poppill{#1}{#2}{white}\ifx\relax#3\relax\else\hspace{6pt}{\popxbold\fontsize{10.6}{12}\selectfont\color{popink}#3}\fi\par\vspace{7pt}}

\newtcolorbox{aretenir}[1][]{popbase,colframe=popgreen,before upper={\popboxhead{Key points}{popgreen}{#1}}}
\newtcolorbox{exemplebox}[1][]{popbase,colframe=poporange,before upper={\popboxhead{Exemple}{poporange}{#1}}}
\newtcolorbox{definition}[1][]{popbase,colframe=popblue,before upper={\popboxhead{Definition}{popblue}{#1}}}
\newtcolorbox{propriete}[1][]{popbase,colframe=poppurple,before upper={\popboxhead{Property}{poppurple}{#1}}}
\newtcolorbox{methode}[1][]{popbase,colframe=popgold,before upper={\popboxhead{Method}{popgold}{#1}}}
\newtcolorbox{attention}[1][]{popbase,colframe=poppink,before upper={\popboxhead{Warning}{poppink}{#1}}}
\newtcolorbox{experience}[1][]{popbase,colframe=popblue,colback=popblueL,before upper={\popboxhead{Experiment}{popblue}{#1}}}
% « Le savais-tu ? » : carte violette pleine (contexte, culture, Cameroun)
\newtcolorbox{savaistu}[1][]{popbase,colback=poppurple,colframe=popink,
  fontupper=\popsemi\fontsize{10.4}{14.2}\selectfont\color{white},
  before upper={\poppill{Did you know?}{white}{poppurple}\ifx\relax#1\relax\else\hspace{6pt}{\popxbold\color{white}#1}\fi\par\vspace{7pt}}}
% Exercice résolu : énoncé en haut, \tcblower puis la solution sur fond crème
\newcounter{popexo}\newcounter{popexr}
\newtcolorbox{exoresolu}[1][]{popbase,colframe=poporange,colbacklower=popcream,
  segmentation style={popink,line width=1.2pt,dashed},
  before upper={\stepcounter{popexr}\popboxhead{Worked example \thepopexr}{poporange}{#1}},
  before lower={\poppill{Solution}{popink}{popcream}\par\vspace{6pt}}}
% Exercice (non résolu en place) — la correction est dans la partie Corrigés
\newtcolorbox{exercice}[1][]{popbase,colframe=popink,
  before upper={\stepcounter{popexo}\popboxhead{Exercise \thepopexo}{popink}{#1}}}
% Corrigé
\newtcolorbox{corrige}[1][]{popbase,colframe=popgreen,colback=popgreenL,
  before upper={\popboxhead{Solution}{popgreen}{#1}}}
% Objectifs / prérequis / bilan
\newtcolorbox{objectifs}{popbase,colframe=popink,colback=poporangeL,
  before upper={\popboxhead{Chapter objectives}{popink}{}}}
\newtcolorbox{prerequis}{popbase,colframe=popink,colback=popblueL,
  before upper={\popboxhead{Prerequisites}{popblue}{}}}
\newtcolorbox{bilan}{popbase,colframe=popink,colback=white,boxrule=2.8pt,
  before upper={\popboxhead{Fiche bilan}{poporange}{L'essentiel à connaître pour l'examen}}}
% Figure encadrée avec légende
\newtcolorbox{popfigure}[1][]{enhanced,breakable=false,colback=white,colframe=popline,boxrule=1.4pt,arc=8pt,
  left=10pt,right=10pt,top=10pt,bottom=8pt,before skip=10pt,after skip=12pt,halign=center,
  lower separated=false,halign lower=center,
  fontlower=\popsemi\fontsize{9}{11.5}\selectfont\color{popmuted},
  #1}
\newcommand{\legende}[1]{\par\vspace{6pt}{\popsemi\fontsize{9}{11.5}\selectfont\color{popmuted}#1}}

% ============================================================
% Signature OnBuch+
% ============================================================
\newcommand{\poplogoinline}[1]{{\poplogo\fontsize{#1}{#1}\selectfont\color{popink}OnBuch\color{poporange}+}}
\newcommand{\signatureonbuch}{%
  \begin{tcolorbox}[enhanced,colback=popink,colframe=popink,boxrule=2.2pt,arc=10pt,
      drop shadow=poporange,left=14pt,right=14pt,top=10pt,bottom=10pt,before skip=0pt,after skip=0pt]
    \tikz[baseline=-0.7ex]\node[circle,fill=popgreen,draw=popcream,line width=1.3pt,minimum size=22pt,inner sep=0]{\popcheck{white}};%
    \hspace{9pt}\begin{minipage}[c]{0.62\linewidth}
      {\popxbold\fontsize{12}{13}\selectfont\color{popcream}Signed\ \textbullet\ {\poplogo OnBuch\color{poporange}+}}\par\vspace{2pt}
      {\raggedright\popsemi\fontsize{8}{10.5}\selectfont\color{popline}\MakeUppercase{OnBuch+ teaching team}\\\MakeUppercase{Follows the official MINESEC syllabus}\par}
    \end{minipage}\hfill
    {\poplogo\fontsize{22}{22}\selectfont\color{popcream}OnBuch\color{poporange}+}
  \end{tcolorbox}%
}

% ============================================================
% Page de garde
% #1 titre · #2 matière · #3 niveau · #4 module · #5 n° leçon · #6 durée ·
% #7 description / objectifs (texte ou itemize)
% ============================================================
\newcommand{\pagegarde}[7]{%
  \thispagestyle{empty}
  \newgeometry{a4paper,margin=1.7cm,top=1.6cm,bottom=1.4cm}%
  \begin{tikzpicture}[remember picture,overlay]
    \fill[poporange!18] ([xshift=-3.2cm,yshift=-3.6cm]current page.north east) circle (4.6cm);
    \fill[poppurple!14] ([xshift=2.4cm,yshift=7.6cm]current page.south west) circle (3.8cm);
  \end{tikzpicture}%
  {\poplogo\fontsize{30}{30}\selectfont\color{popink}OnBuch\color{poporange}+}\hfill
  \raisebox{0.6ex}{\poppill{Signed course \textbullet\ Premium}{popink}{popcream}}\par
  \vspace{1pt}\popsquiggle{poppurple}\par
  \vspace{8pt}
  \begin{tcolorbox}[enhanced,colback=poporange,colframe=popink,boxrule=2.8pt,arc=13pt,
      drop shadow=popink,left=18pt,right=18pt,top=12pt,bottom=13pt,after skip=10pt]
    \poppill{#2 \textbullet\ #3}{popink}{popcream}\hspace{5pt}\poppill{#4}{white}{popink}\par\vspace{8pt}
    {\popdisplay\fontsize{14}{14}\selectfont\color{popink}CHAPTER #5}\par\vspace{4pt}
    {\raggedright\hyphenpenalty=10000\exhyphenpenalty=10000\popdisplay\fontsize{25}{27}\selectfont\color{popcream}\MakeUppercase{#1}\par}
    \vspace{8pt}
    {\popxbold\fontsize{9.5}{9.5}\selectfont\color{popink}\MakeUppercase{Suggested time: #6}}
  \end{tcolorbox}
  \begin{tcolorbox}[enhanced,colback=white,colframe=popink,boxrule=2.2pt,arc=10pt,
      drop shadow=popink,left=16pt,right=16pt,top=10pt,bottom=10pt,after skip=8pt]
    \poppill{In this chapter}{poppurple}{white}\par\vspace{5pt}
    \popsemi\fontsize{9.3}{12.6}\selectfont\color{popink}#7
  \end{tcolorbox}
  \noindent\begin{minipage}[t]{0.487\linewidth}
    \begin{tcolorbox}[enhanced,colback=popgreen,colframe=popink,boxrule=2.2pt,arc=10pt,
        drop shadow=popink,left=11pt,right=11pt,top=10pt,bottom=10pt,height=3.1cm,valign=center]
      \tcbox[colback=white,colframe=popink,boxrule=1.3pt,arc=5pt,boxsep=0pt,nobeforeafter,
        left=3pt,right=3pt,top=3pt,bottom=3pt,valign=center]{\qrcode[height=1.9cm]{https://play.google.com/store/apps/details?id=com.luvvix.onbuch&hl=en}}%
      \hspace{8pt}\begin{minipage}[c]{0.5\linewidth}
        {\popdisplay\fontsize{11}{12.5}\selectfont\color{white}\MakeUppercase{Download OnBuch}}\par\vspace{4pt}
        {\raggedright\popsemi\fontsize{8.4}{11}\selectfont\color{white}Free \textbullet\ Google Play\\AI tutor, past papers, results\par}
      \end{minipage}
    \end{tcolorbox}
  \end{minipage}\hfill
  \begin{minipage}[t]{0.487\linewidth}
    \begin{tcolorbox}[enhanced,colback=poppurple,colframe=popink,boxrule=2.2pt,arc=10pt,
        drop shadow=popink,left=11pt,right=11pt,top=10pt,bottom=10pt,height=3.1cm,valign=center]
      \tikz[baseline=-0.7ex]\node[circle,fill=white,draw=popink,line width=1.3pt,minimum size=1.6cm,inner sep=0]{\popdisplay\fontsize{24}{24}\selectfont\color{poppurple}+};%
      \hspace{8pt}\begin{minipage}[c]{0.6\linewidth}
        {\popdisplay\fontsize{11}{12.5}\selectfont\color{white}\MakeUppercase{Go OnBuch+}}\par\vspace{4pt}
        {\raggedright\popsemi\fontsize{8.4}{11}\selectfont\color{white}All signed courses, unlimited AI tutor, PDF sheets — OnBuch+ tab in the app\par}
      \end{minipage}
    \end{tcolorbox}
  \end{minipage}
  \vspace{8pt}
  \popcontenu
  \vfill
  \signatureonbuch
  \restoregeometry
  \newpage
}

% Rangée « Ce cours contient » (page de garde)
\newcommand{\poptile}[3]{% #1 couleur #2 gros texte #3 légende
  \begin{tcolorbox}[enhanced,colback=#1,colframe=popink,boxrule=2pt,arc=9pt,shadow={2.5pt}{-2.5pt}{0pt}{popink},
      left=6pt,right=6pt,top=9pt,bottom=9pt,halign=center,valign=center,height=1.75cm,width=\linewidth,nobeforeafter]
    {\popdisplay\fontsize{12.5}{13}\selectfont\color{white}\MakeUppercase{#2}}\par\vspace{3pt}
    {\centering\popsemi\fontsize{7.8}{9.5}\selectfont\color{white}#3\par}
  \end{tcolorbox}}
\newcommand{\popcontenu}{%
  \par\noindent{\popxboldls\fontsize{7.5}{7.5}\selectfont\color{popmuted}THIS SIGNED COURSE CONTAINS}\par\vspace{7pt}
  \noindent\begin{minipage}[t]{0.235\linewidth}\poptile{popblue}{Course}{complete, illustrated, step by step}\end{minipage}\hfill
  \begin{minipage}[t]{0.235\linewidth}\poptile{poporange}{Methods}{and worked examples}\end{minipage}\hfill
  \begin{minipage}[t]{0.235\linewidth}\poptile{poppink}{Exercises}{exam style + solutions}\end{minipage}\hfill
  \begin{minipage}[t]{0.235\linewidth}\poptile{popgreen}{Summary sheet}{the essentials to remember}\end{minipage}\par}

% Sommaire
\newtcolorbox{sommaire}{enhanced,colback=white,colframe=popink,boxrule=2.2pt,arc=10pt,
  drop shadow=popink,left=18pt,right=18pt,top=13pt,bottom=13pt,before skip=6pt,after skip=20pt,
  before upper={\poppill{Contents}{poppurple}{white}\par\vspace{9pt}\popsemi\fontsize{10.6}{14.5}\selectfont\color{popink}}}

% Signature de fin de cours — poussée tout en bas de la dernière page
\newcommand{\findecours}{%
  \par\vfill\nopagebreak
  \begin{center}\popsquiggle{poporange}\end{center}\vspace{4pt}
  \signatureonbuch
}
\AtBeginDocument{\def\llbracket{\mathopen{[\mkern-3.5mu[}}\def\rrbracket{\mathclose{]\mkern-3.5mu]}}} % intervalles d'entiers (glyphe ⟦ absent de la police)
% Glyphes absents de Fira Math : redéfinis à partir de glyphes disponibles
\AtBeginDocument{%
  \def\setminus{\mathbin{\backslash}}\let\smallsetminus\setminus
  \def\vdots{\mathord{\vbox{\baselineskip3.2pt\lineskiplimit0pt\kern4pt\hbox{.}\hbox{.}\hbox{.}}}}%
  \def\ddots{\mathinner{\mkern1mu\raise6.5pt\hbox{.}\mkern2mu\raise3.6pt\hbox{.}\mkern2mu\raise0.7pt\hbox{.}\mkern1mu}}%
  \def\triangle{\mathord{\scalebox{0.9}{$\symup{\Delta}$}}}%
  \def\nabla{\mathord{\raisebox{\depth}{\scalebox{1}[-1]{$\symup{\Delta}$}}}}%
  \def\checkmark{\popcheck{popdark}}%
  \def\star{\mathbin{\ast}}\def\bigstar{\mathord{\ast}}%
  \def\diamond{\mathbin{\scalebox{0.6}{$\Diamond$}}}%
  \def\bigcup{\mathop{\mathchoice{\raisebox{-0.3em}{\scalebox{1.7}{$\cup$}}}{\scalebox{1.25}{$\cup$}}{\cup}{\cup}}\displaylimits}%
  \def\bigcap{\mathop{\mathchoice{\raisebox{-0.3em}{\scalebox{1.7}{$\cap$}}}{\scalebox{1.25}{$\cap$}}{\cap}{\cap}}\displaylimits}%
  \def\lesseqqgtr{\mathrel{\lessgtr}}\def\bigoplus{\mathop{\oplus}\displaylimits}%
}
\providecommand{\ssi}{if and only if} % abréviation fréquente
\AtBeginDocument{\providecommand{\wideparen}{\overparen}} % arc de cercle

% Fonctions usuelles que les modèles utilisent sans que amsmath les fournisse
\providecommand{\arsinh}{\operatorname{arsinh}}\providecommand{\arcosh}{\operatorname{arcosh}}
\providecommand{\artanh}{\operatorname{artanh}}\providecommand{\arsech}{\operatorname{arsech}}
\providecommand{\arcsch}{\operatorname{arcsch}}\providecommand{\arcoth}{\operatorname{arcoth}}
\providecommand{\arcsinh}{\operatorname{arsinh}}\providecommand{\arccosh}{\operatorname{arcosh}}
\providecommand{\arctanh}{\operatorname{artanh}}\providecommand{\sech}{\operatorname{sech}}
\providecommand{\csch}{\operatorname{csch}}\providecommand{\arcsec}{\operatorname{arcsec}}
\providecommand{\arccsc}{\operatorname{arccsc}}\providecommand{\arccot}{\operatorname{arccot}}
\providecommand{\sgn}{\operatorname{sgn}}\providecommand{\lcm}{\operatorname{lcm}}
\providecommand{\Var}{\operatorname{Var}}\providecommand{\Cov}{\operatorname{Cov}}

%% FIGFIX-BEGIN v3 — correctifs globaux des figures (docs/onbuch-generation/tools/figfix)
\usepackage{adjustbox}\usepackage{etoolbox}
% toute figure TikZ/pgfplots plus large que la ligne est réduite à la largeur disponible
\BeforeBeginEnvironment{tikzpicture}{\begin{adjustbox}{max width=\linewidth}}
\AfterEndEnvironment{tikzpicture}{\end{adjustbox}}
% légendes pgfplots lisibles par-dessus les courbes
\pgfplotsset{every axis/.append style={legend style={fill=white,fill opacity=0.92,text opacity=1,draw=black!15,inner sep=3pt}}}
% étiquettes d'arêtes (syntaxe edge["..."]) avec fond blanc
\tikzset{every edge quotes/.append style={fill=white,inner sep=1.2pt}}
% bulles de cartes mentales : texte à la ligne dans la bulle, bulle un peu plus grande
\tikzset{every concept/.append style={text width=2.0cm,align=center,minimum size=2.3cm,inner sep=2pt}}
%% FIGFIX-END
\begin{document}

\pagegarde{Elementary Probability}{Pure Mathematics with Mechanics}{Upper Sixth}{Upper Sixth — Pure mathematics with mechanics}{19}{16 periods}{This chapter introduces the language, laws and models of elementary probability, from simple experiments with coins, dice and cards to conditional probability, tree diagrams and Bayes' theorem. Learners will compute theoretical and experimental probabilities, use Venn diagrams and counting techniques, and apply probability to authentic Cameroonian and examination situations.
\vspace{4pt}
\begin{multicols}{2}\small
\begin{itemize}
\item By the end of this chapter I can define the basic vocabulary of probability (experiment, outcome, sample space, event).
\item By the end of this chapter I can list the sample space of one-stage and two-stage experiments systematically.
\item By the end of this chapter I can compute classical probabilities P(E) = n(E)/n(S) for equally likely outcomes.
\item By the end of this chapter I can estimate probabilities empirically from relative frequency and relate them to theoretical values.
\item By the end of this chapter I can apply the addition law for mutually exclusive and non-exclusive events.
\item By the end of this chapter I can apply the multiplication law for independent and dependent events, including conditional probability.
\end{itemize}\end{multicols}}

\begin{sommaire}
\begin{multicols}{2}
\begin{enumerate}
\item Foundations: Experiments, Outcomes, Sample Spaces and Events
\item Classical and Empirical Probability
\item The Addition Law and Venn Diagrams
\item Multiplication Law, Conditional Probability and Independence
\item Tree Diagrams and Multi-Stage Experiments
\item Probability with Cards, Dice, Coins and Counting Techniques
\item Bayes' Theorem
\item Applications, Synthesis and Exam Preparation
\item Integration activity
\item Methods and worked examples
\item Exercises
\item Solutions to the exercises
\item Summary sheet
\end{enumerate}
\end{multicols}
\end{sommaire}

\begin{objectifs}
\begin{itemize}
\item By the end of this chapter I can define the basic vocabulary of probability (experiment, outcome, sample space, event).
\item By the end of this chapter I can list the sample space of one-stage and two-stage experiments systematically.
\item By the end of this chapter I can compute classical probabilities P(E) = n(E)/n(S) for equally likely outcomes.
\item By the end of this chapter I can estimate probabilities empirically from relative frequency and relate them to theoretical values.
\item By the end of this chapter I can apply the addition law for mutually exclusive and non-exclusive events.
\item By the end of this chapter I can apply the multiplication law for independent and dependent events, including conditional probability.
\item By the end of this chapter I can represent and solve probability problems using Venn diagrams and tree diagrams.
\item By the end of this chapter I can use permutations and combinations to count outcomes in probability problems.
\item By the end of this chapter I can state and apply Bayes' theorem to reverse conditional probabilities.
\item By the end of this chapter I can model and solve real-life probability problems (games, quality control, medical testing) in GCE-style questions.
\end{itemize}
\end{objectifs}

\begin{prerequis}
\begin{itemize}
\item Set language and notation: union, intersection, complement, Venn diagrams (Lower Sixth).
\item Counting techniques: factorial notation, permutations and combinations (Lower Sixth).
\item Operations on fractions, decimals and percentages; simplification of ratios.
\item Basic algebra: solving simple linear equations and manipulating fractions with unknowns.
\item Simple statistics: frequency tables and relative frequency (Form 5 / Lower Sixth).
\end{itemize}
\end{prerequis}

\newpage

\coursec{Foundations: Experiments, Outcomes, Sample Spaces and Events}

Every rainy season in Douala, commuters bet on whether the 5 p.m.\ downpour will strike before they reach the motor park; at the same time, a GCE candidate wonders what the chance is of drawing a ``red'' question on probability in the June paper, and a mobile-money agent in Bamenda estimates the likelihood that the next customer will withdraw rather than deposit. All three are, knowingly or not, doing probability. From predicting football pools to assessing the risk of a failed harvest, probability is the mathematics of uncertainty --- and this chapter gives you the exact tools the GCE Advanced Level expects you to master.

Before we can measure uncertainty, we must describe it precisely. What exactly is ``an experiment''? What do we mean by ``an outcome''? How do we list \emph{everything} that could happen, so that nothing is forgotten and nothing is counted twice? This first section builds the vocabulary and the systematic habits on which the whole chapter rests. Master it, and the laws of probability that follow will feel natural; neglect it, and every later formula will rest on sand.

\courssub{Random Experiments, Trials and Outcomes}

\begin{experience}[A first observation]
Take an ordinary coin and toss it once. Before it lands, nobody --- not even the greatest mathematician --- can say with certainty whether it will show Head or Tail. Yet two things are perfectly clear: the \emph{list} of possible results is known in advance (Head or Tail), and if you repeat the toss many times, a regular pattern emerges (each side appears roughly half of the time). This mixture of individual unpredictability and long-run regularity is exactly what probability theory studies.
\end{experience}

\begin{definition}[Random experiment, trial, outcome]
A \cle{random experiment} (or \emph{probabilistic experiment}) is a process or action that satisfies two conditions:
\begin{enumerate}
\item Its result cannot be predicted with certainty before it is performed.
\item The set of all possible results is known in advance.
\end{enumerate}
Each single performance of the experiment is called a \cle{trial}, and each possible result of a trial is called an \cle{outcome} (or \emph{elementary outcome}).
\end{definition}

\begin{definition}[Deterministic experiment]
A \cle{deterministic experiment} is a process whose result is uniquely determined by the initial conditions (e.g.\ calculating $2+3$, or measuring the boiling point of pure water at standard pressure). Probability theory does \emph{not} deal with deterministic experiments.
\end{definition}

\begin{exemplebox}[Identifying experiments, trials and outcomes]
\begin{itemize}
\item Tossing a coin once is a random experiment; the outcomes are $H$ (head) and $T$ (tail).
\item Rolling a fair die is a random experiment; the outcomes are $1, 2, 3, 4, 5, 6$.
\item Drawing one card from a well-shuffled pack of $52$ cards is a random experiment with $52$ outcomes.
\item Choosing one delegate at random from a class of $40$ students is a random experiment with $40$ outcomes.
\end{itemize}
By contrast, ``will it rain in Douala at 5 p.m.\ tomorrow?'' is \emph{not} a repeatable random experiment in the strict sense --- but we shall see later how empirical probability still lets us treat such questions.
\end{exemplebox}

\courssub{The Sample Space}

\begin{definition}[Sample space]
The \cle{sample space} of a random experiment, denoted $S$ (some books write $\Omega$ or $\mathcal{E}$), is the set of \emph{all} possible outcomes of the experiment. The outcomes in $S$ must be:
\begin{itemize}
\item \textbf{Mutually exclusive:} no two outcomes can occur simultaneously in a single trial.
\item \textbf{Exhaustive:} one of the outcomes in $S$ \emph{must} occur in every trial.
\end{itemize}
The number of outcomes in $S$ is written $n(S)$.
\end{definition}

The sample space is the ``universe'' of the experiment: every question we ask about chance will be a question about a part of $S$. Writing $S$ correctly and completely is the first skill of the probabilist.

\begin{methode}[Systematic listing of a sample space]
\begin{enumerate}
\item Identify clearly what one trial of the experiment consists of (one coin? two dice? a coin \emph{and} a die?).
\item Decide how to record an outcome (letters, numbers, ordered pairs).
\item Enumerate \emph{systematically}: fix the first object and vary the second, then change the first, and so on. For two-stage experiments, use a list of ordered pairs, a two-way table, or a grid.
\item Count the outcomes to check: if stage 1 has $m$ results and stage 2 has $n$ results, then $n(S) = m \times n$ (the \cle{product rule}).
\end{enumerate}
\end{methode}

\begin{exemplebox}[One-stage sample spaces]
\begin{itemize}
\item One coin: $S = \{H, T\}$, $n(S) = 2$.
\item One die: $S = \{1, 2, 3, 4, 5, 6\}$, $n(S) = 6$.
\item One card from a pack: $S$ has $n(S) = 52$ outcomes ($13$ values in each of $4$ suits).
\end{itemize}
\end{exemplebox}

\begin{exemplebox}[Two-stage sample spaces]
\textbf{(a) Two coins.} Recording (first coin, second coin):
\[ S = \{HH,\ HT,\ TH,\ TT\}, \qquad n(S) = 2 \times 2 = 4. \]
Note that $HT$ and $TH$ are \emph{different} outcomes: the coins are distinct (one is first, the other second).

\textbf{(b) A coin and a die.} Each outcome is a pair (coin result, die score), giving $n(S) = 2 \times 6 = 12$ outcomes, displayed in the two-way table below.
\end{exemplebox}

\begin{popfigure}
\begin{center}
\begin{tikzpicture}[scale=0.62, every node/.style={font=\small}]
\draw[popline] (0,0) grid (6,2);
\draw[popline, thick] (0,0) rectangle (6,2);
\draw[popline, thick] (0,1) -- (6,1);
\node[font=\small\bfseries, text=popblue] at (0.5,1.5) {$H$};
\node[font=\small\bfseries, text=popblue] at (0.5,0.5) {$T$};
\foreach \x in {1,...,6} {
  \node[font=\small\bfseries, text=popblue] at (\x-0.5,2.35) {\x};
  \node at (\x-0.5,1.5) {$(H,\x)$};
  \node at (\x-0.5,0.5) {$(T,\x)$};
}
\node[font=\small, text=popmuted] at (3,-0.6) {die score};
\node[font=\small, text=popmuted, rotate=90] at (-0.7,1) {coin};
\end{tikzpicture}
\end{center}
\legende{Two-way table for tossing a coin and rolling a die: $n(S) = 12$ equally likely pairs.}
\end{popfigure}

\begin{exemplebox}[Two dice: the $6 \times 6$ grid]
When two dice (say a red one and a blue one) are rolled, an outcome is an ordered pair $(\text{red}, \text{blue})$. There are $6 \times 6 = 36$ outcomes, best displayed as a grid. The figure below shows the full sample space, with the outcomes giving a \emph{total of $8$} highlighted.
\end{exemplebox}

\begin{popfigure}
\begin{center}
\begin{tikzpicture}[scale=0.52, every node/.style={font=\scriptsize}]
% highlight sum = 8 : (2,6),(3,5),(4,4),(5,3),(6,2)
\foreach \p in {(1,5),(2,4),(3,3),(4,2),(5,1)} {
  \fill[poporange, opacity=0.45] \p rectangle ++(1,1);
}
\draw[popline] (0,0) grid (6,6);
\draw[popdark, thick] (0,0) rectangle (6,6);
\foreach \x in {1,...,6} {
  \node[font=\scriptsize\bfseries, text=popblue] at (\x-0.5,6.4) {\x};
  \node[font=\scriptsize\bfseries, text=popblue] at (-0.4,\x-0.5) {\x};
}
\foreach \x in {1,...,6} {
  \foreach \y in {1,...,6} {
    \node at (\x-0.5,\y-0.5) {(\x,\y)};
  }
}
\node[font=\small, text=popmuted] at (3,-0.9) {first die};
\node[font=\small, text=popmuted, rotate=90] at (-1.1,3) {second die};
\end{tikzpicture}
\end{center}
\legende{The $36$ outcomes of two dice. The shaded diagonal is the event ``sum $= 8$'': $(2,6), (3,5), (4,4), (5,3), (6,2)$ --- five outcomes, not one.}
\end{popfigure}

\courssub{Events and Types of Events}

\begin{definition}[Event]
An \cle{event} $E$ is any subset of the sample space $S$: a collection of outcomes sharing a common description. We say the event $E$ \emph{occurs} when the actual outcome of the trial belongs to $E$. The number of outcomes in $E$ is written $n(E)$.
\end{definition}

An event can be described in \emph{words} (``an even score'') or in \emph{set notation} ($E = \{2, 4, 6\}$). The GCE expects you to move fluently between the two.

\begin{exemplebox}[From words to sets]
One die is rolled, $S = \{1,2,3,4,5,6\}$.
\begin{itemize}
\item ``An even score'': $A = \{2, 4, 6\}$.
\item ``A score greater than $4$'': $B = \{5, 6\}$.
\item ``A score of at least $1$'': $C = \{1,2,3,4,5,6\} = S$.
\item ``A score of $7$'': $D = \emptyset$ (the empty set).
\end{itemize}
\end{exemplebox}

\begin{definition}[Simple, compound, sure and impossible events]
Let $S$ be a sample space and $E \subseteq S$ an event.
\begin{itemize}
\item $E$ is a \cle{simple} (or \emph{elementary}) event if it contains exactly one outcome ($n(E)=1$).
\item $E$ is a \cle{compound} event if it contains two or more outcomes ($n(E) \ge 2$).
\item $E$ is a \cle{sure} (or \emph{certain}) event if $E = S$: it must occur in every trial ($P(E)=1$).
\item $E$ is an \cle{impossible} event if $E = \emptyset$: it can never occur ($P(E)=0$).
\item The \cle{complementary event} of $E$, written $E'$ (or $\bar{E}$ or $E^c$), is the set of outcomes of $S$ that are \emph{not} in $E$: the event ``$E$ does not occur''. Clearly $E \cup E' = S$ and $E \cap E' = \emptyset$.
\end{itemize}
\end{definition}

\begin{exemplebox}[Classifying events]
Two coins are tossed: $S = \{HH, HT, TH, TT\}$.
\begin{itemize}
\item $E_1 = \{HH\}$ is simple.
\item $E_2 = $ ``at least one head'' $= \{HH, HT, TH\}$ is compound.
\item $E_3 = $ ``at most two heads'' $= S$ is certain.
\item $E_4 = $ ``three heads'' $= \emptyset$ is impossible.
\item The complement of $E_2$ is $E_2' = \{TT\}$: ``no head at all''.
\end{itemize}
\end{exemplebox}

\begin{attention}[An outcome is not the same thing as an event]
An \emph{outcome} is a single possible result; an \emph{event} is a \emph{set} of outcomes, which may contain one, several, all, or none of them. When a die shows $4$, the outcome is $4$, but several events have occurred simultaneously: ``an even score'' $\{2,4,6\}$, ``a score above $3$'' $\{4,5,6\}$, ``not a $1$'', and so on. In GCE answers, write events with set braces: $E = \{2,4,6\}$, not $E = 2,4,6$.
\end{attention}

\begin{attention}[Do not assume outcomes are equally likely]
Listing $S$ does \emph{not} mean all its members have the same chance. If you record only the \emph{sum} of two dice, the ``outcomes'' $2, 3, \dots, 12$ are \textbf{not} equally likely: a sum of $8$ arises from $5$ of the $36$ equally likely pairs, while a sum of $2$ arises from only $1$. The safe habit is to work with the \emph{equally likely} underlying outcomes (the $36$ ordered pairs) and build events from them.
\end{attention}

\courssub{The Probability Scale: a First Contact}

We shall define probability rigorously in the next section; for now, we need the scale on which it lives.

\begin{definition}[{Probability as a number in $[0,1]$}]
To every event $E$ we attach a number $P(E)$, called the \cle{probability} of $E$, which measures how likely $E$ is to occur. It always satisfies
\[ 0 \le P(E) \le 1. \]
An impossible event has probability $0$; a certain event has probability $1$; the closer $P(E)$ is to $1$, the more likely the event.
\end{definition}

\begin{popfigure}
\begin{center}
\begin{tikzpicture}[scale=1]
\draw[thick, popdark, ->] (0,0) -- (10.6,0);
\foreach \x/\lab in {0/0, 2.5/0.25, 5/0.5, 7.5/0.75, 10/1} {
  \draw[popdark] (\x,0.12) -- (\x,-0.12);
  \node[below, font=\small] at (\x,-0.15) {\lab};
}
\node[above, font=\small, text=popblue] at (0,0.2) {impossible};
\node[above, font=\small, text=poporange] at (2.5,0.2) {unlikely};
\node[above, font=\small, text=poppurple, align=center] at (5,0.2) {even chance};
\node[above, font=\small, text=popgreen] at (7.5,0.2) {likely};
\node[above, font=\small, text=popblue] at (10,0.2) {certain};
\fill[popblue] (0,0) circle (2.2pt);
\fill[poporange] (2.5,0) circle (2.2pt);
\fill[poppurple] (5,0) circle (2.2pt);
\fill[popgreen] (7.5,0) circle (2.2pt);
\fill[popblue] (10,0) circle (2.2pt);
\node[below, font=\small\itshape, text=popmuted] at (5,-0.85) {The probability scale};
\end{tikzpicture}
\end{center}
\legende{The probability scale: every probability lies between $0$ (impossible) and $1$ (certain).}
\end{popfigure}

\begin{aretenir}[Key points of this section]
\begin{itemize}
\item A random experiment has unpredictable single results but a known set of possible results; a deterministic experiment has a predictable result.
\item The sample space $S$ is the set of \emph{all} outcomes; list it systematically (list, table, grid) and check $n(S)$ with the product rule for multi-stage experiments.
\item An event is a \emph{subset} of $S$; events are simple, compound, certain ($=S$), impossible ($=\emptyset$), and every event $E$ has a complement $E'$.
\item Probability is a number in $[0,1]$: $0$ means impossible, $1$ means certain.
\end{itemize}
\end{aretenir}

\begin{exoresolu}[Translating everyday statements into events]
A fair die is rolled once. Write each of the following events in set notation and state whether it is simple, compound, certain or impossible: (a) an even score; (b) a prime score; (c) a score less than $7$; (d) a score of $0$; (e) at least $5$.
\tcblower
The sample space is $S = \{1,2,3,4,5,6\}$.
\begin{itemize}
\item[(a)] $A = \{2,4,6\}$: three outcomes, so $A$ is \textbf{compound}.
\item[(b)] The primes in $S$ are $2, 3, 5$, so $B = \{2,3,5\}$: \textbf{compound}. (Remember: $1$ is \emph{not} prime.)
\item[(c)] $C = \{1,2,3,4,5,6\} = S$: \textbf{certain}.
\item[(d)] $D = \emptyset$: \textbf{impossible}.
\item[(e)] ``At least $5$'' means $5$ or more: $E = \{5,6\}$: \textbf{compound}.
\end{itemize}
\end{exoresolu}

\begin{exoresolu}[Building a sample space and reading off events]
A coin is tossed and a die is rolled. (a) Write down the sample space and give $n(S)$. (b) List the outcomes of the event $A$: ``a head and an even number''. (c) List the event $B$: ``a tail or a number greater than $4$''. (d) Describe $A'$ in words.
\tcblower
(a) $S = \{(H,1),(H,2),(H,3),(H,4),(H,5),(H,6),(T,1),(T,2),(T,3),(T,4),(T,5),(T,6)\}$, so $n(S) = 2 \times 6 = 12$.

(b) A head together with an even die score: $A = \{(H,2),(H,4),(H,6)\}$.

(c) A tail (any score) or a score above $4$ (any coin): $B = \{(T,1),(T,2),(T,3),(T,4),(T,5),(T,6),(H,5),(H,6)\}$.

(d) $A'$ is ``not (a head and an even number)'', i.e.\ ``a tail, or a head with an odd number''.
\end{exoresolu}

\begin{exercice}[Vocabulary and sample spaces]
\begin{enumerate}
\item A student is chosen at random from a class of $25$ girls and $15$ boys. What is the sample space if we record only the student's name? What is $n(S)$?
\item Two fair dice, one red and one blue, are thrown. List the outcomes of the events: (a) ``both dice show the same number''; (b) ``the sum is $10$ or more''; (c) ``the red die shows $3$''. State which of these are simple and which are compound.
\item A card is drawn from a standard pack of $52$ cards. Express in words and in set notation (using suits and values) the complement of the event ``drawing a heart''.
\item Place on the probability scale a reasonable position for each statement: (a) ``the sun will rise tomorrow''; (b) ``it will rain in Douala at least once next August''; (c) ``a tossed fair coin lands on its edge''.
\end{enumerate}
\end{exercice}

\begin{corrige}[Exercise 1]
\begin{enumerate}
\item $S$ is the set of the $40$ students' names; $n(S) = 25 + 15 = 40$.
\item (a) $\{(1,1),(2,2),(3,3),(4,4),(5,5),(6,6)\}$ --- compound. (b) $\{(4,6),(5,5),(5,6),(6,4),(6,5),(6,6)\}$ --- compound. (c) $\{(3,1),(3,2),(3,3),(3,4),(3,5),(3,6)\}$ --- compound. None is simple, since each contains more than one outcome.
\item The complement is ``drawing a card that is not a heart'', i.e.\ the set of the $39$ spades, diamonds and clubs.
\item (a) essentially at $1$ (certain); (b) very close to $1$ (very likely); (c) essentially at $0$ (practically impossible).
\end{enumerate}
\end{corrige}

\begin{savaistu}[Why ``at least'' and ``at most'' matter at the GCE]
Examiners report that a large share of lost marks in probability comes from misreading small words: ``at least one head'' \emph{includes} the case of two heads, while ``exactly one head'' excludes it. A reliable reflex: translate the phrase into an inequality first (``at least one head'' means ``number of heads $\ge 1$''), then list the outcomes satisfying it. This habit, learned now, will pay off in every section of this chapter.
\end{savaistu}

We now have the language of probability: experiments, outcomes, sample spaces and events, together with the scale on which chance is measured. In the next section we attach numbers to this language, defining classical (theoretical) probability for equally likely outcomes and empirical probability from observed frequencies.

\coursec{Classical and Empirical Probability}

In the previous section we built the language of probability: we defined an \cle{experiment}, its \cle{outcomes}, the \cle{sample space} $S$, and \cle{events} as subsets of $S$. We now answer the natural next question: \emph{how do we assign a number to an event to measure how likely it is?} There are two great roads to this number. The first is purely logical: when all outcomes are equally likely, we simply \emph{count}. This is the \cle{classical} (or theoretical) probability. The second is experimental: we repeat the experiment many times and observe how often the event occurs. This is the \cle{empirical} (or experimental) probability. The deep and beautiful fact, which we shall meet informally as the \emph{law of large numbers}, is that the two roads lead to the same destination.

\courssub{Classical (Theoretical) Probability}

\textbf{Intuition.} A fair die is thrown once. There is no reason for any face to be favoured over another, so each of the six faces is \emph{equally likely}. If we ask for the probability of an even number, three of the six equally likely outcomes ($2, 4, 6$) suit us. It is natural to say the probability is $\frac{3}{6} = \frac{1}{2}$. This counting idea is the whole of classical probability.

\begin{definition}[Classical probability]
Let $S$ be a finite sample space in which all outcomes are \cle{equally likely}. The \cle{classical probability} of an event $E$ is
\[
P(E) = \frac{n(E)}{n(S)} = \frac{\text{number of outcomes favourable to } E}{\text{total number of possible outcomes}}.
\]
\end{definition}

The hypothesis ``equally likely'' is essential: the formula is a \emph{definition} built on symmetry, not a theorem that applies to every experiment. A fair coin, a fair die, a well-shuffled pack of cards, identical balls thoroughly mixed in a bag — these are the standard situations where the formula applies. In examination questions, the words \emph{fair}, \emph{unbiased}, \emph{at random} and \emph{identical except for colour} are the signals that equiprobability may be assumed.

\begin{exemplebox}[Balls in a bag]
A bag contains $5$ red balls and $7$ blue balls, identical except for colour. One ball is drawn at random. Then $n(S) = 12$. If $R$ is the event ``the ball is red'', $P(R) = \frac{5}{12}$; if $B$ is the event ``the ball is blue'', $P(B) = \frac{7}{12}$. Note that $P(R) + P(B) = 1$, as it must, since every ball is either red or blue.
\end{exemplebox}

\begin{exoresolu}[Two dice]
Two fair dice are thrown. Find the probability that the sum of the scores is $9$.
\tcblower
The sample space consists of $6 \times 6 = 36$ equally likely ordered pairs. The outcomes giving a sum of $9$ are
\[
(3,6),\ (4,5),\ (5,4),\ (6,3),
\]
so $n(E) = 4$. Hence
\[
P(\text{sum is } 9) = \frac{4}{36} = \frac{1}{9}.
\]
\end{exoresolu}

\begin{attention}[Equally likely means equally likely]
When two dice are thrown, the \emph{sums} $2, 3, \dots, 12$ are \textbf{not} equally likely: there is only one way to make $2$ but six ways to make $7$. Never apply $\frac{n(E)}{n(S)}$ to outcomes that are not equally likely. Always reduce the situation to the underlying equally likely outcomes (here, the $36$ ordered pairs). A two-way table or a list of the $36$ pairs is the safest way to count.
\end{attention}

\textbf{Properties of classical probability.} Since $E \subseteq S$, we have $0 \le n(E) \le n(S)$, and dividing by $n(S) > 0$ gives the fundamental bounds. The empty event contains no outcome, and the whole sample space contains all of them. Finally, $E$ and its complement $E'$ partition $S$, so $n(E) + n(E') = n(S)$.

\begin{propriete}[Basic properties of probability]
For any event $E$ in a finite equiprobable sample space $S$:
\[
0 \le P(E) \le 1, \qquad P(S) = 1, \qquad P(\varnothing) = 0, \qquad P(E') = 1 - P(E).
\]
These properties (in particular $P(E') = 1 - P(E)$) are examinable and follow directly from counting, since $n(E) + n(E') = n(S)$.
\end{propriete}

An event with probability $0$ is called \cle{impossible} and an event with probability $1$ is called \cle{certain}. The last formula is a powerful working tool: when the event $E$ is complicated but its complement is simple, compute $P(E')$ and subtract from $1$.

\begin{exemplebox}[Using the complement]
A fair coin is tossed three times. What is the probability of obtaining \emph{at least one} head? Here $n(S) = 2^3 = 8$. The complement of ``at least one head'' is ``no head at all'', i.e.\ the single outcome $TTT$. Hence
\[
P(\text{at least one head}) = 1 - \frac{1}{8} = \frac{7}{8}.
\]
Listing all favourable outcomes would have been longer and more error-prone.
\end{exemplebox}

\begin{exoresolu}[Probability with a pack of cards]
A card is drawn at random from a well-shuffled pack of $52$ cards. Find the probability that it is (a) a king; (b) a heart; (c) not a spade.
\tcblower
Here $n(S) = 52$, all cards equally likely.
\begin{enumerate}
\item There are $4$ kings, so $P(\text{king}) = \frac{4}{52} = \frac{1}{13}$.
\item There are $13$ hearts, so $P(\text{heart}) = \frac{13}{52} = \frac{1}{4}$.
\item Using the complement: $P(\text{not a spade}) = 1 - \frac{13}{52} = \frac{3}{4}$.
\end{enumerate}
\end{exoresolu}

\begin{savaistu}[Odds — an extension you may meet]
In games and betting, likelihood is sometimes expressed as \cle{odds} rather than probability. The \emph{odds in favour} of $E$ are the ratio $n(E) : n(E')$, and the \emph{odds against} $E$ are $n(E') : n(E)$. For example, drawing an ace from a pack of $52$ cards has odds in favour of $4 : 48$, i.e.\ $1 : 12$, corresponding to $P = \frac{1}{13}$. Conversion: odds $a : b$ in favour correspond to $P(E) = \frac{a}{a+b}$. This is an extension beyond the strict syllabus, but it appears occasionally in GCE-style questions.
\end{savaistu}

\courssub{Empirical (Experimental) Probability}

\textbf{Intuition.} A drawing pin is thrown in the air. It can land \emph{point up} or \emph{point down}, but there is no symmetry to tell us the two outcomes are equally likely — and indeed they are not. Classical probability is powerless here. The only honest way to estimate the probability of ``point up'' is to throw the pin many times and observe.

\begin{definition}[Empirical probability]
If an experiment is repeated $n$ times under identical conditions and the event $E$ occurs $f$ times, the \cle{relative frequency} of $E$ is
\[
\text{relative frequency of } E = \frac{f}{n} = \frac{\text{number of times } E \text{ occurs}}{\text{number of trials}}.
\]
The \cle{empirical probability} of $E$ is the relative frequency obtained from a large number of trials; it serves as an estimate of the true probability.
\end{definition}

\begin{methode}[Designing and recording a probability experiment]
\begin{enumerate}
\item \textbf{Define} the experiment, the trial, and the event $E$ precisely (e.g.\ toss a coin; $E$ = ``head'').
\item \textbf{Prepare a table} with three columns: outcome, tally, frequency; add a fourth column for relative frequency $\frac{f}{n}$.
\item \textbf{Perform} a fixed number $n$ of trials under identical conditions, recording each result immediately.
\item \textbf{Compute} the relative frequencies and check that they sum to $1$.
\item \textbf{Compare} with the theoretical probability (when it exists) and comment, bearing in mind the number of trials.
\end{enumerate}
\end{methode}

\begin{experience}[Tossing a coin 50 times]
Work in pairs. Toss a coin $50$ times, recording the result of each toss. After every $10$ tosses, compute the \emph{cumulative} relative frequency of heads: $\frac{\text{heads so far}}{\text{tosses so far}}$. A typical set of results might be:

\begin{center}
\begin{tabularx}{\linewidth}{|l|X|X|X|X|X|}
\hline
\rowcolor{popblueL} Tosses completed & $10$ & $20$ & $30$ & $40$ & $50$ \\
\hline Heads so far & $7$ & $11$ & $14$ & $19$ & $24$ \\
\hline Relative frequency & $0.70$ & $0.55$ & $0.47$ & $0.48$ & $0.48$ \\
\hline
\end{tabularx}
\end{center}

Observe how the relative frequency swings wildly at first ($0.70$ after $10$ tosses) and then settles near $0.5$ as the number of tosses grows. Repeat the experiment with a drawing pin, where the theoretical value is unknown: the stabilised relative frequency is then your best estimate of the probability.
\end{experience}

\begin{popfigure}
\begin{center}
\begin{tikzpicture}
\begin{axis}[
  width=0.85\linewidth, height=6cm,
  xlabel={Number of tosses}, ylabel={Cumulative relative frequency of heads},
  xmin=0, xmax=200, ymin=0.2, ymax=0.8,
  xtick={0,50,100,150,200}, ytick={0.2,0.3,0.4,0.5,0.6,0.7,0.8},
  grid=both, grid style={popline, very thin},
]
\addplot[popblue, thick, mark=*, mark size=1pt] coordinates {
  (1,1) (2,0.5) (5,0.8) (10,0.7) (15,0.6) (20,0.55) (25,0.6)
  (30,0.533) (40,0.525) (50,0.48) (60,0.517) (75,0.507)
  (100,0.49) (125,0.504) (150,0.493) (175,0.503) (200,0.5)
};
\addplot[poporange, dashed, thick] coordinates {(0,0.5) (200,0.5)};
\node[poporange, anchor=south west, font=\small] at (axis cs:105,0.51) {theoretical $P = 0.5$};
\end{axis}
\end{tikzpicture}
\end{center}
\legende{Cumulative relative frequency of ``heads'' against the number of tosses: the fluctuations die out and the curve converges towards $0.5$.}
\end{popfigure}

This stabilisation is not a coincidence; it is one of the most important facts in all of probability.

\begin{propriete}[Law of large numbers (informal statement)]
As the number of trials of an experiment increases indefinitely, the relative frequency of an event $E$ stabilises towards a fixed value, which is the true probability $P(E)$. The statement is informal here; its precise proof belongs to university mathematics and is not examinable.
\end{propriete}

The law of large numbers is the bridge between the two halves of this section: it guarantees that empirical probability, computed from \emph{many} trials, is a reliable estimate of the theoretical probability.

\begin{attention}[Few trials prove nothing]
If you toss a fair coin $4$ times and get $3$ heads, the relative frequency of heads is $0.75$ — far from $0.5$ — yet the coin may be perfectly fair. Short runs fluctuate enormously. Never conclude that an object is biased from a handful of trials; only a \emph{large} number of trials, with a relative frequency stabilising clearly away from the theoretical value, gives serious evidence of bias.
\end{attention}

\courssub{Comparing Theory and Experiment: Is This Die Fair?}

A practical use of empirical probability is to test whether a real object behaves like its theoretical model. Suppose we suspect a die of being loaded. We throw it $600$ times and record:

\begin{center}
\begin{tabularx}{\linewidth}{|l|X|X|X|X|X|X|}
\hline
\rowcolor{popblueL} Face & $1$ & $2$ & $3$ & $4$ & $5$ & $6$ \\
\hline Frequency & $95$ & $102$ & $98$ & $90$ & $101$ & $114$ \\
\hline Relative frequency & $0.158$ & $0.170$ & $0.163$ & $0.150$ & $0.168$ & $0.190$ \\
\hline
\end{tabularx}
\end{center}

For a fair die, the theoretical probability of each face is $\frac{1}{6} \approx 0.167$. The observed relative frequencies all lie close to this value; the largest deviation, for face $6$, is about $0.023$, which is entirely plausible for $600$ throws. We conclude that the data give \emph{no convincing evidence} that the die is biased. Had face $6$ appeared, say, $180$ times (relative frequency $0.30$), suspicion would be justified.

\begin{popfigure}
\begin{center}
\begin{tikzpicture}
\begin{axis}[
  width=0.85\linewidth, height=6cm,
  ybar, bar width=10pt,
  xlabel={Face of the die}, ylabel={Probability},
  symbolic x coords={1,2,3,4,5,6},
  xtick=data, ymin=0, ymax=0.22,
  legend style={at={(0.5,1.02)}, anchor=south, legend columns=2, font=\small},
  grid=both, grid style={popline, very thin},
]
\addplot[fill=popblue, draw=popdark] coordinates {(1,0.167) (2,0.167) (3,0.167) (4,0.167) (5,0.167) (6,0.167)};
\addplot[fill=poporange, draw=popdark] coordinates {(1,0.158) (2,0.170) (3,0.163) (4,0.150) (5,0.168) (6,0.190)};
\legend{Theoretical $\frac{1}{6}$, Experimental ($600$ throws)}
\end{axis}
\end{tikzpicture}
\end{center}
\legende{Theoretical probability versus experimental relative frequency for each face of a die thrown $600$ times: the bars are close, so the die appears fair.}
\end{popfigure}

\begin{exoresolu}[Estimating a probability from data]
A trader at Mokolo market tosses a bottle top $200$ times; it lands ``open side up'' $86$ times. (a) Estimate the probability that the top lands open side up. (b) Why would the classical formula $\frac{1}{2}$ be wrong here?
\tcblower
(a) The empirical probability is the relative frequency:
\[
P(\text{open side up}) \approx \frac{86}{200} = 0.43.
\]
(b) The two outcomes ``open side up'' and ``open side down'' are \emph{not equally likely}: the shape of the bottle top breaks the symmetry. Classical probability requires equally likely outcomes, so $\frac{n(E)}{n(S)} = \frac{1}{2}$ has no justification. Only experiment can estimate this probability, and $200$ trials give a reasonably stable estimate.
\end{exoresolu}

\begin{exoresolu}[Expected frequency]
A fair die is thrown $240$ times. How many times would you \emph{expect} the face $5$ to appear?
\tcblower
The theoretical probability of a $5$ is $\frac{1}{6}$. The \cle{expected frequency} is
\[
\text{expected frequency} = n \times P(E) = 240 \times \frac{1}{6} = 40.
\]
This is an average over many repetitions of the experiment: in any single run of $240$ throws, the actual count will usually differ a little from $40$.
\end{exoresolu}

\begin{exercice}[]
\begin{enumerate}
\item A bag contains $4$ white, $3$ green and $5$ black balls. One ball is drawn at random. Find the probability that it is (a) green; (b) not white.
\item Two fair coins are tossed. Using the complement, find the probability of at least one tail.
\item A card is drawn at random from a pack of $52$ cards. Find the probability that it is (a) a red card; (b) an ace; (c) not a face card (king, queen or jack).
\item A coin is tossed $1000$ times and lands heads $512$ times. (a) Compute the relative frequency of heads. (b) Is there strong evidence that the coin is biased? Justify your answer.
\item A die is thrown $300$ times and the face $6$ appears $95$ times. Comment on the fairness of the die.
\item The probability that a seed germinates is estimated at $0.8$. If a farmer plants $500$ seeds, how many does he expect to germinate?
\end{enumerate}
\end{exercice}

\begin{corrige}[Exercise 1]
\begin{enumerate}
\item $n(S) = 12$. (a) $P(\text{green}) = \frac{3}{12} = \frac{1}{4}$. (b) $P(\text{not white}) = 1 - \frac{4}{12} = \frac{2}{3}$.
\item $n(S) = 4$; complement ``no tail'' $= \{HH\}$, so $P = 1 - \frac{1}{4} = \frac{3}{4}$.
\item (a) $\frac{26}{52} = \frac{1}{2}$. (b) $\frac{4}{52} = \frac{1}{13}$. (c) There are $12$ face cards, so $P = 1 - \frac{12}{52} = \frac{40}{52} = \frac{10}{13}$.
\item (a) $\frac{512}{1000} = 0.512$. (b) No: $0.512$ is close to $0.5$ and such a fluctuation is plausible in $1000$ tosses of a fair coin.
\item Relative frequency $\frac{95}{300} \approx 0.317$, nearly double the theoretical $\frac{1}{6} \approx 0.167$. With $300$ trials this is a large deviation: there is serious evidence that the die is biased towards $6$.
\item Expected number $= 500 \times 0.8 = 400$ seeds.
\end{enumerate}
\end{corrige}

\begin{aretenir}[Key points of the section]
\begin{itemize}
\item Classical probability: $P(E) = \dfrac{n(E)}{n(S)}$, valid only when outcomes are \cle{equally likely}.
\item Always: $0 \le P(E) \le 1$, $P(S) = 1$, $P(\varnothing) = 0$, and $P(E') = 1 - P(E)$ — use the complement for ``at least one'' questions.
\item Empirical probability is the relative frequency $\frac{f}{n}$ from repeated trials; it is the only option when outcomes are not equally likely.
\item Law of large numbers: relative frequency stabilises towards the true probability as $n$ grows — so judge bias only from \emph{many} trials.
\item Expected frequency in $n$ trials: $n \times P(E)$.
\end{itemize}
\end{aretenir}

\coursec{The Addition Law and Venn Diagrams}

In the previous section we learnt how to measure the probability of a \emph{single} event, either by counting equally likely outcomes (classical probability) or by repeating an experiment many times (empirical probability). Real questions, however, rarely concern one event alone. A student in Yaound\'e may ask: \emph{``What is the probability that a randomly chosen classmate plays football \textbf{or} handball?''} A GCE candidate may ask: \emph{``What is the probability of scoring grade A in at least one of two papers?''} To answer such questions we need a rule for combining probabilities: the \cle{addition law}. And to apply it correctly, we need a picture that shows clearly how events overlap: the \cle{Venn diagram}. That is the work of this section.

\courssub{From ``or'' to addition: the intuition}

Suppose a fair die is thrown once. Let $A$ = ``the score is even'' $=\{2,4,6\}$ and $B$ = ``the score is a multiple of $3$'' $=\{3,6\}$. The event ``$A$ \textbf{or} $B$'' is the \cle{union}
\[ A\cup B=\{2,3,4,6\},\qquad P(A\cup B)=\frac{4}{6}=\frac{2}{3}. \]
Notice that $P(A)+P(B)=\dfrac{3}{6}+\dfrac{2}{6}=\dfrac{5}{6}$, which is \emph{too big}. Why? Because the outcome $6$ belongs to \emph{both} events and has been counted twice. Subtracting the overlap once repairs the count:
\[ P(A\cup B)=P(A)+P(B)-P(A\cap B)=\frac{3}{6}+\frac{2}{6}-\frac{1}{6}=\frac{4}{6}. \]
This simple observation is the whole secret of the addition law.

\courssub{Mutually exclusive events}

\begin{definition}[Mutually exclusive events]
Two events $A$ and $B$ are \cle{mutually exclusive} (or \emph{disjoint}) if they cannot occur at the same time, i.e.
\[ A\cap B=\varnothing,\qquad\text{equivalently}\qquad P(A\cap B)=0. \]
\end{definition}

\begin{exemplebox}[Testing for exclusivity]
A card is drawn from an ordinary pack of $52$ cards.
\begin{itemize}
\item $A$ = ``the card is a King'', $B$ = ``the card is a Queen'': a card cannot be both, so $A\cap B=\varnothing$; $A$ and $B$ are mutually exclusive.
\item $A$ = ``the card is a King'', $C$ = ``the card is a heart'': the King of hearts belongs to both, so $A$ and $C$ are \textbf{not} mutually exclusive.
\end{itemize}
\end{exemplebox}

\begin{popfigure}
\begin{center}
\begin{tikzpicture}[scale=0.9]
\draw[thick, popink] (-3.4,-2.1) rectangle (3.4,2.1);
\node[popink] at (3.0,1.75) {$\mathcal{E}$};
\draw[thick, popblue, fill=popblueL] (-1.8,0) circle (1.0);
\draw[thick, poporange, fill=popgold!40] (1.8,0) circle (1.0);
\node at (-1.8,0) {$A$};
\node at (1.8,0) {$B$};
\node at (0,0) {$\varnothing$};
\end{tikzpicture}
\end{center}
\legende{Mutually exclusive events: the circles for $A$ and $B$ do not overlap, because $A\cap B=\varnothing$.}
\end{popfigure}

\courssub{The addition law}

\begin{propriete}[Addition law of probability]
Let $A$ and $B$ be two events of the same experiment.
\begin{itemize}
\item \textbf{Exclusive case.} If $A$ and $B$ are mutually exclusive,
\[ P(A\cup B)=P(A)+P(B). \]
\item \textbf{General case.} For \emph{any} two events,
\[ P(A\cup B)=P(A)+P(B)-P(A\cap B). \]
\end{itemize}
The general formula is examinable and its justification (below) is instructive.
\end{propriete}

\textbf{Justification by counting.} Suppose the sample space $\mathcal{E}$ has $n$ equally likely outcomes. Write $n(A)$, $n(B)$, $n(A\cap B)$ for the numbers of outcomes in each set. When we form $n(A)+n(B)$, every outcome lying in \emph{both} $A$ and $B$ is counted twice, so it must be removed once:
\[ n(A\cup B)=n(A)+n(B)-n(A\cap B). \]
Dividing each term by $n$ and using the classical definition $P(E)=\dfrac{n(E)}{n}$ gives
\[ P(A\cup B)=P(A)+P(B)-P(A\cap B). \]
If $A\cap B=\varnothing$ then $P(A\cap B)=0$ and the general law reduces to the exclusive law: the exclusive case is just a special case of the general one.

\begin{attention}[Do not add blindly]
Adding $P(A)+P(B)$ \textbf{without checking whether $A$ and $B$ overlap} double-counts the intersection $A\cap B$ and gives an answer that is too large — sometimes even greater than $1$, which is impossible for a probability. Always ask first: \emph{``Can the two events happen together?''} If yes, subtract $P(A\cap B)$.
\end{attention}

\begin{exoresolu}[Addition law with a pack of cards]
A card is drawn at random from a pack of $52$ cards. Find the probability that it is (a) a King or a Queen; (b) a King or a heart.
\tcblower
(a) Let $K$ = ``King'', $Q$ = ``Queen''. These are mutually exclusive, so
\[ P(K\cup Q)=P(K)+P(Q)=\frac{4}{52}+\frac{4}{52}=\frac{8}{52}=\frac{2}{13}. \]
(b) Let $H$ = ``heart''. $K$ and $H$ are \emph{not} exclusive since $K\cap H=\{\text{King of hearts}\}$, so $P(K\cap H)=\dfrac{1}{52}$. By the general addition law,
\[ P(K\cup H)=P(K)+P(H)-P(K\cap H)=\frac{4}{52}+\frac{13}{52}-\frac{1}{52}=\frac{16}{52}=\frac{4}{13}. \]
\end{exoresolu}

\courssub{Complementary and exhaustive events}

\begin{definition}[Complement; exhaustive events]
\begin{itemize}
\item The \cle{complement} of an event $E$, written $E'$ (or $\bar{E}$), is the event ``$E$ does not occur''. It contains every outcome of $\mathcal{E}$ that is not in $E$.
\item Events $E_1,E_2,\dots,E_k$ are \cle{exhaustive} if at least one of them must occur, i.e. $E_1\cup E_2\cup\cdots\cup E_k=\mathcal{E}$.
\end{itemize}
In particular $E$ and $E'$ are always both mutually exclusive \emph{and} exhaustive.
\end{definition}

\begin{propriete}[Complement rule]
For every event $E$,
\[ P(E)+P(E')=1,\qquad\text{so}\qquad P(E')=1-P(E). \]
\end{propriete}

\textbf{Proof.} $E$ and $E'$ are mutually exclusive, so by the exclusive addition law $P(E)+P(E')=P(E\cup E')=P(\mathcal{E})=1$. \hfill$\blacksquare$

\begin{methode}[``At least one'' problems: use the complement]
\begin{enumerate}
\item Identify the event $E$ = ``at least one success''.
\item Recognise that its complement $E'$ = ``\textbf{no} success at all'' is usually much easier to compute.
\item Compute $P(E')$, then conclude $P(E)=1-P(E')$.
\end{enumerate}
\end{methode}

\begin{exemplebox}[At least one six]
A fair die is thrown twice. What is the probability of obtaining \emph{at least one} six?
Listing all favourable pairs among the $36$ outcomes is tedious. Instead, $P(\text{no six})=\dfrac{5}{6}\times\dfrac{5}{6}=\dfrac{25}{36}$, hence
\[ P(\text{at least one six})=1-\frac{25}{36}=\frac{11}{36}. \]
\end{exemplebox}

\courssub{Representing probabilities with two-set Venn diagrams}

A Venn diagram displays the sample space $\mathcal{E}$ as a rectangle and the events $A$, $B$ as overlapping circles inside it. The plane is cut into \textbf{four regions}, each with a precise meaning:

\begin{center}
\begin{tabularx}{\linewidth}{|l|X|}
\hline
\rowcolor{popblueL} \textbf{Region} & \textbf{Meaning} \\
\hline
$A\cap B$ & both $A$ and $B$ occur \\
\hline
$A\cap B'$ ($A$ only) & $A$ occurs but $B$ does not \\
\hline
$A'\cap B$ ($B$ only) & $B$ occurs but $A$ does not \\
\hline
$A'\cap B'=(A\cup B)'$ & neither $A$ nor $B$ occurs \\
\hline
\end{tabularx}
\end{center}

\begin{popfigure}
\begin{center}
\begin{tikzpicture}[scale=1.0]
\draw[thick, popink] (-3.6,-2.3) rectangle (3.6,2.3);
\node[popink] at (3.2,1.95) {$\mathcal{E}$};
% Draw circles
\draw[thick, popblue] (-0.7,0) circle (1.35);
\draw[thick, poporange] (0.7,0) circle (1.35);
% Fill regions with distinct colors
% A only (blue)
\begin{scope}
\clip (-0.7,0) circle (1.35);
\clip (0.7,0) circle (1.35);
\fill[popblueL] (-3.6,-2.3) rectangle (3.6,2.3);
\end{scope}
% B only (orange)
\begin{scope}
\clip (0.7,0) circle (1.35);
\clip (-0.7,0) circle (1.35) (-3.6,-2.3) rectangle (3.6,2.3);
\fill[popgold!40] (-3.6,-2.3) rectangle (3.6,2.3);
\end{scope}
% Intersection (green)
\begin{scope}
\clip (-0.7,0) circle (1.35);
\clip (0.7,0) circle (1.35);
\fill[popgreenL] (-3.6,-2.3) rectangle (3.6,2.3);
\end{scope}
% Redraw circles on top
\draw[thick, popblue] (-0.7,0) circle (1.35);
\draw[thick, poporange] (0.7,0) circle (1.35);
% Labels
\node at (0,0) {\small $A\cap B$};
\node at (-1.45,0) {\small $A$ only};
\node at (1.45,0) {\small $B$ only};
\node at (2.7,-1.7) {\small neither};
\node[popblue] at (-1.5,1.15) {$A$};
\node[poporange] at (1.5,1.15) {$B$};
\end{tikzpicture}
\end{center}
\legende{The four regions of a two-set Venn diagram: $A$ only (blue), $A\cap B$ (green), $B$ only (orange), and neither (white outside both circles).}
\end{popfigure}

\begin{methode}[Completing a Venn diagram from the inside outwards]
\begin{enumerate}
\item Start with the \textbf{intersection}: write the number (or probability) of $A\cap B$ in the overlap.
\item Fill ``$A$ only'' $=n(A)-n(A\cap B)$ and ``$B$ only'' $=n(B)-n(A\cap B)$.
\item Fill ``neither'' $=n(\mathcal{E})-n(A\cup B)$, where $n(A\cup B)$ is the sum of the three inner regions.
\item Check that the four regions add up to $n(\mathcal{E})$.
\end{enumerate}
Never write $n(A)$ itself inside the diagram: $n(A)$ is the \emph{sum} of two regions.
\end{methode}

\courssub{A Cameroon-context case study}

\begin{experience}[Football and handball in a Form Five class]
In a Form Five class of $40$ pupils in Bamenda, $25$ play football ($F$), $18$ play handball ($H$), and $10$ play both sports. A pupil is chosen at random. We investigate the probabilities of the various combinations.
\end{experience}

Working from the inside outwards: $n(F\cap H)=10$; football only $=25-10=15$; handball only $=18-10=8$; neither $=40-(15+10+8)=7$.

\begin{popfigure}
\begin{center}
\begin{tikzpicture}[scale=1.0]
\draw[thick, popink] (-3.6,-2.3) rectangle (3.6,2.3);
\node[popink] at (3.2,1.95) {$\mathcal{E}$};
% Circles
\draw[thick, popblue] (-0.7,0) circle (1.35);
\draw[thick, poporange] (0.7,0) circle (1.35);
% Fill regions
% F only
\begin{scope}
\clip (-0.7,0) circle (1.35);
\clip (0.7,0) circle (1.35);
\fill[popblueL] (-3.6,-2.3) rectangle (3.6,2.3);
\end{scope}
% H only
\begin{scope}
\clip (0.7,0) circle (1.35);
\clip (-0.7,0) circle (1.35) (-3.6,-2.3) rectangle (3.6,2.3);
\fill[popgold!40] (-3.6,-2.3) rectangle (3.6,2.3);
\end{scope}
% Intersection
\begin{scope}
\clip (-0.7,0) circle (1.35);
\clip (0.7,0) circle (1.35);
\fill[popgreenL] (-3.6,-2.3) rectangle (3.6,2.3);
\end{scope}
% Redraw circles
\draw[thick, popblue] (-0.7,0) circle (1.35);
\draw[thick, poporange] (0.7,0) circle (1.35);
% Numbers
\node at (0,0) {$10$};
\node at (-1.4,0) {$15$};
\node at (1.4,0) {$8$};
\node at (2.8,-1.7) {$7$};
\node[popblue] at (-1.6,1.2) {$F$};
\node[poporange] at (1.6,1.2) {$H$};
\end{tikzpicture}
\end{center}
\legende{Venn diagram for the class: $15$ football only, $10$ both, $8$ handball only, $7$ neither.}
\end{popfigure}

Reading directly from the diagram (each pupil is equally likely, so divide by $40$):
\[ P(F\cap H)=\frac{10}{40}=\frac14,\quad P(F\cup H)=\frac{33}{40},\quad P(F\text{ only})=\frac{15}{40}=\frac38,\quad P(\text{neither})=\frac{7}{40}. \]
Check with the addition law: $P(F\cup H)=\dfrac{25}{40}+\dfrac{18}{40}-\dfrac{10}{40}=\dfrac{33}{40}$. \checkmark

\begin{exoresolu}[GCE-style: probabilities given algebraically]
In a certain town, the probability that a household owns a television is $0.6$, the probability that it owns a radio is $0.7$, and the probability that it owns both is $0.4$. A household is chosen at random. Find the probability that it owns (a) a television or a radio; (b) a television only; (c) neither.
\tcblower
Let $T$ = ``owns a television'', $R$ = ``owns a radio''. Given: $P(T)=0.6$, $P(R)=0.7$, $P(T\cap R)=0.4$.

(a) By the general addition law,
\[ P(T\cup R)=0.6+0.7-0.4=0.9. \]
(b) ``Television only'' is $T\cap R'$:
\[ P(T\cap R')=P(T)-P(T\cap R)=0.6-0.4=0.2. \]
(c) ``Neither'' is the complement of the union:
\[ P\bigl((T\cup R)'\bigr)=1-P(T\cup R)=1-0.9=0.1. \]
(Placed on a Venn diagram, the four regions are $0.2$, $0.4$, $0.3$ and $0.1$, which sum to $1$.)
\end{exoresolu}

\begin{exoresolu}[GCE-style: finding an unknown from the addition law]
Events $A$ and $B$ are such that $P(A)=0.45$, $P(B)=0.5$ and $P(A\cup B)=0.75$. (a) Find $P(A\cap B)$. (b) Find $P(A'\cap B)$. (c) State, with a reason, whether $A$ and $B$ are mutually exclusive.
\tcblower
(a) Rearranging the addition law,
\[ P(A\cap B)=P(A)+P(B)-P(A\cup B)=0.45+0.5-0.75=0.2. \]
(b) $P(A'\cap B)=P(B)-P(A\cap B)=0.5-0.2=0.3$.
(c) Since $P(A\cap B)=0.2\neq 0$, the events are \textbf{not} mutually exclusive.
\end{exoresolu}

\begin{attention}[Reading ``only'' and ``neither'' correctly]
\begin{itemize}
\item ``$A$ only'' is \textbf{not} $P(A)$: it is $P(A)-P(A\cap B)$.
\item ``Neither'' is the complement of the \emph{union}, not of the intersection: $P(\text{neither})=1-P(A\cup B)$.
\item The four regions of the diagram must sum to $1$ (or to $n(\mathcal{E})$ when working with counts). Use this as a check.
\end{itemize}
\end{attention}

\begin{exercice}[]
A bag contains $3$ red, $4$ blue and $5$ green counters. One counter is drawn at random. Find the probability that it is (a) red or blue; (b) not green.
\end{exercice}

\begin{exercice}[]
Of $50$ students in an Upper Sixth class in Douala, $30$ study French, $22$ study Computer Science and $12$ study both. A student is chosen at random. Draw a Venn diagram and find the probability that the student studies (a) French or Computer Science; (b) French only; (c) neither subject.
\end{exercice}

\begin{exercice}[]
Given $P(A)=0.55$, $P(B)=0.35$ and $P(A\cup B)=0.7$, find $P(A\cap B)$ and $P(A\cap B')$. Are $A$ and $B$ mutually exclusive? Justify.
\end{exercice}

\begin{corrige}[Exercise 1]
Total $=12$ counters. (a) Red and blue are mutually exclusive: $P=\dfrac{3}{12}+\dfrac{4}{12}=\dfrac{7}{12}$. (b) Complement: $P(\text{not green})=1-\dfrac{5}{12}=\dfrac{7}{12}$.
\end{corrige}

\begin{corrige}[Exercise 2]
Inside outwards: both $=12$; French only $=30-12=18$; Computer Science only $=22-12=10$; neither $=50-(18+12+10)=10$.
(a) $P(F\cup CS)=\dfrac{40}{50}=\dfrac45$ (check: $\tfrac{30}{50}+\tfrac{22}{50}-\tfrac{12}{50}=\tfrac{40}{50}$). (b) $P(F\text{ only})=\dfrac{18}{50}=\dfrac{9}{25}$. (c) $P(\text{neither})=\dfrac{10}{50}=\dfrac15$.
\end{corrige}

\begin{corrige}[Exercise 3]
$P(A\cap B)=0.55+0.35-0.7=0.2$; $P(A\cap B')=0.55-0.2=0.35$. Since $P(A\cap B)=0.2\neq0$, $A$ and $B$ are not mutually exclusive.
\end{corrige}

\begin{aretenir}[Key points of the section]
\begin{itemize}
\item \textbf{Addition law (general):} $P(A\cup B)=P(A)+P(B)-P(A\cap B)$; the subtraction removes the double-counted overlap.
\item \textbf{Mutually exclusive} means $A\cap B=\varnothing$; then $P(A\cup B)=P(A)+P(B)$.
\item \textbf{Complement:} $P(E')=1-P(E)$; for ``at least one'', compute $1-P(\text{none})$.
\item \textbf{Venn diagram:} fill from the intersection outwards; the four regions are $A\cap B$, $A$ only, $B$ only, neither, and they sum to the whole sample space.
\item Never add probabilities without first checking whether the events overlap.
\end{itemize}
\end{aretenir}

\begin{savaistu}[Why Venn diagrams?]
The diagrams are named after the English logician John Venn, who popularised them around 1880 while studying logic. Today they are a standard tool of the GCE Advanced Level examination: many probability questions become almost mechanical once the regions are filled correctly. Master the ``inside-outwards'' method and you will earn these marks quickly and safely.
\end{savaistu}

\coursec{Multiplication Law, Conditional Probability and Independence}

In the previous section we studied the \cle{addition law}, which answers the question ``what is the probability that $A$ \emph{or} $B$ occurs?''. We now turn to the equally important question ``what is the probability that $A$ \emph{and} $B$ occur?''. Answering it leads us to one of the most powerful ideas in all of probability: the notion of \cle{conditional probability}, which measures how the occurrence of one event changes our belief about another. This distinction between events that influence each other and events that do not is at the heart of the \cle{multiplication law} and of the GCE Advanced Level examination.

\courssub{The Multiplication Law for Independent Events}

\begin{experience}[Tossing a coin and rolling a die]
Toss a fair coin and roll a fair die. Does the result of the coin influence the die in any way? Clearly not. The probability of obtaining a Head is $\tfrac12$, and the probability of obtaining a six is $\tfrac16$, whatever the coin shows. Intuitively, the probability of obtaining \emph{both} a Head and a six should be a fraction of a fraction: half of the time the coin gives a Head, and among those cases the die gives a six one sixth of the time. Hence
\[ P(\text{Head and six}) = \frac{1}{2}\times\frac{1}{6} = \frac{1}{12}. \]
We can check this by listing the $12$ equally likely outcomes $(H,1),(H,2),\dots,(T,6)$: exactly one of them is favourable, giving $\tfrac{1}{12}$. The intuition and the enumeration agree.
\end{experience}

This example illustrates events that do not affect each other. Such events are called \cle{independent}.

\begin{definition}[Independent events]
Two events $A$ and $B$ are \textbf{independent} if the occurrence of one does not change the probability of the other. Formally, $A$ and $B$ are independent if and only if
\[ P(A\cap B) = P(A)\times P(B). \]
\end{definition}

\begin{propriete}[Multiplication law — independent case]
If $A$ and $B$ are independent events, then
\[ P(A\cap B) = P(A)\,P(B). \]
More generally, if $A_1, A_2, \dots, A_n$ are mutually independent, then
\[ P(A_1\cap A_2\cap \cdots \cap A_n) = P(A_1)\,P(A_2)\cdots P(A_n). \]
\end{propriete}

\begin{exemplebox}[Two independent machines]
A factory in Douala has two machines that work independently. The probability that machine $A$ breaks down on a given day is $0.1$ and the probability that machine $B$ breaks down is $0.2$. The probability that \emph{both} break down on the same day is
\[ P(A\cap B) = 0.1\times 0.2 = 0.02. \]
The probability that \emph{neither} breaks down is
\[ P(A'\cap B') = P(A')P(B') = 0.9\times 0.8 = 0.72, \]
since $A'$ and $B'$ are also independent.
\end{exemplebox}

\begin{attention}[Do not multiply blindly]
The formula $P(A\cap B)=P(A)P(B)$ is valid \emph{only} for independent events. Applying it to dependent events (for example two cards drawn without replacement) is one of the most frequent errors at the GCE. Always ask yourself first: does the first outcome change the conditions of the second?
\end{attention}

\courssub{Sampling With and Without Replacement}

Dependence arises most naturally in \cle{sampling} problems. A bag contains $5$ red and $3$ blue balls. Two balls are drawn, one after the other.

\textbf{With replacement.} The first ball is put back before the second draw. The bag is restored to its original state, so the two draws are independent:
\[ P(\text{both red}) = \frac{5}{8}\times\frac{5}{8} = \frac{25}{64}. \]

\textbf{Without replacement.} The first ball is kept. If it was red, only $4$ red balls remain among $7$ balls, so the probability of a second red is now $\tfrac47$, not $\tfrac58$:
\[ P(\text{both red}) = \frac{5}{8}\times\frac{4}{7} = \frac{20}{56} = \frac{5}{14}. \]

\begin{attention}[Without replacement, probabilities change]
In sampling \textbf{without replacement}, the composition of the bag (or pack of cards) changes after each draw: both the number of favourable items and the total number of items decrease. The second probability must be computed \emph{given} the result of the first draw. This is exactly the idea of conditional probability, which we now formalise.
\end{attention}

\courssub{Conditional Probability}

Suppose we know that event $B$ has occurred. The sample space is then \emph{restricted} to $B$: only the outcomes inside $B$ remain possible. Within this reduced sample space, the proportion corresponding to $A$ is the share of $A\cap B$ inside $B$.

\begin{definition}[Conditional probability]
Let $A$ and $B$ be events with $P(B)>0$. The \textbf{conditional probability of $A$ given $B$}, written $P(A\mid B)$, is
\[ P(A\mid B) = \frac{P(A\cap B)}{P(B)}. \]
It is the probability that $A$ occurs \emph{knowing that} $B$ has occurred.
\end{definition}

\begin{popfigure}
\begin{center}
\begin{tikzpicture}[scale=1.0]
\draw[thick, rounded corners] (-0.5,-1.6) rectangle (6.6,2.4);
\node at (6.1,2.0) {$\mathcal{E}$};
\draw[fill=popblueL, fill opacity=0.6, thick] (1.6,0.4) circle (1.5);
\draw[fill=poporange, fill opacity=0.35, thick] (3.9,0.4) circle (1.5);
\node at (0.6,0.4) {$A$};
\node at (4.9,0.4) {$B$};
\node at (2.75,0.4) {$A\cap B$};
\draw[->, thick, popdark] (5.6,-1.1) -- (4.6,-0.5);
\node[popdark, align=center] at (5.7,-1.3) {given $B$: only\\ this disc remains};
\end{tikzpicture}
\legende{$P(A\mid B)$ is the proportion of the disc $B$ occupied by $A\cap B$: the sample space is restricted to $B$.}
\end{center}
\end{popfigure}

\begin{exemplebox}[Reading a restricted sample space]
A fair die is rolled. Let $A=\{\text{score is }2\}$ and $B=\{\text{score is even}\}=\{2,4,6\}$. Given that the score is even, only three outcomes remain possible, all equally likely, and one of them is $2$:
\[ P(A\mid B) = \frac{P(A\cap B)}{P(B)} = \frac{1/6}{3/6} = \frac{1}{3}. \]
Notice that $P(A\mid B)=\tfrac13 > P(A)=\tfrac16$: the information ``the score is even'' has increased the probability of $A$.
\end{exemplebox}

\courssub{Conditional Probability from Two-Way Tables}

Examination questions very often present data in a \cle{two-way contingency table}. Conditional probabilities are then read by restricting attention to a single row or column.

\begin{popfigure}
\begin{center}
\begin{tikzpicture}
\matrix (m) [matrix of nodes, nodes={minimum width=2.5cm, minimum height=1cm, draw, anchor=center}, column sep=-\pgflinewidth, row sep=-\pgflinewidth] {
& \textbf{Pass} & \textbf{Fail} & \textbf{Total} \\
\textbf{Boys} & 18 & 12 & 30 \\
\textbf{Girls} & 22 & 8 & 30 \\
\textbf{Total} & 40 & 20 & 60 \\
};
\end{tikzpicture}
\legende{Results of $60$ candidates: e.g.\ $P(\text{pass}\mid\text{boy})=\dfrac{18}{30}$, read entirely within the Boys row.}
\end{center}
\end{popfigure}

\begin{exoresolu}[Conditional probabilities from a table]
The table above summarises the results of $60$ candidates. A candidate is chosen at random. Find (a) $P(\text{pass}\mid\text{boy})$; (b) $P(\text{boy}\mid\text{fail})$; (c) test whether ``passing'' and ``being a boy'' are independent.
\tcblower
\textbf{(a)} Restrict the sample space to the $30$ boys: $18$ of them passed, so
\[ P(\text{pass}\mid\text{boy}) = \frac{18}{30} = \frac{3}{5}. \]
\textbf{(b)} Restrict to the $20$ candidates who failed: $12$ are boys, so
\[ P(\text{boy}\mid\text{fail}) = \frac{12}{20} = \frac{3}{5}. \]
\textbf{(c)} $P(\text{pass}) = \dfrac{18+22}{60} = \dfrac{40}{60} = \dfrac{2}{3}$. Since
\[ P(\text{pass}\mid\text{boy}) = \frac{3}{5} \neq \frac{2}{3} = P(\text{pass}), \]
the two events are \textbf{not independent}: boys pass at a different rate from the group as a whole. (Equivalently, $P(\text{pass}\cap\text{boy})=\tfrac{18}{60}=0.3$ while $P(\text{pass})P(\text{boy})=\tfrac23\times\tfrac12=\tfrac13\approx 0.333$.)
\end{exoresolu}

\courssub{The General Multiplication Law}

Rearranging the definition of conditional probability gives the general rule for the intersection of any two events, dependent or not.

\begin{propriete}[General multiplication law]
For any events $A$ and $B$ with $P(A)>0$,
\[ P(A\cap B) = P(A)\times P(B\mid A). \]
For three events, $P(A\cap B\cap C) = P(A)\,P(B\mid A)\,P(C\mid A\cap B)$. When $A$ and $B$ are independent, $P(B\mid A)=P(B)$ and the law reduces to $P(A\cap B)=P(A)P(B)$.
\end{propriete}

\begin{exemplebox}[Two cards without replacement]
Two cards are drawn successively without replacement from a pack of $52$ cards. The probability that both are aces is
\[ P(\text{1st ace})\times P(\text{2nd ace}\mid\text{1st ace}) = \frac{4}{52}\times\frac{3}{51} = \frac{12}{2652} = \frac{1}{221}. \]
After the first ace is drawn, only $3$ aces remain among $51$ cards: the dependence is fully captured by the conditional probability.
\end{exemplebox}

\courssub{Testing for Independence}

\begin{definition}[Formal test of independence]
Events $A$ and $B$ (with non-zero probabilities) are independent if any one of the following equivalent conditions holds:
\[ P(A\cap B) = P(A)P(B), \qquad P(A\mid B) = P(A), \qquad P(B\mid A) = P(B). \]
In an examination, compute both sides of one of these equalities and compare.
\end{definition}

\begin{attention}[Mutually exclusive is NOT independent]
Do not confuse \cle{mutually exclusive} events ($A\cap B=\varnothing$, so $P(A\cap B)=0$) with \cle{independent} events ($P(A\cap B)=P(A)P(B)$). If $A$ and $B$ are mutually exclusive with $P(A)>0$ and $P(B)>0$, then
\[ P(A\cap B) = 0 \neq P(A)P(B) > 0, \]
so they are \textbf{never independent}. Indeed, if $B$ occurs then $A$ cannot occur: knowing $B$ gives complete information about $A$, which is the opposite of independence.
\end{attention}

\begin{exoresolu}[GCE-style: from $P(A\cup B)$ to $P(A\mid B)$ and a test of independence]
Events $A$ and $B$ are such that $P(A)=0.5$, $P(B)=0.4$ and $P(A\cup B)=0.7$.
(a) Find $P(A\cap B)$. (b) Find $P(A\mid B)$. (c) Determine whether $A$ and $B$ are independent.
\tcblower
\textbf{(a)} By the addition law,
\[ P(A\cap B) = P(A)+P(B)-P(A\cup B) = 0.5+0.4-0.7 = 0.2. \]
\textbf{(b)} By the definition of conditional probability,
\[ P(A\mid B) = \frac{P(A\cap B)}{P(B)} = \frac{0.2}{0.4} = 0.5. \]
\textbf{(c)} Test: $P(A)P(B) = 0.5\times 0.4 = 0.2 = P(A\cap B)$. Therefore $A$ and $B$ \textbf{are independent}. This is confirmed by $P(A\mid B)=0.5=P(A)$.
\end{exoresolu}

\begin{methode}[Tackling a multiplication-law problem]
\begin{enumerate}
\item Identify the events $A$ and $B$ and write down all given probabilities.
\item Ask: are the events independent (separate experiments, sampling with replacement) or dependent (sampling without replacement, related characteristics)?
\item Independent: use $P(A\cap B)=P(A)P(B)$. Dependent: use $P(A\cap B)=P(A)P(B\mid A)$, adjusting the second probability to the new situation.
\item For ``given that'' questions, restrict the sample space and use $P(A\mid B)=\dfrac{P(A\cap B)}{P(B)}$.
\item To test independence, compare $P(A\cap B)$ with $P(A)P(B)$ and conclude in a sentence.
\end{enumerate}
\end{methode}

\begin{exercice}[Practice]
\begin{enumerate}
\item A bag contains $4$ white and $6$ black balls. Two balls are drawn at random. Find the probability that both are white (a) with replacement; (b) without replacement.
\item Given $P(A)=0.6$, $P(B)=0.5$ and $P(A\cup B)=0.8$, find $P(A\mid B)$ and $P(B\mid A)$, and state, with justification, whether $A$ and $B$ are independent.
\item In a school, $45\%$ of students study Physics and $30\%$ study both Physics and Chemistry. Find the probability that a Physics student also studies Chemistry.
\end{enumerate}
\end{exercice}

\begin{corrige}[Exercise 1]
\textbf{1.} (a) With replacement, the draws are independent: $\tfrac{4}{10}\times\tfrac{4}{10}=\tfrac{16}{100}=\tfrac{4}{25}$. (b) Without replacement: $\tfrac{4}{10}\times\tfrac{3}{9}=\tfrac{12}{90}=\tfrac{2}{15}$.
\textbf{2.} $P(A\cap B)=0.6+0.5-0.8=0.3$. Then $P(A\mid B)=\tfrac{0.3}{0.5}=0.6$ and $P(B\mid A)=\tfrac{0.3}{0.6}=0.5$. Since $P(A\cap B)=0.3=0.6\times0.5=P(A)P(B)$, the events are independent.
\textbf{3.} $P(\text{Chem}\mid\text{Phys})=\dfrac{P(\text{Chem}\cap\text{Phys})}{P(\text{Phys})}=\dfrac{0.30}{0.45}=\dfrac{2}{3}$.
\end{corrige}

\begin{aretenir}[Key points of this section]
\begin{itemize}
\item Independent events: $P(A\cap B)=P(A)P(B)$; the occurrence of one does not affect the other.
\item Conditional probability: $P(A\mid B)=\dfrac{P(A\cap B)}{P(B)}$ — restrict the sample space to $B$.
\item General multiplication law: $P(A\cap B)=P(A)P(B\mid A)$, essential for sampling without replacement.
\item Test of independence: compare $P(A\cap B)$ with $P(A)P(B)$, or $P(A\mid B)$ with $P(A)$.
\item Mutually exclusive events with positive probabilities are never independent.
\end{itemize}
\end{aretenir}

The multiplication law and conditional probability are the engine behind \cle{tree diagrams}, which we study in the next section as a systematic tool for multi-stage experiments.

\coursec{Tree Diagrams and Multi-Stage Experiments}

In the previous section, we mastered the Multiplication Law $P(A \cap B) = P(A) \times P(B|A)$ and the concept of conditional probability. These tools are algebraic; they tell us \emph{how} to calculate, but for experiments with several stages, keeping track of all the conditional probabilities can become confusing. A \cle{probability tree diagram} (or simply a \cle{tree diagram}) provides a visual map of a multi-stage experiment. It organises the sample space, displays the conditional probabilities on the branches, and reduces the calculation of compound probabilities to following paths. This is an indispensable tool for GCE Advanced Level problems, especially those involving sampling without replacement, repeated trials, or reverse conditional probabilities (which lead naturally to Bayes' theorem).

\courssub{Constructing Probability Trees for Two-Stage Experiments}

Consider a two-stage experiment: drawing two balls, one after the other, from a bag containing 5 red and 3 blue balls. The outcome of the first draw affects the probabilities for the second draw (sampling \emph{without replacement}). A tree diagram represents this perfectly.

\begin{definition}[Probability Tree Diagram]
A \cle{probability tree diagram} is a graphical representation of a multi-stage random experiment.
\begin{itemize}
    \item Each \cle{node} (junction point) represents a stage of the experiment.
    \item Each \cle{branch} represents a possible outcome at that stage.
    \item The \cle{probability} of that outcome, \emph{conditional on the path taken to reach that node}, is written on the branch.
    \item The \cle{final outcomes} (elementary events of the combined experiment) are found at the ends of the terminal branches (the \emph{leaves}).
\end{itemize}
\end{definition}

\begin{propriete}[Fundamental Rules of a Probability Tree]
Let a node represent a stage where outcomes $B_1, B_2, \dots, B_k$ are possible, given the history leading to that node.
\begin{enumerate}
    \item \textbf{Sum Rule:} The probabilities on the branches leaving a single node sum to 1.
    \[ \sum_{i=1}^k P(B_i \mid \text{history}) = 1 \]
    \item \cle{Multiplication Rule (Along a Path):} The probability of reaching a specific final outcome (a complete path from root to leaf) is the \textbf{product} of the probabilities on the branches of that path.
    \[ P(\text{Path}) = P(\text{1st branch}) \times P(\text{2nd branch}) \times \cdots \]
    \item \cle{Addition Rule (Across Paths):} The probability of a compound event (a set of final outcomes) is the \textbf{sum} of the probabilities of the individual paths that constitute the event.
\end{enumerate}
\textbf{Justification:} The Multiplication Rule along a path is exactly the General Multiplication Law $P(A \cap B) = P(A)\,P(B|A)$ applied sequentially: $P(A \cap B \cap C) = P(A)\,P(B|A)\,P(C|A \cap B)$. The Sum Rule at a node reflects the fact that the outcomes at that stage form a partition of the (conditional) sample space, so their probabilities are exhaustive and mutually exclusive.
\end{propriete}

\begin{methode}[Building and Reading a Probability Tree]
\begin{enumerate}
    \item \textbf{Identify the stages.} How many sequential actions? (e.g., Draw 1, Draw 2).
    \item \textbf{Draw the first-level branches.} List all possible outcomes of the first stage. Label the branches with their (unconditional) probabilities.
    \item \textbf{Draw subsequent levels.} From the end of \emph{each} first-stage branch, draw branches for the second-stage outcomes. \textbf{Crucial:} label these branches with \emph{conditional} probabilities, assuming the first-stage outcome has happened.
    \item \textbf{Check the Sum Rule.} At every node, ensure the outgoing branch probabilities add to 1.
    \item \textbf{Calculate path probabilities.} Multiply along each complete path to the end and write the result at the leaf. The sum of all leaf probabilities must be 1.
    \item \textbf{Answer the question.} Identify which leaves correspond to the required event and add their probabilities.
\end{enumerate}
\end{methode}

\begin{exoresolu}[Two Balls Drawn Without Replacement]
A bag contains 5 Red (R) and 3 Blue (B) balls. Two balls are drawn at random, one after the other, \textbf{without replacement}. Construct the probability tree and find:
\begin{enumerate}
    \item $P(\text{both Red})$;
    \item $P(\text{one ball of each colour})$;
    \item $P(\text{the second ball is Red})$.
\end{enumerate}
\tcblower
\textbf{Tree diagram:} see the figure below. First draw: $P(R_1)=\frac{5}{8}$, $P(B_1)=\frac{3}{8}$. After a Red is drawn, 4 Red and 3 Blue remain out of 7; after a Blue is drawn, 5 Red and 2 Blue remain out of 7.

\textbf{1.} $P(\text{both Red}) = P(R_1 \cap R_2) = \dfrac{5}{8} \times \dfrac{4}{7} = \dfrac{20}{56} = \dfrac{5}{14}$.

\textbf{2.} $P(\text{one of each}) = P(R_1 \cap B_2) + P(B_1 \cap R_2) = \dfrac{5}{8}\times\dfrac{3}{7} + \dfrac{3}{8}\times\dfrac{5}{7} = \dfrac{15}{56} + \dfrac{15}{56} = \dfrac{30}{56} = \dfrac{15}{28}$.

\textbf{3.} $P(\text{second is Red}) = P(R_1 \cap R_2) + P(B_1 \cap R_2) = \dfrac{20}{56} + \dfrac{15}{56} = \dfrac{35}{56} = \dfrac{5}{8}$.

\emph{Observation:} $P(\text{second is Red}) = P(\text{first is Red}) = \frac{5}{8}$. In sampling without replacement, the \emph{unconditional} probability of drawing a given colour is the same at every draw --- a useful check.
\end{exoresolu}

\begin{popfigure}
\begin{tikzpicture}[
    grow=right,
    level distance=3.2cm,
    level 1/.style={sibling distance=3.4cm},
    level 2/.style={sibling distance=1.7cm},
    edge from parent/.style={draw, thick},
    every node/.style={font=\small}
]
\node[circle, draw, inner sep=2pt, fill=popblueL] {Start}
    child { node[circle, draw, inner sep=2pt, fill=popgreenL] (B1) {B}
        edge from parent node[below, pos=0.5] {$\frac{3}{8}$}
        child { node[rectangle, draw, inner sep=2pt, fill=popgreenL] (B2b) {B}
            edge from parent node[below, pos=0.5] {$\frac{2}{7}$}
        }
        child { node[rectangle, draw, inner sep=2pt, fill=popgreenL] (R2b) {R}
            edge from parent node[above, pos=0.5] {$\frac{5}{7}$}
        }
    }
    child { node[circle, draw, inner sep=2pt, fill=popgreenL] (R1) {R}
        edge from parent node[above, pos=0.5] {$\frac{5}{8}$}
        child { node[rectangle, draw, inner sep=2pt, fill=popgreenL] (B2) {B}
            edge from parent node[below, pos=0.5] {$\frac{3}{7}$}
        }
        child { node[rectangle, draw, inner sep=2pt, fill=popgreenL] (R2) {R}
            edge from parent node[above, pos=0.5] {$\frac{4}{7}$}
        }
    };
\node[right=4pt of R2, align=left] {RR:\ $\frac{5}{8}\times\frac{4}{7}=\frac{20}{56}$};
\node[right=4pt of B2, align=left] {RB:\ $\frac{5}{8}\times\frac{3}{7}=\frac{15}{56}$};
\node[right=4pt of R2b, align=left] {BR:\ $\frac{3}{8}\times\frac{5}{7}=\frac{15}{56}$};
\node[right=4pt of B2b, align=left] {BB:\ $\frac{3}{8}\times\frac{2}{7}=\frac{6}{56}$};
\end{tikzpicture}
\legende{Probability tree for drawing two balls without replacement from 5 Red and 3 Blue. Second-level probabilities are conditional on the first draw; leaf probabilities sum to $\frac{20+15+15+6}{56}=1$.}
\end{popfigure}

\begin{attention}[Without Replacement: Adjust Numerator AND Denominator!]
The most common error in tree diagrams is forgetting to update the contents of the bag on the second-stage branches when sampling is \textbf{without replacement}.
\begin{itemize}
    \item \textbf{With replacement:} second-level probabilities are identical to first-level probabilities (independent stages).
    \item \textbf{Without replacement:} the total decreases by 1, and the count of the colour already drawn decreases by 1.
\end{itemize}
Always ask: ``Given what happened on the first branch, what is in the bag \emph{now}?''
\end{attention}

\courssub{Three-Stage Trees and `At Least One' Problems}

Tree diagrams extend naturally to three or more stages. The rules remain identical: multiply along branches, add across final outcomes. A classic GCE problem involves repeated trials (best-of-three matches, repeated tosses, inspections) where we need the probability of \emph{at least one} success.

\begin{methode}[`At Least One' via the Complement on a Tree]
For an event ``at least one success in $n$ trials'':
\begin{enumerate}
    \item Draw the tree (or just the relevant part).
    \item Identify the \textbf{single path} corresponding to ``no success at all'' (failure on every trial).
    \item Calculate the probability of that single path, $P(\text{all failures})$.
    \item Use the complement: $P(\text{at least one success}) = 1 - P(\text{all failures})$.
\end{enumerate}
This is almost always faster than adding the probabilities of all paths containing at least one success.
\end{methode}

\begin{exoresolu}[Best-of-Three Tournament]
Two players, A and B, play a series of independent games; the first to win 2 games wins the match. The probability that A wins any single game is $0.6$. Find the probability that:
\begin{enumerate}
    \item A wins the match;
    \item the match lasts exactly 3 games.
\end{enumerate}
\tcblower
The match stops as soon as a player has 2 wins, so some paths have length 2 and others length 3. The terminal paths and their probabilities are:
\begin{itemize}
    \item AA: $0.6 \times 0.6 = 0.36$ \quad (A wins the match)
    \item ABA: $0.6 \times 0.4 \times 0.6 = 0.144$ \quad (A wins the match)
    \item BAA: $0.4 \times 0.6 \times 0.6 = 0.144$ \quad (A wins the match)
    \item ABB: $0.6 \times 0.4 \times 0.4 = 0.096$ \quad (B wins the match)
    \item BAB: $0.4 \times 0.6 \times 0.4 = 0.096$ \quad (B wins the match)
    \item BB: $0.4 \times 0.4 = 0.16$ \quad (B wins the match)
\end{itemize}
Check: $0.36 + 0.144 + 0.144 + 0.096 + 0.096 + 0.16 = 1$. \checkmark

\textbf{1.} $P(\text{A wins match}) = 0.36 + 0.144 + 0.144 = \mathbf{0.648}$.

\textbf{2.} $P(\text{match lasts 3 games}) = 0.144 + 0.144 + 0.096 + 0.096 = \mathbf{0.48}$.
\end{exoresolu}

\begin{popfigure}
\begin{tikzpicture}[
    grow=right,
    level distance=2.6cm,
    level 1/.style={sibling distance=4.4cm},
    level 2/.style={sibling distance=2.4cm},
    level 3/.style={sibling distance=1.2cm},
    edge from parent/.style={draw, thick},
    every node/.style={font=\small}
]
\node[circle, draw, inner sep=2pt, fill=popblueL] {Start}
    child { node[circle, draw, inner sep=2pt, fill=popgreenL] (B1) {B}
        edge from parent node[below, pos=0.5] {$0.4$}
        child { node[rectangle, draw, inner sep=2pt, fill=popgreenL] (B2b) {B}
            edge from parent node[below, pos=0.5] {$0.4$}
        }
        child { node[circle, draw, inner sep=2pt, fill=popgreenL] (A2b) {A}
            edge from parent node[above, pos=0.5] {$0.6$}
            child { node[rectangle, draw, inner sep=2pt, fill=popgreenL] (B3b) {B}
                edge from parent node[below, pos=0.5] {$0.4$}
            }
            child { node[rectangle, draw, inner sep=2pt, fill=popgreenL] (A3b) {A}
                edge from parent node[above, pos=0.5] {$0.6$}
            }
        }
    }
    child { node[circle, draw, inner sep=2pt, fill=popgreenL] (A1) {A}
        edge from parent node[above, pos=0.5] {$0.6$}
        child { node[circle, draw, inner sep=2pt, fill=popgreenL] (B2a) {B}
            edge from parent node[below, pos=0.5] {$0.4$}
            child { node[rectangle, draw, inner sep=2pt, fill=popgreenL] (B3a) {B}
                edge from parent node[below, pos=0.5] {$0.4$}
            }
            child { node[rectangle, draw, inner sep=2pt, fill=popgreenL] (A3) {A}
                edge from parent node[above, pos=0.5] {$0.6$}
            }
        }
        child { node[rectangle, draw, inner sep=2pt, fill=popgreenL] (A2) {A}
            edge from parent node[above, pos=0.5] {$0.6$}
        }
    };
\node[right=4pt of A2] {AA:\ $0.36$};
\node[right=4pt of A3] {ABA:\ $0.144$};
\node[right=4pt of B3a] {ABB:\ $0.096$};
\node[right=4pt of A3b] {BAA:\ $0.144$};
\node[right=4pt of B3b] {BAB:\ $0.096$};
\node[right=4pt of B2b] {BB:\ $0.16$};
\end{tikzpicture}
\legende{Tree for a best-of-three match. The match stops early when a player reaches 2 wins (paths AA and BB have length 2).}
\end{popfigure}

\courssub{Conditional Probability from Tree Diagrams}

One of the most powerful uses of a tree diagram is finding \emph{reverse} conditional probabilities: $P(\text{first-stage event} \mid \text{second-stage event})$. The tree naturally gives $P(\text{first} \cap \text{second})$ (a path probability) and $P(\text{second})$ (a sum of path probabilities). The conditional probability is then their ratio.

\begin{propriete}[Conditional Probability on a Tree]
For a two-stage experiment with first-stage outcomes $A_i$ and second-stage outcomes $B_j$:
\[ P(A_i \mid B_j) = \frac{P(A_i \cap B_j)}{P(B_j)} = \frac{\text{probability of the path } (A_i \to B_j)}{\text{sum of probabilities of all paths ending in } B_j}. \]
Visually: \textbf{restrict attention to the leaves corresponding to $B_j$.} The conditional probability $P(A_i \mid B_j)$ is the proportion of that group's total probability contributed by the path through $A_i$.
\end{propriete}

\begin{exoresolu}[Reverse Conditional Probability]
Using the bag with 5 Red and 3 Blue balls drawn without replacement (first example), find the probability that the \textbf{first} ball was Red, given that the \textbf{second} ball is Red; that is, find $P(R_1 \mid R_2)$.
\tcblower
From the tree:
\begin{itemize}
    \item $P(R_1 \cap R_2) = \dfrac{5}{8} \times \dfrac{4}{7} = \dfrac{20}{56}$;
    \item $P(B_1 \cap R_2) = \dfrac{3}{8} \times \dfrac{5}{7} = \dfrac{15}{56}$;
    \item $P(R_2) = \dfrac{20}{56} + \dfrac{15}{56} = \dfrac{35}{56}$.
\end{itemize}
\[ P(R_1 \mid R_2) = \frac{P(R_1 \cap R_2)}{P(R_2)} = \frac{20/56}{35/56} = \frac{20}{35} = \frac{4}{7}. \]
\textbf{Interpretation:} given that the second ball is Red, we look only at the two paths ending in Red. The first ball was Red in $\frac{20}{35} = \frac{4}{7}$ of that probability mass.
\end{exoresolu}

\courssub{Real-Life Models and Forming Equations}

Tree diagrams model many real-world sequential processes: quality control (inspecting items one by one), medical screening (test, then confirmatory test), weather forecasting (rain today $\to$ rain tomorrow), or multi-stage interviews. A favourite GCE technique gives an \emph{overall} probability and asks for an unknown branch probability: form an equation from the tree and solve it.

\begin{exoresolu}[Quality Control: Finding an Unknown Probability]
A machine produces components independently, each with probability $d$ of being defective. Components are packed in boxes of 2, and a box is rejected if \emph{at least one} component is defective. The probability that a box is rejected is $0.19$. Find $d$.
\tcblower
\textbf{Model:} Stage 1: component 1 --- $D$ with probability $d$, $G$ with probability $1-d$. Stage 2: component 2 --- same probabilities (independence). The tree has four paths: $DD, DG, GD, GG$.
``Box accepted'' $=$ no defective $=$ path $GG$ only.
\begin{align*}
P(\text{accepted}) &= P(GG) = (1-d)^2,\\
P(\text{rejected}) &= 1 - (1-d)^2.
\end{align*}
Given $P(\text{rejected}) = 0.19$:
\begin{align*}
1 - (1-d)^2 &= 0.19\\
(1-d)^2 &= 0.81\\
1-d &= 0.9 \quad (\text{taking the positive root, since } 0 \le d \le 1)\\
d &= \mathbf{0.1}.
\end{align*}
\end{exoresolu}

\begin{attention}[Multiply Along a Path, Add Across Paths]
A persistent error is \textbf{adding} probabilities along a single path instead of \textbf{multiplying} them.
\begin{itemize}
    \item \textbf{Along a path (AND):} multiply. ``First Red \textbf{and} second Red'' $\to \frac{5}{8} \times \frac{4}{7}$.
    \item \textbf{Across paths (OR):} add. ``RR \textbf{or} BR'' $\to \frac{20}{56} + \frac{15}{56}$.
\end{itemize}
If you ever write $\frac{5}{8} + \frac{4}{7}$ for a two-stage event, stop: you are mixing stages. Probabilities along a path are conditional and must be multiplied.
\end{attention}

\begin{exoresolu}[Weather Model: Dependent Days]
In a certain region, if it rains on a given day, the probability it rains the next day is $0.7$; if it is fine on a given day, the probability it rains the next day is $0.3$. The probability it rains on Monday is $0.4$. Find the probability that it rains on Wednesday.
\tcblower
Let $R$ = rain, $F$ = fine. The transition probabilities are:
$P(R_{\text{next}}|R) = 0.7$, $P(F_{\text{next}}|R) = 0.3$, $P(R_{\text{next}}|F) = 0.3$, $P(F_{\text{next}}|F) = 0.7$.

The paths from Monday to rain on Wednesday:
\begin{enumerate}
    \item $R \to R \to R$: $0.4 \times 0.7 \times 0.7 = 0.196$
    \item $R \to F \to R$: $0.4 \times 0.3 \times 0.3 = 0.036$
    \item $F \to R \to R$: $0.6 \times 0.3 \times 0.7 = 0.126$
    \item $F \to F \to R$: $0.6 \times 0.7 \times 0.3 = 0.126$
\end{enumerate}
\[ P(R_{\text{Wed}}) = 0.196 + 0.036 + 0.126 + 0.126 = \mathbf{0.484}. \]
\end{exoresolu}

\courssub{Synthesis and Exam Checklist}

\begin{aretenir}[Mastering Tree Diagrams for the GCE]
\begin{itemize}
    \item \textbf{Draw it.} Even when not explicitly asked, sketch a mini-tree: it prevents errors in identifying conditional probabilities.
    \item \textbf{Label clearly.} Write probabilities \emph{on the branches}. Distinguish $P(B)$ (first level) from $P(B|A)$ (second level).
    \item \textbf{Check sums.} At every node, outgoing probabilities sum to 1; all leaf probabilities sum to 1.
    \item \textbf{Without replacement:} update numerators \textbf{and} denominators on later branches.
    \item \textbf{`At least one':} use $1 - P(\text{none})$; find the single ``all failures'' path.
    \item \textbf{Reverse conditional:} $P(\text{cause}|\text{effect}) = \dfrac{P(\text{cause} \cap \text{effect})}{P(\text{effect})}$; on the tree, (path probability) $\div$ (total of the relevant leaves).
    \item \textbf{Unknown probability:} call it $p$; express path probabilities in terms of $p$; use the given total to form an equation; solve; check $0 \le p \le 1$.
\end{itemize}
\end{aretenir}

\begin{exercice}[Consolidation Exercise]
\begin{enumerate}
    \item A box contains 4 black pens and 6 blue pens. Two pens are taken at random without replacement. Draw the tree diagram and find the probability that (a) both are black, (b) they are of different colours, (c) the second is black given that the first is blue.
    \item In a driving test, the probability of passing the theory part is $0.8$. A candidate who passes theory then takes the practical part, which they pass with probability $0.7$. A candidate who fails theory cannot take the practical. Draw the tree and find the probability that a randomly chosen candidate obtains the licence (passes both parts).
    \item A biased coin has $P(H) = p$. It is tossed 3 times independently. The probability of obtaining \emph{at least one} Head is $0.973$. Form an equation in $p$ and solve it.
    \item (GCE style) A student travels to school by bus or on foot. If he takes the bus one day, the probability he takes the bus the next day is $0.7$. If he walks one day, the probability he walks the next day is $0.4$. On Monday he takes the bus. Find the probability that he walks on Wednesday.
\end{enumerate}
\end{exercice}

\begin{corrige}[Exercise 1]
First level: Bk $\left(\frac{4}{10}\right)$, Bl $\left(\frac{6}{10}\right)$.
Second level from Bk: Bk $\left(\frac{3}{9}\right)$, Bl $\left(\frac{6}{9}\right)$.
Second level from Bl: Bk $\left(\frac{4}{9}\right)$, Bl $\left(\frac{5}{9}\right)$.
\begin{enumerate}
    \item[(a)] $P(\text{Bk,Bk}) = \dfrac{4}{10}\times\dfrac{3}{9} = \dfrac{12}{90} = \dfrac{2}{15}$.
    \item[(b)] $P(\text{different}) = \dfrac{4}{10}\times\dfrac{6}{9} + \dfrac{6}{10}\times\dfrac{4}{9} = \dfrac{24+24}{90} = \dfrac{48}{90} = \dfrac{8}{15}$.
    \item[(c)] $P(\text{Bk}_2 \mid \text{Bl}_1) = \dfrac{4}{9}$ (read directly from the branch leaving the Bl node).
\end{enumerate}
\end{corrige}

\begin{corrige}[Exercise 2]
Tree: Theory Pass $(0.8) \to$ Practical Pass $(0.7)$ or Practical Fail $(0.3)$; Theory Fail $(0.2) \to$ no practical (licence failed).
\[ P(\text{licence}) = P(\text{pass theory}) \times P(\text{pass practical} \mid \text{pass theory}) = 0.8 \times 0.7 = \mathbf{0.56}. \]
\end{corrige}

\begin{corrige}[Exercise 3]
``At least one Head'' is the complement of ``no Heads'', i.e. three Tails, a single path of probability $(1-p)^3$.
\begin{align*}
1 - (1-p)^3 &= 0.973\\
(1-p)^3 &= 0.027\\
1-p &= \sqrt[3]{0.027} = 0.3\\
p &= \mathbf{0.7}.
\end{align*}
Check: $0 \le 0.7 \le 1$. \checkmark
\end{corrige}

\begin{corrige}[Exercise 4]
Monday: Bus with probability $1$.
Tuesday: Bus $(0.7)$, Walk $(0.3)$.
From Tuesday Bus: Wednesday Walk $(0.3)$. From Tuesday Walk: Wednesday Walk $(0.4)$.
Paths ending in ``Walk on Wednesday'':
\begin{enumerate}
    \item Bus $\to$ Bus $\to$ Walk: $1 \times 0.7 \times 0.3 = 0.21$;
    \item Bus $\to$ Walk $\to$ Walk: $1 \times 0.3 \times 0.4 = 0.12$.
\end{enumerate}
\[ P(\text{walks on Wednesday}) = 0.21 + 0.12 = \mathbf{0.33}. \]
\end{corrige}

\coursec{Probability with Cards, Dice, Coins and Counting Techniques}

In the previous section we mastered tree diagrams, which handle \emph{sequential} experiments beautifully. But many GCE questions involve a different kind of situation: a hand of cards dealt \emph{all at once}, a committee \emph{selected} from a group, two dice thrown \emph{simultaneously}. Here there is no natural ``first branch, second branch''. Instead, we count. This section combines the classical definition $P(E)=\dfrac{\text{number of favourable outcomes}}{\text{total number of outcomes}}$ with the powerful counting tools you met in the chapter on permutations and combinations. The standard objects --- coins, dice and the 52-card pack --- appear in almost every GCE paper, so we begin by fixing their exact structure.

\courssub{Standard Facts: Coins, Dice and the Pack of 52 Cards}

\begin{definition}[Standard objects of probability]
\begin{itemize}
\item A \cle{fair coin} has two equally likely faces: Head (H) and Tail (T).
\item A \cle{fair die} is a cube whose faces are numbered $1,2,3,4,5,6$, each equally likely. Two opposite faces always sum to $7$.
\item A \cle{standard pack of cards} contains $52$ cards, divided into $4$ \cle{suits} of $13$ \cle{ranks} each.
\end{itemize}
\end{definition}

\begin{aretenir}[Structure of a standard 52-card pack]
\begin{itemize}
\item \textbf{4 suits}: Spades ($\spadesuit$, black), Clubs ($\clubsuit$, black), Hearts ($\heartsuit$, red), Diamonds ($\diamondsuit$, red). So $26$ red cards and $26$ black cards.
\item \textbf{13 ranks per suit}: $2,3,4,5,6,7,8,9,10$, Jack (J), Queen (Q), King (K), Ace (A).
\item \textbf{Face (court) cards}: J, Q, K of each suit, giving $12$ face cards in total. The Ace is \emph{not} a face card.
\item \textbf{Aces}: $4$ in total. Picture cards (J, Q, K): $12$. Number cards ($2$ to $10$): $36$.
\end{itemize}
\end{aretenir}

\begin{popfigure}
\begin{center}
\begin{tikzpicture}[scale=0.62, every node/.style={font=\scriptsize}]
\foreach \i/\rank in {0/A,1/2,2/3,3/4,4/5,5/6,6/7,7/8,8/9,9/10,10/J,11/Q,12/K}{
  \draw[fill=popink, draw=popink] (\i*0.85,3) rectangle (\i*0.85+0.7,4);
  \node[white] at (\i*0.85+0.35,3.5) {\rank};
  \draw[fill=popdark, draw=popdark] (\i*0.85,1.5) rectangle (\i*0.85+0.7,2.5);
  \node[white] at (\i*0.85+0.35,2) {\rank};
  \draw[fill=poporange, draw=poporange] (\i*0.85,0) rectangle (\i*0.85+0.7,1);
  \node[white] at (\i*0.85+0.35,0.5) {\rank};
  \draw[fill=poppink, draw=poppink] (\i*0.85,-1.5) rectangle (\i*0.85+0.7,-0.5);
  \node[white] at (\i*0.85+0.35,-1) {\rank};
}
\node[anchor=west, font=\small\bfseries] at (11.6,3.5) {Spades (black)};
\node[anchor=west, font=\small\bfseries] at (11.6,2) {Clubs (black)};
\node[anchor=west, font=\small\bfseries] at (11.6,0.5) {Hearts (red)};
\node[anchor=west, font=\small\bfseries] at (11.6,-1) {Diamonds (red)};
\node[anchor=west, font=\small] at (11.6,-2.3) {13 ranks per suit};
\node[anchor=west, font=\small] at (11.6,-2.9) {J, Q, K = face cards};
\end{tikzpicture}
\end{center}
\legende{The structure of a standard pack: 4 suits $\times$ 13 ranks $= 52$ cards; two black suits and two red suits.}
\end{popfigure}

\begin{exemplebox}[Single-draw probabilities from a pack]
One card is drawn at random from a well-shuffled pack. Then
\[ P(\text{ace})=\frac{4}{52}=\frac{1}{13}, \qquad P(\text{red face card})=\frac{6}{52}=\frac{3}{26}, \qquad P(\text{heart or king})=\frac{13+4-1}{52}=\frac{16}{52}=\frac{4}{13}, \]
where in the last computation we subtracted the King of Hearts, counted twice (addition law).
\end{exemplebox}

\begin{exoresolu}[Two cards drawn without replacement]
Two cards are drawn at random, one after the other \emph{without} replacement, from a standard pack. Find the probability that (a) both are aces; (b) the first is a spade and the second is red.
\tcblower
\textbf{(a)} After the first ace is drawn, $3$ aces remain out of $51$ cards:
\[ P(\text{both aces})=\frac{4}{52}\times\frac{3}{51}=\frac{12}{2652}=\frac{1}{221}. \]
\textbf{(b)} The first card is a spade (hence black), so all $26$ red cards remain in the pack of $51$:
\[ P=\frac{13}{52}\times\frac{26}{51}=\frac{1}{4}\times\frac{26}{51}=\frac{13}{102}. \]
Note how the multiplication law handles the changing sample space --- exactly the logic of a tree diagram, compressed into one line.
\end{exoresolu}

\courssub{Combined Experiments: Two Dice, a Die and a Coin}

When a die and a coin are thrown together, the sample space has $6\times 2 = 12$ equally likely outcomes $(n, F)$ with $n\in\{1,\dots,6\}$ and $F\in\{H,T\}$. When \emph{two dice} are thrown, there are $6\times 6 = 36$ equally likely \emph{ordered pairs} $(i,j)$. This is the crucial point: the 36 pairs are equally likely, but the 11 possible sums $2,3,\dots,12$ are \emph{not}.

\begin{attention}[Sums of two dice are NOT equally likely]
There is only one way to obtain a sum of $2$, namely $(1,1)$, but six ways to obtain a sum of $7$: $(1,6),(2,5),(3,4),(4,3),(5,2),(6,1)$. Hence $P(\text{sum}=7)=\dfrac{6}{36}=\dfrac16$ while $P(\text{sum}=2)=\dfrac{1}{36}$. Never write $P(\text{sum}=7)=\dfrac{1}{11}$. Always count \emph{ordered pairs}, not sums.
\end{attention}

\begin{popfigure}
\begin{center}
\begin{tikzpicture}[scale=0.72, every node/.style={font=\scriptsize}]
\foreach \i in {1,...,6}{
  \foreach \j in {1,...,6}{
    \pgfmathparse{(\i+\j>=9) ? 1 : 0}
    \ifnum\pgfmathresult=1
      \draw[fill=popblueL, draw=gray] (\i-1,\j-1) rectangle (\i,\j);
    \else
      \draw[fill=white, draw=gray] (\i-1,\j-1) rectangle (\i,\j);
    \fi
    \pgfmathparse{int(\i+\j)}
    \node at (\i-0.5,\j-0.5) {\pgfmathresult};
  }
}
\foreach \i in {1,...,6}{
  \node at (\i-0.5,-0.4) {\i};
  \node at (-0.4,\i-0.5) {\i};
}
\node[rotate=90] at (-1,2.5) {Die 1};
\node at (2.5,-1) {Die 2};
\end{tikzpicture}
\end{center}
\legende{The 36 equally likely outcomes of two dice; each cell shows the sum. Shaded cells: the event ``sum $\geq 9$'' (10 cells), so $P(\text{sum}\geq 9)=\frac{10}{36}=\frac{5}{18}$. The diagonal cells $(i,i)$ are the 6 doubles.}
\end{popfigure}

\begin{exemplebox}[Reading the two-dice grid]
From the grid: $P(\text{sum}\geq 9)=\dfrac{10}{36}=\dfrac{5}{18}$; $P(\text{a double})=\dfrac{6}{36}=\dfrac16$; $P(\text{sum}=7)=\dfrac{6}{36}=\dfrac16$; for $P(|\text{difference}|=2)$ the favourable pairs are $(1,3),(2,4),(3,5),(4,6)$ and their reversals, $8$ pairs in all, so $P=\dfrac{8}{36}=\dfrac29$.
\end{exemplebox}

\begin{exoresolu}[A die and a coin]
A fair die and a fair coin are thrown together. Find the probability of obtaining (a) a head and an even number; (b) a tail or a number greater than 4.
\tcblower
The 12 outcomes are equally likely.
\textbf{(a)} Favourable: $(2,H),(4,H),(6,H)$, so $P=\dfrac{3}{12}=\dfrac14$. Alternatively, by independence of the die and the coin: $\dfrac12\times\dfrac36=\dfrac14$.
\textbf{(b)} $P(T)=\dfrac{6}{12}$, $P(n>4)=\dfrac{4}{12}$, $P\bigl(T\cap\{n>4\}\bigr)=\dfrac{2}{12}$. By the addition law:
\[ P=\frac{6}{12}+\frac{4}{12}-\frac{2}{12}=\frac{8}{12}=\frac23. \]
\end{exoresolu}

\courssub{Permutations or Combinations? The Revision Link}

Suppose a committee of $3$ is chosen from $10$ people. Does the order of selection matter? No: the committee $\{\text{Ali, Bih, Chen}\}$ is the same however we list its members. But if we choose a \emph{chairperson, secretary and treasurer}, then order matters: Ali as chair and Bih as secretary is different from Bih as chair and Ali as secretary.

\begin{aretenir}[The two counting tools]
\begin{itemize}
\item \cle{Combination} (order does not matter): the number of ways of choosing $r$ objects from $n$ distinct objects is
\[ \binom{n}{r} = {}^{n}C_{r} = \frac{n!}{r!(n-r)!}. \]
\item \cle{Permutation} (order matters): the number of ordered arrangements of $r$ objects chosen from $n$ distinct objects is
\[ {}^{n}P_{r} = \frac{n!}{(n-r)!}. \]
\end{itemize}
\end{aretenir}

\begin{methode}[Choosing between ${}^nP_r$ and ${}^nC_r$ in probability]
\begin{enumerate}
\item Ask: ``If I swap two of the chosen objects, do I get a \emph{different} outcome of the experiment?'' If \textbf{yes}, use permutations; if \textbf{no}, use combinations.
\item Selections (committees, hands of cards, groups of objects drawn together): use ${}^nC_r$.
\item Arrangements (words from letters, codes from digits, rankings, ordered draws): use ${}^nP_r$ or the multiplication principle.
\item \textbf{Golden rule}: use the \emph{same} type of counting in the numerator and in the denominator.
\end{enumerate}
\end{methode}

\begin{attention}[Never mix ${}^nP_r$ and ${}^nC_r$ in the same fraction]
If the denominator counts \emph{unordered} selections (combinations), the numerator must also count unordered selections. Mixing ${}^nP_r$ on top with ${}^nC_r$ below does not give a number between $0$ and $1$ in general and is a classic source of wrong answers. Consistency guarantees that each favourable outcome corresponds to exactly one counted object on each side of the fraction.
\end{attention}

\begin{propriete}[Probability by counting]
If an experiment consists of selecting $r$ objects from $n$ distinct objects, all selections being equally likely, then for any event $E$,
\[ P(E)=\frac{\text{number of favourable selections}}{\text{total number of selections}}. \]
\end{propriete}

\courssub{Worked Examples with Counting Techniques}

\begin{exoresolu}[A committee containing at least one woman]
A committee of $4$ people is to be chosen at random from $6$ men and $5$ women. Find the probability that the committee contains (a) exactly two women; (b) at least one woman.
\tcblower
Total number of committees (order irrelevant): $\dbinom{11}{4}=330$.
\textbf{(a)} Choose $2$ women from $5$ and $2$ men from $6$:
\[ P(\text{exactly 2 women})=\frac{\dbinom{5}{2}\dbinom{6}{2}}{\dbinom{11}{4}}=\frac{10\times 15}{330}=\frac{150}{330}=\frac{5}{11}. \]
\textbf{(b)} ``At least one woman'' means $1, 2, 3$ or $4$ women --- four cases. The complement ``no woman'' (all $4$ chosen from the $6$ men) is a single case:
\[ P(\text{at least one woman})=1-\frac{\dbinom{6}{4}}{\dbinom{11}{4}}=1-\frac{15}{330}=1-\frac{1}{22}=\frac{21}{22}. \]
Complementary counting turned four computations into one.
\end{exoresolu}

\begin{methode}[Complementary counting for ``at least one'' problems]
\begin{enumerate}
\item Identify the complement of the event: ``at least one'' becomes ``none''.
\item Count the selections in the complement (usually a single case).
\item Compute $P(E)=1-\dfrac{\text{selections with none}}{\text{total selections}}$.
\item Use this method whenever the direct event splits into three or more cases.
\end{enumerate}
\end{methode}

\begin{exoresolu}[A hand of cards with exactly two aces]
A hand of $5$ cards is dealt at random from a standard pack. Find the probability that the hand contains exactly two aces.
\tcblower
A hand is an unordered selection, so we use combinations throughout. Total hands: $\dbinom{52}{5}=2\,598\,960$.
Favourable hands: choose $2$ aces from $4$, and $3$ non-aces from the $48$ remaining cards:
\[ P=\frac{\dbinom{4}{2}\dbinom{48}{3}}{\dbinom{52}{5}}=\frac{6\times 17\,296}{2\,598\,960}=\frac{103\,776}{2\,598\,960}=\frac{2162}{54\,145}\approx 0.0399. \]
So roughly one hand in $25$ contains exactly two aces.
\end{exoresolu}

\begin{exoresolu}[Arrangements with restrictions]
The letters of the word \textbf{CHANCE} are arranged in a row at random. Find the probability that the arrangement begins and ends with a consonant.
\tcblower
CHANCE has $6$ letters, the letter C appearing twice; the consonants are C, H, N. Total distinct arrangements: $\dfrac{6!}{2!}=360$.
Favourable arrangements: the two end letters must be consonants. Count by cases on the end letters:
\begin{itemize}
\item Ends are H and N (in either order): $2$ choices for the ordered pair of ends; the middle $4$ letters (C, C, A, E) arrange in $\dfrac{4!}{2!}=12$ ways: $2\times 12 = 24$.
\item Ends are C and H, or C and N: for each such pair, $2$ orders for the ends; the middle letters are then $4$ distinct letters, arranging in $4! = 24$ ways. Two such pairs: $2\times 2\times 24 = 96$.
\item Ends are both C: $1$ way for the ends; the middle letters H, A, N, E arrange in $4!=24$ ways: $24$.
\end{itemize}
Total favourable: $24+96+24=144$. Hence
\[ P=\frac{144}{360}=\frac{2}{5}. \]
\end{exoresolu}

\begin{exoresolu}[Digits and complementary counting]
Four different digits are chosen at random from $\{1,2,3,4,5,6,7,8,9\}$ and arranged in the order chosen to form a 4-digit number. Find the probability that the number contains at least one even digit.
\tcblower
Order matters (the arrangement forms the number), so use permutations consistently. Total numbers: ${}^{9}P_{4}=\dfrac{9!}{5!}=3024$.
Complement: no even digit, i.e.\ all four digits chosen from the $5$ odd digits $\{1,3,5,7,9\}$: ${}^{5}P_{4}=120$.
\[ P(\text{at least one even})=1-\frac{120}{3024}=1-\frac{5}{126}=\frac{121}{126}. \]
\end{exoresolu}

\courssub{GCE-Style Mixed Exercises}

\begin{exercice}[Mixed set: counting, addition and multiplication laws]
\begin{enumerate}
\item Two fair dice are thrown. Find the probability that the sum is $8$ \emph{or} a double is obtained.
\item A bag contains $5$ red and $4$ blue balls. Three balls are drawn at random without replacement. Find the probability that they are all of the same colour, (a) using a tree-type multiplication argument; (b) using combinations.
\item A delegation of $5$ students is chosen from $8$ boys and $7$ girls. Find the probability that the delegation contains at least $2$ girls.
\item A card is drawn from a pack, replaced, and a second card drawn. Find the probability that exactly one of the two cards is a king. Compare with the answer when the first card is \emph{not} replaced.
\item How many distinct arrangements of the letters of \textbf{BAMENDA} are there? If one arrangement is chosen at random, find the probability that the two A's are together.
\end{enumerate}
\end{exercice}

\begin{corrige}[Exercise 1]
\textbf{1.} $P(\text{sum }8)=\frac{5}{36}$ (pairs $(2,6),(3,5),(4,4),(5,3),(6,2)$); $P(\text{double})=\frac{6}{36}$; intersection: only $(4,4)$, so $\frac{1}{36}$. Addition law: $P=\dfrac{5+6-1}{36}=\dfrac{10}{36}=\dfrac{5}{18}$.

\textbf{2.} (a) $P(\text{all red})=\dfrac{5}{9}\cdot\dfrac{4}{8}\cdot\dfrac{3}{7}=\dfrac{60}{504}=\dfrac{5}{42}$; $P(\text{all blue})=\dfrac{4}{9}\cdot\dfrac{3}{8}\cdot\dfrac{2}{7}=\dfrac{24}{504}=\dfrac{1}{21}$. Total: $\dfrac{5}{42}+\dfrac{1}{21}=\dfrac{7}{42}=\dfrac16$.
(b) $\dfrac{\dbinom{5}{3}+\dbinom{4}{3}}{\dbinom{9}{3}}=\dfrac{10+4}{84}=\dfrac{14}{84}=\dfrac16$. Both methods agree, as they must.

\textbf{3.} Complement: $0$ or $1$ girl.
\[ P=1-\frac{\dbinom{7}{0}\dbinom{8}{5}+\dbinom{7}{1}\dbinom{8}{4}}{\dbinom{15}{5}}=1-\frac{56+7\times 70}{3003}=1-\frac{546}{3003}=1-\frac{2}{11}=\frac{9}{11}. \]

\textbf{4.} With replacement (independence): $P=2\times\dfrac{4}{52}\times\dfrac{48}{52}=\dfrac{384}{2704}=\dfrac{24}{169}\approx 0.142$. Without replacement: $P=2\times\dfrac{4}{52}\times\dfrac{48}{51}=\dfrac{384}{2652}=\dfrac{32}{221}\approx 0.145$. The probability is slightly larger without replacement, since removing a non-king raises the proportion of kings left in the pack.

\textbf{5.} BAMENDA has $7$ letters with A repeated twice: $\dfrac{7!}{2!}=2520$ distinct arrangements. Treat the two A's as a single block: $6$ objects, giving $6!=720$ arrangements. Hence $P=\dfrac{720}{2520}=\dfrac{2}{7}$.
\end{corrige}

\begin{attention}[Final checklist for counting-based probability]
\begin{itemize}
\item Are the counted objects equally likely? (Ordered pairs of dice: yes. Sums of dice: no.)
\item Same counting type in numerator and denominator.
\item With or without replacement? This changes both the totals and the independence structure.
\item ``At least one'' $\Rightarrow$ think complement first.
\end{itemize}
\end{attention}

With these counting techniques in hand, we are ready for the next section, where conditional probability reaches its most powerful form: Bayes' theorem, which lets us \emph{reverse} a conditional probability --- reasoning from an observed effect back to its most likely cause.

\coursec{Bayes' Theorem}

In the previous section we used permutations and combinations to count favourable outcomes and compute probabilities in games of chance. In all those problems the flow of information was ``forward'': we knew the composition of a bag, a machine or a population, and we asked for the probability of an outcome. Real life, however, very often forces us to reason \emph{backwards}. A doctor knows how reliable a test is and observes a positive result; she wants the probability that the patient is actually sick. A factory manager finds a defective item on the conveyor belt; he wants to know which machine most likely produced it. This reversal of conditional probability is exactly the business of \cle{Bayes' theorem}, one of the most powerful and most frequently examined tools of the chapter.

\courssub{Motivation: reversing a conditional probability}

\begin{experience}[The screening test dilemma]
A disease affects $1\%$ of a population. A screening test detects the disease correctly in $99\%$ of sick people, but it also gives a (false) positive result for $5\%$ of healthy people. A person from Yaound\'e tests positive. What is the probability that the person really has the disease?

Most students instinctively answer ``about $99\%$''. We shall see that the true answer is below $17\%$! The given information is $P(\text{positive}\mid\text{disease})$, but the question asks for $P(\text{disease}\mid\text{positive})$. These two conditional probabilities run in opposite directions, and confusing them is one of the costliest errors in probability.
\end{experience}

\begin{attention}[Confusing $P(A\mid B)$ with $P(B\mid A)$]
Confusing $P(A\mid B)$ with $P(B\mid A)$ is the single most common mistake in this topic. They are generally \textbf{different} numbers:
\[
P(A\mid B)=\frac{P(A\cap B)}{P(B)},\qquad P(B\mid A)=\frac{P(A\cap B)}{P(A)}.
\]
The numerators agree, but the denominators differ. For example, $P(\text{positive test}\mid\text{disease})=0.99$ while $P(\text{disease}\mid\text{positive test})\approx 0.167$ in the situation above. Always identify carefully which event is the \emph{condition} (the given information) and which is the event whose probability is required.
\end{attention}

\courssub{The Law of Total Probability}

Before reversing a conditional probability we must be able to compute the probability of the ``observed'' event itself, when it can happen through several distinct routes.

\begin{definition}[Partition of a sample space]
Events $A_1,A_2,\dots,A_n$ form a \cle{partition} of the sample space $S$ if they are pairwise mutually exclusive and exhaustive, i.e.
\[
A_i\cap A_j=\varnothing\ (i\neq j),\qquad A_1\cup A_2\cup\cdots\cup A_n=S.
\]
In particular $P(A_1)+P(A_2)+\cdots+P(A_n)=1$.
\end{definition}

\begin{propriete}[Law of total probability]
Let $\{A_1,A_2,\dots,A_n\}$ be a partition of $S$ with $P(A_i)>0$ for each $i$, and let $B$ be any event. Then
\[
P(B)=P(B\mid A_1)P(A_1)+P(B\mid A_2)P(A_2)+\cdots+P(B\mid A_n)P(A_n)=\sum_{i=1}^{n}P(B\mid A_i)\,P(A_i).
\]
\emph{Proof (instructive, not required for the exam).} Since the $A_i$ cover $S$, $B=(B\cap A_1)\cup(B\cap A_2)\cup\cdots\cup(B\cap A_n)$, and these pieces are mutually exclusive. Hence $P(B)=\sum_i P(B\cap A_i)=\sum_i P(B\mid A_i)P(A_i)$ by the multiplication law. $\blacksquare$
\end{propriete}

\begin{popfigure}
\begin{tikzpicture}[scale=1.0]
\draw[thick, popdark] (0,0) rectangle (9,5);
\node[popdark] at (0.35,4.6) {$S$};
\draw[thick, popblue] (3.2,0) -- (3.2,5);
\draw[thick, popblue] (6.1,0) -- (6.1,5);
\node[popblue] at (1.6,4.35) {$A_1$};
\node[popblue] at (4.65,4.35) {$A_2$};
\node[popblue] at (7.55,4.35) {$A_3$};
\draw[thick, poporange, fill=poporange, fill opacity=0.15] (4.9,2.1) ellipse (3.4 and 1.35);
\node[poporange] at (7.6,1.15) {$B$};
\fill[poporange] (2.2,2.3) circle (0.05);
\fill[poporange] (4.4,1.8) circle (0.05);
\fill[poporange] (6.9,2.5) circle (0.05);
\node[popmuted, align=center, text width=2.6cm] at (1.6,0.6) {$B\cap A_1$};
\node[popmuted, align=center, text width=2.6cm] at (4.65,0.6) {$B\cap A_2$};
\node[popmuted, align=center, text width=2.6cm] at (7.55,0.6) {$B\cap A_3$};
\end{tikzpicture}
\legende{The event $B$ crosses every part of the partition $\{A_1,A_2,A_3\}$; its probability is the sum of the contributions $P(B\cap A_i)=P(B\mid A_i)P(A_i)$.}
\end{popfigure}

\courssub{Bayes' Theorem}

\begin{propriete}[Bayes' theorem]
Let $\{A_1,A_2,\dots,A_n\}$ be a partition of $S$ with $P(A_i)>0$, and let $B$ be an event with $P(B)>0$. Then for each $i$,
\[
P(A_i\mid B)=\frac{P(B\mid A_i)\,P(A_i)}{\displaystyle\sum_{j=1}^{n}P(B\mid A_j)\,P(A_j)}.
\]
\emph{Derivation (examinable in two lines).} By the definition of conditional probability and the multiplication law,
\[
P(A_i\mid B)=\frac{P(A_i\cap B)}{P(B)}=\frac{P(B\mid A_i)P(A_i)}{P(B)},
\]
and the denominator is expanded by the law of total probability. $\blacksquare$
\end{propriete}

\begin{aretenir}[Reading Bayes' formula]
\begin{itemize}
\item $P(A_i)$ is the \cle{prior} probability (before observing $B$);
\item $P(B\mid A_i)$ is the \cle{likelihood} (how well $A_i$ explains the observation);
\item $P(A_i\mid B)$ is the \cle{posterior} probability (after observing $B$);
\item the denominator is simply $P(B)$, the total probability of the observed event.
\end{itemize}
\end{aretenir}

\courssub{Solving Bayes problems with a tree diagram}

The formula is compact, but at GCE level the \textbf{tree diagram is the recommended method}: it organises the computation and makes the denominator impossible to forget.

\begin{methode}[Bayes via a tree diagram]
\begin{enumerate}
\item Draw a two-stage tree: first the ``cause'' branches $A_1,A_2,\dots$ with the prior probabilities, then the ``result'' branches ($B$ and $B'$) with the conditional probabilities.
\item Compute the probability of \textbf{every path that ends at the observed event} $B$ (multiply along each path).
\item Add these path probabilities: this is $P(B)$, the denominator.
\item The required answer is
\[
P(A_i\mid B)=\frac{\text{probability of the path through }A_i\text{ ending at }B}{\text{sum of all path probabilities ending at }B}.
\]
\end{enumerate}
\end{methode}

\begin{attention}[Incomplete denominator]
A frequent error is to divide by only one or two of the paths leading to the observed event. The denominator must contain \textbf{all} branches that produce $B$. If the test can be positive for both sick and healthy people, both paths appear in $P(\text{positive})$.
\end{attention}

\begin{exoresolu}[Two machines in a factory]
A factory in Douala uses two machines, $M_1$ and $M_2$. Machine $M_1$ produces $60\%$ of the items and $M_2$ produces $40\%$. Of the items from $M_1$, $2\%$ are defective; of those from $M_2$, $5\%$ are defective. An item is chosen at random and found to be defective. Find the probability that it was produced by machine $M_1$.
\tcblower
\textbf{Solution.} Let $D$ = ``the item is defective''. From the tree,
\[
P(M_1\cap D)=0.6\times 0.02=0.012,\qquad P(M_2\cap D)=0.4\times 0.05=0.020.
\]
Total probability of a defective item:
\[
P(D)=0.012+0.020=0.032.
\]
By Bayes' theorem,
\[
P(M_1\mid D)=\frac{P(M_1\cap D)}{P(D)}=\frac{0.012}{0.032}=\frac{3}{8}=0.375.
\]
So although $M_1$ produces the majority of items, a defective item is \emph{more} likely ($0.625$) to have come from the less reliable machine $M_2$.
\end{exoresolu}

\begin{popfigure}
\begin{tikzpicture}[grow=right, level distance=3.2cm,
level 1/.style={sibling distance=3.4cm},
level 2/.style={sibling distance=1.5cm},
every node/.style={font=\small}]
\node {Start}
child { node {$M_2\ (0.4)$}
  child { node[popmuted] {$D'$: $0.4\times0.95=0.380$} edge from parent node[below, pos=0.6] {$0.95$} }
  child { node[poporange] {$D$: $0.4\times0.05=0.020$} edge from parent node[above, poporange, pos=0.6] {$0.05$} }
  edge from parent node[below, pos=0.6] {$0.4$} }
child { node {$M_1\ (0.6)$}
  child { node[popmuted] {$D'$: $0.6\times0.98=0.588$} edge from parent node[below, pos=0.6] {$0.98$} }
  child { node[poporange] {$D$: $0.6\times0.02=0.012$} edge from parent node[above, poporange, pos=0.6] {$0.02$} }
  edge from parent node[above, pos=0.6] {$0.6$} };
\end{tikzpicture}
\legende{Tree for the two-machine problem. The two orange paths lead to the observed event $D$; $P(M_1\mid D)=\dfrac{0.012}{0.012+0.020}$.}
\end{popfigure}

\begin{exoresolu}[Medical screening test]
Return to the opening situation: a disease affects $1\%$ of a population; the test is positive for $99\%$ of sick people and (falsely) for $5\%$ of healthy people. A person tests positive. Find the probability that the person has the disease.
\tcblower
\textbf{Solution.} Let $M$ = ``has the disease'', $T$ = ``tests positive''. The paths ending at $T$ are:
\[
P(M\cap T)=0.01\times 0.99=0.0099,\qquad P(M'\cap T)=0.99\times 0.05=0.0495.
\]
Hence $P(T)=0.0099+0.0495=0.0594$ and
\[
P(M\mid T)=\frac{0.0099}{0.0594}=\frac{1}{6}\approx 0.167.
\]
Despite a ``$99\%$ reliable'' test, a positive result means only about a $17\%$ chance of disease.
\end{exoresolu}

\textbf{Why is the answer so small?} Because the disease is \emph{rare}: among $10\,000$ people, about $100$ are sick (of whom $99$ test positive), while $9\,900$ are healthy (of whom $495$ test positive falsely). The $99$ true positives are drowned among $495$ false positives. The lesson: when the prior probability $P(M)$ is tiny, even a good test yields a modest posterior probability. This is why positive screening results are always confirmed by a second, more specific test.

\begin{exoresolu}[Students from two schools]
School $A$ supplies $45\%$ and school $B$ supplies $55\%$ of the candidates of a centre. The pass rate is $80\%$ in school $A$ and $70\%$ in school $B$. A candidate chosen at random has passed. Find the probability that the candidate came from school $A$.
\tcblower
\textbf{Solution.} Let $P$ = ``the candidate passed''.
\[
P(A\cap P)=0.45\times0.80=0.36,\qquad P(B\cap P)=0.55\times0.70=0.385.
\]
\[
P(P)=0.36+0.385=0.745,\qquad
P(A\mid P)=\frac{0.36}{0.745}\approx 0.483.
\]
The candidate is slightly more likely to come from school $B$, because $B$ supplies more candidates even though its pass rate is lower.
\end{exoresolu}

\courssub{GCE exam technique}

\begin{methode}[Recognising a Bayes question]
A question calls for Bayes' theorem whenever:
\begin{itemize}
\item the data give conditional probabilities in one direction (e.g.\ $P(\text{defective}\mid\text{machine})$), and
\item the question asks for the \textbf{reversed} conditional probability, typically worded: ``\emph{given that} the item is defective, find the probability that it came from machine $A$'', or ``\emph{if} the test is positive, find the probability that the patient is sick''.
\end{itemize}
Then: draw the tree, mark all paths ending at the given event, and divide the relevant path probability by their sum.
\end{methode}

\begin{exercice}[Bayes practice]
\begin{enumerate}
\item Three boxes $B_1,B_2,B_3$ contain respectively $2$ defective out of $10$, $3$ out of $10$ and $1$ out of $10$ items. A box is chosen at random and an item drawn from it is defective. Find the probability that it came from $B_2$.
\item In a town, $30\%$ of commuters use taxis and $70\%$ use buses. $10\%$ of taxi users and $25\%$ of bus users arrive late. A commuter arrives late; find the probability that he used a bus.
\item A test detects a virus in $95\%$ of carriers and gives false positives for $2\%$ of non-carriers. If $0.5\%$ of the population carries the virus, find the probability that a person who tests positive is a carrier.
\end{enumerate}
\end{exercice}

\begin{corrige}[Exercise 1]
\textbf{1.} Paths to ``defective'': $\tfrac13\cdot\tfrac{2}{10}=\tfrac{2}{30}$, $\tfrac13\cdot\tfrac{3}{10}=\tfrac{3}{30}$, $\tfrac13\cdot\tfrac{1}{10}=\tfrac{1}{30}$. Total $=\tfrac{6}{30}=\tfrac15$. Hence $P(B_2\mid D)=\tfrac{3/30}{6/30}=\tfrac12$.

\textbf{2.} $P(\text{late})=0.3(0.1)+0.7(0.25)=0.03+0.175=0.205$. $P(\text{bus}\mid\text{late})=\dfrac{0.175}{0.205}=\dfrac{35}{41}\approx0.854$.

\textbf{3.} $P(T)=0.005(0.95)+0.995(0.02)=0.00475+0.0199=0.02465$. $P(\text{carrier}\mid T)=\dfrac{0.00475}{0.02465}\approx0.193$. Again the rare disease gives a low posterior probability.
\end{corrige}

\begin{aretenir}[Key points of the section]
\begin{itemize}
\item Total probability: $P(B)=\sum_i P(B\mid A_i)P(A_i)$ for a partition $\{A_i\}$.
\item Bayes: $P(A_i\mid B)=\dfrac{P(B\mid A_i)P(A_i)}{\sum_j P(B\mid A_j)P(A_j)}$.
\item With a tree: answer $=\dfrac{\text{path through }A_i\text{ to }B}{\text{all paths to }B}$.
\item $P(A\mid B)\neq P(B\mid A)$ in general; a positive test for a rare disease can still mean a small probability of illness.
\end{itemize}
\end{aretenir}

\begin{savaistu}[Thomas Bayes and modern life]
The theorem is named after the Reverend Thomas Bayes, whose work was published after his death in the 18th century. Today Bayesian reasoning hides inside spam filters, medical diagnosis software, weather forecasting and even the algorithms that recommend videos on your phone: each new observation updates a prior belief into a posterior one, exactly as in this lesson.
\end{savaistu}

\coursec{Applications, Synthesis and Exam Preparation}

Having mastered Bayes' theorem, we now step back to see how all the tools of this chapter fit together. Probability is not just a collection of formulas; it is a language for quantifying uncertainty in the real world. In this final section we explore authentic applications, develop a systematic strategy for choosing the right tool, work through GCE-style multi-part questions, and run a clinic on the most common exam errors.

\courssub{Real-Life Applications of Probability}

Probability theory was born from games of chance, but today it underpins decision-making in engineering, medicine, finance, and daily life in Cameroon. We illustrate four key domains.

\begin{experience}[Games of Chance: Fairness and Expectation]
Consider a simple dice game popular in school clubs: you pay \SI{100}{FCFA} to roll a fair die. If you roll a 6 you win \SI{500}{FCFA}; otherwise you win nothing. Is the game fair?
\begin{itemize}
    \item Sample space: $\Omega = \{1,2,3,4,5,6\}$, all outcomes equally likely.
    \item Event $W$ = ``win'' = $\{6\}$, so $P(W)=\dfrac{1}{6}$.
    \item Net gain $X$: $X=400$ FCFA if you win ($500-100$), $X=-100$ FCFA if you lose.
    \item Average (expected) net gain per play:
    \[
    E(X) = 400\times\frac{1}{6} + (-100)\times\frac{5}{6} = \frac{400-500}{6} = -\frac{100}{6} \approx -16.67\ \text{FCFA}.
    \]
\end{itemize}
The negative average gain means the game is \emph{unfair} to the player; on average you lose about 17 FCFA per play. In a \cle{fair game} the average net gain is $0$. This idea of \cle{expected value} is the quantitative backbone of ``fairness''. (Expectation is developed further in the Statistics option; here we use it only informally, as a weighted average of gains.)
\end{experience}

\begin{experience}[Quality Control: Acceptance Sampling]
A factory in Douala produces batches of 1000 bolts. The quality manager selects 5 bolts at random \emph{without replacement}; the batch is accepted only if all 5 are non-defective. Suppose a batch actually contains 20 defective bolts (a 2\% defect rate). What is the probability the batch is accepted?
\begin{itemize}
    \item We choose 5 bolts from 1000 and require all 5 to come from the 980 good bolts. Since selection is without replacement and order does not matter, we use combinations:
    \[
    P(\text{accept}) = \frac{\dbinom{980}{5}\dbinom{20}{0}}{\dbinom{1000}{5}} = \frac{\dbinom{980}{5}}{\dbinom{1000}{5}}.
    \]
    \item Because the sample (5) is very small compared with the population (1000), the probabilities on the second, third, \dots\ draws barely change, so we may approximate the draws as independent with $p = 0.98$ for a good bolt:
    \[
    P(\text{accept}) \approx (0.98)^5 \approx 0.9039.
    \]
\end{itemize}
So there is about a 90.4\% chance the batch passes even though it has a 2\% defect rate. A quality engineer would repeat this calculation for many possible defect rates to judge whether the sampling plan is strict enough — an important industrial application of probability.
\end{experience}

\begin{experience}[Everyday Cameroonian Contexts]
\begin{enumerate}
    \item \textbf{Transport delays:} A taxi driver knows that on a rainy day, the probability of a traffic jam on the Yaound\'e--Douala highway is 0.7. If there is a jam, he arrives late with probability 0.9; if no jam, he arrives late with probability 0.1. He leaves on a rainy day. What is the probability he arrives on time?\\
    \emph{Solution sketch (total probability):} $P(\text{on time}) = 0.7\times 0.1 + 0.3\times 0.9 = 0.07 + 0.27 = 0.34$.
    \item \textbf{Rainfall forecast:} The meteorological service forecasts rain with probability 0.8 when it actually rains, and forecasts rain with probability 0.1 when it does not rain. It rains on 30\% of days in August. You hear a forecast of rain. What is the probability it actually rains?\\
    \emph{Solution sketch (Bayes):} $P(R|F) = \dfrac{0.8\times 0.3}{0.8\times 0.3 + 0.1\times 0.7} = \dfrac{0.24}{0.31} \approx 0.774$.
    \item \textbf{Mobile network failure:} Your phone shows ``No Service''. The probability of a tower outage in your area is 0.02. If the tower is down, you certainly have no service. If the tower is up, your phone might still show no service due to a handset fault with probability 0.01. Given ``No Service'', what is the probability the tower is down?\\
    \emph{Solution sketch (Bayes):} $P(\text{down}|\text{no service}) = \dfrac{1\times 0.02}{1\times 0.02 + 0.01\times 0.98} = \dfrac{0.02}{0.0298} \approx 0.671$.
    \item \textbf{Exam guessing strategy:} A multiple-choice question has 5 options, only one correct. You can eliminate 2 options for sure. Should you guess among the remaining 3 or leave it blank? Assume $+4$ marks for a correct answer, $-1$ for a wrong answer, $0$ for blank. Average gain if you guess: $\dfrac{1}{3}\times 4 + \dfrac{2}{3}\times(-1) = \dfrac{2}{3} > 0$. So guessing is profitable on average: \textbf{guess!}
\end{enumerate}
\end{experience}

\begin{aretenir}[Communicating Risk: Fractions, Decimals, Percentages, Odds]
The same probability can be expressed in different forms. Be fluent in converting between them.
\begin{itemize}
    \item \textbf{Fraction:} $\dfrac{3}{8}$ (exact, preferred in exams).
    \item \textbf{Decimal:} $0.375$ (exact if terminating, otherwise rounded).
    \item \textbf{Percentage:} $37.5\%$ (multiply the decimal by 100).
    \item \textbf{Odds in favour:} $\dfrac{P}{1-P} = \dfrac{3/8}{5/8} = \dfrac{3}{5}$, read ``3 to 5''.
    \item \textbf{Odds against:} $\dfrac{1-P}{P} = \dfrac{5}{3}$, read ``5 to 3''.
\end{itemize}
\textbf{Exam convention:} Unless the question asks for a percentage or decimal, give your final answer as a \textbf{simplified fraction}. Odds are rarely asked but may appear in word problems.
\end{aretenir}

\begin{popfigure}
\begin{tikzpicture}
\begin{axis}[
    ybar,
    bar width=12pt,
    width=0.9\linewidth,
    height=7cm,
    symbolic x coords={Batch accepted, Batch rejected},
    xtick=data,
    xlabel={Outcome of sampling 5 bolts},
    ylabel={Probability},
    ymin=0, ymax=1,
    nodes near coords,
    nodes near coords align={vertical},
    enlarge x limits=0.3,
    ymajorgrids=true,
    grid style=dashed,
]
\addplot[fill=popblue!70] coordinates {(Batch accepted,0.9039) (Batch rejected,0.0961)};
\end{axis}
\end{tikzpicture}
\legende{Quality control: probability that a batch with 2\% defectives passes a 5-item acceptance sample (independence approximation).}
\end{popfigure}

\courssub{Synthesis: Choosing the Right Tool — A Decision Strategy}

Faced with a probability problem, many students freeze because they do not know whether to use a Venn diagram, a tree diagram, a formula, counting techniques, or Bayes' theorem. The following method and flowchart will guide you.

\begin{methode}[Problem-Solving Strategy for Probability]
\begin{enumerate}
    \item \textbf{Read carefully and identify the random experiment.} What is being observed? (e.g., drawing cards, testing components, weather forecast).
    \item \textbf{Define the sample space $\Omega$ and the events of interest.} Give each event a clear letter and a verbal description (e.g., $A$ = ``first bolt is defective'', $B$ = ``forecast says rain''). \textbf{Marks are awarded for clear notation.}
    \item \textbf{Determine the type of probability:} Classical (equally likely outcomes), Empirical (given frequencies), or given probabilities in a model.
    \item \textbf{Choose the representation/tool:}
        \begin{itemize}
            \item \textbf{Venn diagram} if you have two or three events with given $P(A)$, $P(B)$, $P(A\cap B)$ and need unions, intersections, complements.
            \item \textbf{Tree diagram} if the experiment has sequential stages (with or without replacement) or you need conditional probabilities $P(B|A)$.
            \item \textbf{Counting techniques (permutations/combinations)} if the sample space is large but structured (cards, lottery, committees) and outcomes are equally likely.
            \item \textbf{Bayes' theorem} if you are given ``reverse'' conditional probabilities (e.g., $P(\text{test positive} | \text{disease})$) and asked for $P(\text{disease} | \text{test positive})$.
        \end{itemize}
    \item \textbf{Set up the calculation using the appropriate law(s):} Addition, Multiplication, Total Probability, Bayes.
    \item \textbf{Compute step by step, keeping fractions exact until the final answer.}
    \item \textbf{Check:} Is the answer in $[0,1]$? Does it make sense in context? Is it simplified?
\end{enumerate}
\end{methode}

\begin{popfigure}
\begin{tikzpicture}[
    decision/.style={diamond, draw=popdark, thick, fill=popblue!20, text width=3.4cm, align=center, inner sep=2pt},
    process/.style={rectangle, draw=popdark, thick, fill=popgreen!20, text width=3.4cm, align=center, inner sep=2pt, rounded corners},
    startstop/.style={ellipse, draw=popdark, thick, fill=poporange!30, text width=3cm, align=center, inner sep=2pt},
    arrow/.style={->, thick, popdark}
]
\node[startstop] (start) {Start: read problem};
\node[process, below=1.1cm of start] (define) {Define $\Omega$ and events $A,B,\dots$ in words};
\node[decision, below=1.1cm of define] (seq) {Sequential stages? (with/without replacement)};
\node[process, left=2.6cm of seq] (tree) {Draw tree diagram; use multiplication law};
\node[decision, below=1.1cm of seq] (count) {Equally likely outcomes, large $\Omega$?};
\node[process, left=2.6cm of count] (comb) {Use combinations or permutations};
\node[decision, right=2.6cm of count] (bayes) {Given $P(B|A)$, want $P(A|B)$?};
\node[process, right=2.6cm of bayes] (bayesbox) {Apply Bayes' theorem};
\node[decision, below=1.1cm of count] (venn) {Two/three events, given intersections?};
\node[process, left=2.6cm of venn] (vennbox) {Draw Venn diagram; use addition law};
\node[process, below=1.1cm of venn] (compute) {Compute step by step; keep fractions exact};
\node[startstop, below=1.1cm of compute] (end) {Final answer (simplified fraction)};

\draw[arrow] (start) -- (define);
\draw[arrow] (define) -- (seq);
\draw[arrow] (seq) -- node[left, pos=0.5] {Yes} (tree);
\draw[arrow] (seq) -- node[right, pos=0.5] {No} (count);
\draw[arrow] (count) -- node[left, pos=0.5] {Yes} (comb);
\draw[arrow] (count) -- node[above, pos=0.5] {No} (bayes);
\draw[arrow] (bayes) -- node[above, pos=0.5] {Yes} (bayesbox);
\draw[arrow] (bayes) |- node[right, pos=0.25] {No} (venn);
\draw[arrow] (venn) -- node[left, pos=0.5] {Yes} (vennbox);
\draw[arrow] (venn) -- node[right, pos=0.5] {No} (compute);
\draw[arrow] (tree) |- (compute);
\draw[arrow] (comb) |- (compute);
\draw[arrow] (bayesbox) |- (compute);
\draw[arrow] (vennbox) |- (compute);
\draw[arrow] (compute) -- (end);
\end{tikzpicture}
\legende{Flowchart: which probability tool should I use?}
\end{popfigure}

\courssub{GCE Past-Paper Style Questions}

The following multi-part questions combine several laws from the chapter. They are typical of the structured questions in the GCE Advanced Level papers.

\begin{exoresolu}[Structured Problem: Factory Components]
A factory has three machines $M_1$, $M_2$, $M_3$ producing identical components. $M_1$ produces 50\% of the output, $M_2$ produces 30\%, and $M_3$ produces 20\%. The probabilities that a component is defective given it came from $M_1$, $M_2$, $M_3$ are 0.02, 0.03, 0.04 respectively. A component is chosen at random from the day's output.
\begin{enumerate}
    \item Draw a tree diagram to represent this information.
    \item Find the probability that the component is defective.
    \item Given that the component is defective, find the probability it was produced by $M_2$.
    \item The day's output is 10,000 components, of which the expected number of defectives is 270. Two components are selected at random, one after the other, without replacement. Given that the first is defective, find the probability that the second is also defective.
\end{enumerate}
\tcblower
\textbf{Solution.}
\begin{enumerate}
    \item \textbf{Tree diagram:} First set of branches: machines $M_1$ (0.5), $M_2$ (0.3), $M_3$ (0.2). From each machine, a second pair of branches: Defective ($D$) with probabilities 0.02, 0.03, 0.04 respectively, and Non-defective ($D'$) with 0.98, 0.97, 0.96. The joint probabilities at the ends of the branches are:
    \begin{align*}
        P(M_1 \cap D) &= 0.5 \times 0.02 = 0.01, \\
        P(M_2 \cap D) &= 0.3 \times 0.03 = 0.009, \\
        P(M_3 \cap D) &= 0.2 \times 0.04 = 0.008.
    \end{align*}
    \item \textbf{Total probability of defective:}
    \[
    P(D) = 0.01 + 0.009 + 0.008 = 0.027 = \frac{27}{1000}.
    \]
    \item \textbf{Bayes' theorem for $P(M_2|D)$:}
    \[
    P(M_2|D) = \frac{P(M_2 \cap D)}{P(D)} = \frac{0.009}{0.027} = \frac{1}{3}.
    \]
    \item \textbf{Second component defective given the first is defective (without replacement).}\\
    After one defective component is removed, 9,999 components remain, of which 269 are defective. Hence
    \[
    P(\text{2nd defective} \mid \text{1st defective}) = \frac{269}{9999} \approx 0.0269.
    \]
    Note how close this is to $P(D)=0.027$: when the population is large, removing one item changes the probabilities only slightly. The question tests whether you recognise the without-replacement adjustment.
\end{enumerate}
\end{exoresolu}

\begin{exoresolu}[Structured Problem: Cards and Conditional Probability]
A standard pack of 52 cards is shuffled. Two cards are drawn at random without replacement.
\begin{enumerate}
    \item Find the probability that both cards are aces.
    \item Find the probability that the second card is an ace given that the first card is an ace.
    \item Find the probability that the first card is an ace given that the second card is an ace.
    \item Are the events ``first card is an ace'' and ``second card is an ace'' independent? Justify.
\end{enumerate}
\tcblower
\textbf{Solution.}
Let $A_1$ = ``first card is an ace'', $A_2$ = ``second card is an ace''.
\begin{enumerate}
    \item $P(A_1 \cap A_2) = \dfrac{4}{52} \times \dfrac{3}{51} = \dfrac{12}{2652} = \dfrac{1}{221}$.
    \item $P(A_2|A_1) = \dfrac{3}{51} = \dfrac{1}{17}$.
    \item We use $P(A_1|A_2) = \dfrac{P(A_1 \cap A_2)}{P(A_2)}$. The marginal probability that the second card is an ace equals that of the first card: $P(A_2) = P(A_1) = \dfrac{4}{52} = \dfrac{1}{13}$ (each position in a shuffled pack is equally likely to hold any card). Therefore
    \[
    P(A_1|A_2) = \frac{1/221}{1/13} = \frac{13}{221} = \frac{1}{17}.
    \]
    \item Check independence: $P(A_1 \cap A_2) = \dfrac{1}{221}$, while $P(A_1)\,P(A_2) = \dfrac{1}{13} \times \dfrac{1}{13} = \dfrac{1}{169}$. Since $\dfrac{1}{221} \neq \dfrac{1}{169}$, the events are \textbf{not independent}. (Intuitively: knowing the first card is an ace reduces the number of aces available for the second draw.)
\end{enumerate}
\end{exoresolu}

\begin{exoresolu}[Structured Problem: Venn Diagram and Algebra]
In a survey of 200 Upper Sixth students in a Cameroonian lyc\'ee:
\begin{itemize}
    \item 120 study Mathematics ($M$),
    \item 90 study Physics ($P$),
    \item 70 study Chemistry ($C$),
    \item 40 study Mathematics and Physics,
    \item 30 study Mathematics and Chemistry,
    \item 20 study Physics and Chemistry,
    \item 10 study all three subjects.
\end{itemize}
A student is selected at random. Find the probability that the student:
\begin{enumerate}
    \item studies Mathematics but not Physics,
    \item studies exactly one of the three subjects,
    \item studies Physics given that they study Mathematics,
    \item studies Chemistry given that they study neither Mathematics nor Physics.
\end{enumerate}
\tcblower
\textbf{Solution.}
Draw a three-circle Venn diagram inside the rectangle $\Omega$ and fill the regions \emph{from the centre outwards}:
\begin{itemize}
    \item Centre (all three): $10$.
    \item $M \cap P$ only: $40 - 10 = 30$; \quad $M \cap C$ only: $30 - 10 = 20$; \quad $P \cap C$ only: $20 - 10 = 10$.
    \item $M$ only: $120 - (30+20+10) = 60$; \quad $P$ only: $90 - (30+10+10) = 40$; \quad $C$ only: $70 - (20+10+10) = 30$.
\end{itemize}
Total in at least one subject: $60+40+30+30+20+10+10 = 200$. So every student studies at least one of the three subjects (the region outside the circles is empty).
\begin{enumerate}
    \item $P(M \cap P') = \dfrac{(M \text{ only}) + (M \cap C \text{ only})}{200} = \dfrac{60+20}{200} = \dfrac{80}{200} = \dfrac{2}{5}$.
    \item $P(\text{exactly one subject}) = \dfrac{60+40+30}{200} = \dfrac{130}{200} = \dfrac{13}{20}$.
    \item $P(P|M) = \dfrac{P(M \cap P)}{P(M)} = \dfrac{40}{120} = \dfrac{1}{3}$.
    \item ``Neither Mathematics nor Physics'' is the region $C$ only, with 30 students; all of them study Chemistry. Hence
    \[
    P\bigl(C \mid (M \cup P)'\bigr) = \frac{30}{30} = 1.
    \]
\end{enumerate}
\end{exoresolu}

\courssub{Common Exam Errors Clinic}

We now review the most frequent mistakes made by candidates in GCE Probability questions. Each warning corresponds to a trap identified in earlier sections.

\begin{attention}[Rounding Too Early]
In multi-step calculations (especially Bayes' theorem or tree diagrams with several branches), \textbf{never round intermediate probabilities}. Keep them as exact fractions (e.g., $\frac{27}{1000}$) until the very last step. Rounding $0.027$ to $0.03$ early can change a final answer from $\frac{1}{3}$ to $0.33$, losing accuracy marks.
\end{attention}

\begin{attention}[{Probability Outside $[0,1]$}]
If your final answer is negative or greater than 1, \textbf{you have made an error}. Common causes:
\begin{itemize}
    \item Adding probabilities of non-mutually-exclusive events without subtracting the intersection.
    \item Using $P(A|B) = \dfrac{P(A)}{P(B)}$ instead of $\dfrac{P(A \cap B)}{P(B)}$.
    \item Arithmetic slips in combination calculations.
\end{itemize}
Always do a sanity check: $0 \le P(\text{event}) \le 1$.
\end{attention}

\begin{attention}[Confusing ``And'' with ``Or'', ``Given'' with ``And'']
\begin{itemize}
    \item ``A and B'' $\rightarrow$ intersection $A \cap B$ $\rightarrow$ Multiplication Law.
    \item ``A or B'' (inclusive) $\rightarrow$ union $A \cup B$ $\rightarrow$ Addition Law.
    \item ``A given B'' $\rightarrow$ conditional probability $P(A|B) = \dfrac{P(A \cap B)}{P(B)}$.
    \item ``A and B are independent'' $\rightarrow$ $P(A \cap B) = P(A)\,P(B)$.
\end{itemize}
Read the wording \emph{very} carefully. ``The probability that a student studies Physics given that they study Mathematics'' is $P(P|M)$, not $P(P \cap M)$.
\end{attention}

\begin{attention}[With vs. Without Replacement]
Tree diagrams must show the correct probabilities on the second set of branches.
\begin{itemize}
    \item \textbf{With replacement:} probabilities on the second branches are identical to those on the first branches.
    \item \textbf{Without replacement:} probabilities on the second branches change (the denominator decreases by 1, and the numerator decreases by 1 if the first item drawn was of that type).
\end{itemize}
If the problem says ``two balls are drawn at random'', assume \textbf{without replacement} unless stated otherwise.
\end{attention}

\begin{attention}[Misusing Combinations vs. Permutations]
Use \textbf{combinations} ($\binom{n}{r}$) when the order of selection does \textbf{not} matter (committees, hands of cards, samples).
Use \textbf{permutations} (${}^nP_r$) when order \textbf{does} matter (arrangements, race rankings, codes).
In probability with equally likely outcomes, you can use either \emph{as long as you are consistent} in numerator and denominator. But combinations are safer for card and ball problems.
\end{attention}

\begin{attention}[Forgetting to Define Events]
Examiners allocate marks for \textbf{clear definition of events}. Always start with:
``Let $A$ be the event that \dots'' and ``Let $B$ be the event that \dots''
Do not just write $P(A|B)$ without saying what $A$ and $B$ are.
\end{attention}

\courssub{Exam Tips Checklist}

\begin{aretenir}[Final Exam Checklist for Probability]
\begin{itemize}
    \item \textbf{Notation:} Define every event in words. Use $P(A)$, $P(A \cap B)$, $P(A \cup B)$, $P(A|B)$, $P(A')$ correctly.
    \item \textbf{Fractions:} Give final answers as simplified fractions (e.g., $\frac{13}{20}$, not $0.65$ or $65\%$, unless asked).
    \item \textbf{Tree diagrams:} Label all branches with probabilities. Check that the branches leaving each node sum to 1.
    \item \textbf{Venn diagrams:} Draw the rectangle ($\Omega$), label the circles, fill regions from the centre outwards, and show the number (or probability) in each region.
    \item \textbf{Bayes' theorem:} Write the formula $P(A|B) = \dfrac{P(B|A)\,P(A)}{P(B)}$ explicitly before substituting values.
    \item \textbf{Independence vs mutual exclusivity:} These are different. Independent: $P(A \cap B)=P(A)\,P(B)$. Mutually exclusive: $P(A \cap B)=0$. Two events of non-zero probability cannot be both independent and mutually exclusive.
    \item \textbf{Total probability:} When using $P(B) = \sum_i P(B|A_i)\,P(A_i)$, ensure the $A_i$ are mutually exclusive and exhaustive.
    \item \textbf{Time management:} Probability questions are often high-mark (8--12 marks). Do not rush; a clear, step-by-step solution earns method marks even if the final arithmetic has a slip.
\end{itemize}
\end{aretenir}

\begin{exercice}[Mixed Practice for Revision]
\begin{enumerate}
    \item A box contains 5 red and 3 blue balls. Two balls are drawn at random without replacement. Find the probability that (a) both are red, (b) exactly one is red, (c) the second is red given the first is blue.
    \item In a certain town, 40\% of households have a car, 60\% have a motorcycle, and 20\% have both. A household is chosen at random. Find the probability it has (a) a car or a motorcycle (or both), (b) a car but not a motorcycle, (c) neither, (d) a motorcycle given it has a car.
    \item A diagnostic test for a disease has sensitivity 95\% (probability of a positive result given the disease) and specificity 90\% (probability of a negative result given no disease). The disease prevalence is 2\%. A person tests positive. Find the probability they actually have the disease.
    \item A committee of 4 is to be chosen from 6 men and 4 women. Find the probability the committee contains (a) 2 men and 2 women, (b) at least 3 women.
    \item Events $A$ and $B$ are such that $P(A)=0.4$, $P(B)=0.5$, and $P(A \cup B)=0.7$. Determine whether $A$ and $B$ are (a) mutually exclusive, (b) independent.
\end{enumerate}
\end{exercice}

\begin{corrige}[Exercise 1]
(a) $\dfrac{5}{8} \times \dfrac{4}{7} = \dfrac{20}{56} = \dfrac{5}{14}$. \quad (b) $P(RB) + P(BR) = \dfrac{5}{8}\times\dfrac{3}{7} + \dfrac{3}{8}\times\dfrac{5}{7} = \dfrac{30}{56} = \dfrac{15}{28}$. \quad (c) After a blue is drawn, 5 red remain out of 7 balls: $\dfrac{5}{7}$.
\end{corrige}

\begin{corrige}[Exercise 2]
Let $C$ = ``has a car'', $M$ = ``has a motorcycle''. $P(C)=0.4$, $P(M)=0.6$, $P(C \cap M)=0.2$.
(a) $P(C \cup M) = 0.4+0.6-0.2 = 0.8$. \quad (b) $P(C \cap M') = 0.4-0.2 = 0.2$. \quad (c) $P\bigl((C \cup M)'\bigr) = 1-0.8 = 0.2$. \quad (d) $P(M|C) = \dfrac{0.2}{0.4} = 0.5$.
\end{corrige}

\begin{corrige}[Exercise 3]
Let $D$ = ``has the disease'', $T$ = ``tests positive''. $P(D)=0.02$, $P(T|D)=0.95$, $P(T'|D')=0.90$ so $P(T|D')=0.10$.
\[
P(T) = P(T|D)P(D) + P(T|D')P(D') = 0.95\times 0.02 + 0.10\times 0.98 = 0.019 + 0.098 = 0.117.
\]
\[
P(D|T) = \frac{0.019}{0.117} = \frac{19}{117} \approx 0.162.
\]
The probability is surprisingly low (about 16\%): because the disease is rare, most positive results come from the large healthy population. This is the well-known \emph{base-rate} effect in medical screening.
\end{corrige}

\begin{corrige}[Exercise 4]
Total number of committees: $\dbinom{10}{4} = 210$.
(a) $\dfrac{\dbinom{6}{2}\dbinom{4}{2}}{210} = \dfrac{15 \times 6}{210} = \dfrac{90}{210} = \dfrac{3}{7}$.
(b) ``At least 3 women'' = exactly 3 women or exactly 4 women.
3 women: $\dfrac{\dbinom{6}{1}\dbinom{4}{3}}{210} = \dfrac{6 \times 4}{210} = \dfrac{24}{210}$; \quad 4 women: $\dfrac{\dbinom{6}{0}\dbinom{4}{4}}{210} = \dfrac{1}{210}$.
Total: $\dfrac{25}{210} = \dfrac{5}{42}$.
\end{corrige}

\begin{corrige}[Exercise 5]
$P(A \cap B) = P(A) + P(B) - P(A \cup B) = 0.4 + 0.5 - 0.7 = 0.2$.
(a) Mutually exclusive $\iff P(A \cap B)=0$. Here $0.2 \neq 0$, so $A$ and $B$ are \textbf{not mutually exclusive}.
(b) Independent $\iff P(A \cap B) = P(A)\,P(B) = 0.4 \times 0.5 = 0.2$. Here $0.2 = 0.2$, so $A$ and $B$ \textbf{are independent}.
\end{corrige}

\vspace{1em}
\noindent\textbf{End of Chapter PMM19.} You have now covered the complete syllabus for Elementary Probability. Revise the summary boxes, re-work the \texttt{exoresolu} examples without looking at the solutions, and practise past GCE papers under timed conditions. Good luck!

\newpage

\coursec{Integration activity}

\courssub{Aims of the activity}

This integration activity closes the chapter \emph{Elementary Probability} by mobilising, in one single authentic case study, the principal notions studied in the lessons:

\begin{itemize}
\item construct a \cle{sample space} and identify \cle{events} in a concrete situation;
\item represent a two-stage experiment by a \cle{tree diagram} and apply the \cle{multiplication law} along branches and the \cle{addition law} across branches;
\item compute a \cle{conditional probability} and apply \cle{Bayes' theorem} to ``reverse'' a conditional probability;
\item estimate a probability \cle{empirically} by simulation and compare it with the \cle{theoretical} value;
\item interpret numerical results and communicate a justified recommendation in ordinary language, as expected in the GCE Advanced Level examination and in professional life.
\end{itemize}

The activity is designed for \textbf{group work} (3 to 4 learners per group), over a double period. Each group needs: a flipchart sheet, markers, a small opaque bag, and 20 identical slips of paper.

\courssub{Situation}

\textbf{The Mobile-Money Agent's Dilemma.}\par
Madam Eposi runs a busy mobile-money kiosk at Mile 17, Buea. Every day, customers come to deposit money onto their mobile wallets or to withdraw cash from them. The kiosk depends on a mobile network connection to validate each transaction, and the network is not always reliable: when the connection drops in the middle of a transaction, the transaction \emph{fails} and the customer must wait or come back later. Madam Eposi is wondering whether it is worth paying for a backup internet connection.

Her record book over the last months gives the following reliable statistics:

\begin{itemize}
\item $60\%$ of all transactions are \textbf{deposits} and $40\%$ are \textbf{withdrawals};
\item a deposit fails (because of a network drop) in $5\%$ of cases;
\item a withdrawal fails in $12\%$ of cases (withdrawals take longer to process, so they are more exposed to network drops).
\end{itemize}

A transaction is chosen at random from her record book. We consider the events:
\begin{itemize}
\item $D$: ``the transaction is a deposit'';
\item $W$: ``the transaction is a withdrawal'';
\item $F$: ``the transaction fails''.
\end{itemize}

Madam Eposi estimates that a backup connection would cost her about $25\,000\ \mathrm{FCFA}$ per month, and that failed transactions cost her, in lost commissions and lost customers, about $200\ \mathrm{FCFA}$ per failed transaction on average. She handles roughly $2\,000$ transactions per month. The backup connection would reduce the failure rate of \emph{both} types of transaction to about $1\%$.

Your group is consulted as a team of young analysts. You must model the situation, quantify the risk of failure, determine where failures come from, test your model by simulation, and finally advise Madam Eposi.

\courssub{Tasks}

\begin{enumerate}
\item \textbf{Modelling.} State the sample space of the two-stage experiment ``type of transaction, then outcome of transaction''. List its four elementary outcomes. Give the values of $P(D)$, $P(W)$, $P(F\mid D)$, $P(F\mid W)$, $P(\overline{F}\mid D)$ and $P(\overline{F}\mid W)$.

\item \textbf{Tree diagram.} Draw a complete probability tree for this experiment, with the first level of branches for the type of transaction and the second level for the outcome (fail / succeed). Write the probability on every branch and the probability of every complete path.

\item \textbf{Overall failure rate.} Using the law of total probability (addition across the paths leading to $F$), compute $P(F)$, the probability that a randomly chosen transaction fails. Express the result as a percentage.

\item \textbf{Where do failures come from?} A customer complains: ``Your network has just swallowed my transaction!'' Using Bayes' theorem, compute the probability that this failed transaction was a withdrawal, i.e.\ $P(W\mid F)$. Compute also $P(D\mid F)$ and comment on the comparison with the original proportions $60\%/40\%$.

\item \textbf{Simulation.} To estimate these probabilities empirically, carry out the following simulation. Prepare 20 slips: 12 marked $D$ and 8 marked $W$ (these reproduce the proportions $60\%/40\%$).
\begin{itemize}
\item Stage 1: draw one slip at random to determine the type of transaction, note it, and return it to the bag.
\item Stage 2: if the slip was $D$, simulate the outcome by drawing from a second set of 20 slips of which exactly 1 is marked $F$ (failure probability $5\%$); if the slip was $W$, draw from a set of 25 slips of which exactly 3 are marked $F$ (failure probability $12\%$).
\end{itemize}
Repeat the complete two-stage trial $50$ times. Record the number of failures $n_F$ and, among the failures, the number of withdrawals $n_{W\cap F}$. Compute the experimental frequencies $\widehat{P}(F)=\dfrac{n_F}{50}$ and $\widehat{P}(W\mid F)=\dfrac{n_{W\cap F}}{n_F}$.

\item \textbf{Comparison.} Compare your experimental frequencies with the theoretical values found in questions 3 and 4. Explain why the two need not coincide, and state what you would expect if the number of trials were increased from $50$ to $5\,000$.

\item \textbf{Advisory note.} Write a short advisory note (8 to 12 lines) to Madam Eposi. Using the theoretical failure rate, estimate the expected number of failed transactions per month and their monthly cost. Compare with the cost of the backup connection (which would bring the failure rate down to $1\%$ for all transactions) and conclude with a clear, justified recommendation.
\end{enumerate}

\courssub{Model answer}

\textbf{1. Modelling.}\par
The experiment has two stages: first the type of transaction ($D$ or $W$), then the outcome ($F$ or $\overline{F}$). The sample space is
\[
S=\{(D,F),\ (D,\overline{F}),\ (W,F),\ (W,\overline{F})\}.
\]
From the data:
\[
P(D)=0.6,\quad P(W)=0.4,\quad P(F\mid D)=0.05,\quad P(F\mid W)=0.12,
\]
\[
P(\overline{F}\mid D)=1-0.05=0.95,\qquad P(\overline{F}\mid W)=1-0.12=0.88.
\]

\textbf{2. Tree diagram.}\par

\begin{popfigure}
\begin{tikzpicture}[grow=right, level distance=3.2cm,
level 1/.style={sibling distance=2.6cm},
level 2/.style={sibling distance=1.3cm},
every node/.style={font=\small}]
\node {Start}
child { node {$W$} 
  child { node {$\overline{F}$} edge from parent node[below] {$0.88$} }
  child { node {$F$} edge from parent node[above] {$0.12$} }
  edge from parent node[above] {$0.4$} }
child { node {$D$}
  child { node {$\overline{F}$} edge from parent node[below] {$0.95$} }
  child { node {$F$} edge from parent node[above] {$0.05$} }
  edge from parent node[above] {$0.6$} };
\end{tikzpicture}
\legende{Probability tree for the mobile-money transactions.}
\end{popfigure}

The path probabilities (multiplication law) are:
\begin{align*}
P(D\cap F)&=0.6\times 0.05=0.03, & P(D\cap\overline{F})&=0.6\times 0.95=0.57,\\
P(W\cap F)&=0.4\times 0.12=0.048, & P(W\cap\overline{F})&=0.4\times 0.88=0.352.
\end{align*}
Check: $0.03+0.57+0.048+0.352=1$.

\textbf{3. Overall failure rate.}\par
The event $F$ is reached through two mutually exclusive paths, so by the law of total probability:
\[
P(F)=P(D\cap F)+P(W\cap F)=0.03+0.048=0.078.
\]
A randomly chosen transaction therefore fails with probability $0.078$, i.e.\ $7.8\%$.

\textbf{4. Where do failures come from?}\par
By Bayes' theorem:
\[
P(W\mid F)=\frac{P(W\cap F)}{P(F)}=\frac{0.048}{0.078}=\frac{48}{78}=\frac{8}{13}\approx 0.615.
\]
Similarly
\[
P(D\mid F)=\frac{0.03}{0.078}=\frac{30}{78}=\frac{5}{13}\approx 0.385.
\]
\textbf{Comment:} although only $40\%$ of all transactions are withdrawals, about $61.5\%$ of the \emph{failed} transactions are withdrawals. Withdrawals are over-represented among failures because their conditional failure rate ($12\%$) is much higher than that of deposits ($5\%$). This is exactly the kind of ``reversal'' that Bayes' theorem quantifies.

\textbf{5. Simulation.}\par
The exact outcome depends on the draws; the following is a typical set of results for 50 trials:
\begin{center}
\begin{tabular}{|l|c|c|}
\hline
\rowcolor{popblueL} Quantity & Theoretical value & Typical experimental value \\
\hline
$P(F)$ & $0.078$ & $\widehat{P}(F)=\dfrac{4}{50}=0.08$ \\
\hline
$P(W\mid F)$ & $0.615$ & $\widehat{P}(W\mid F)=\dfrac{3}{4}=0.75$ \\
\hline
\end{tabular}
\end{center}
(Your group's values will differ; any values reasonably close to the theoretical ones are acceptable.)

\textbf{6. Comparison.}\par
The experimental frequencies are random: with only $50$ trials, chance fluctuations are large, so $\widehat{P}(F)$ and $\widehat{P}(W\mid F)$ generally differ from the theoretical probabilities — especially $P(W\mid F)$, which is estimated from very few failures. By the \cle{law of large numbers}, as the number of trials increases the experimental frequencies stabilise towards the theoretical probabilities: with $5\,000$ trials we would expect $\widehat{P}(F)$ to be very close to $0.078$ and $\widehat{P}(W\mid F)$ very close to $0.615$.

\textbf{7. Advisory note (model).}\par

\begin{quote}
\emph{To Madam Eposi, Mobile-Money Kiosk, Mile 17, Buea.}

\emph{Dear Madam, our analysis of your records shows that about $7.8\%$ of your transactions fail because of network drops. Out of roughly $2\,000$ transactions per month, this represents about $0.078\times 2\,000=156$ failed transactions, costing you about $156\times 200=31\,200\ \mathrm{FCFA}$ per month in lost commissions and lost customers. We also found that more than $6$ failed transactions out of $10$ are withdrawals, so withdrawal customers suffer most. A backup connection at $25\,000\ \mathrm{FCFA}$ per month would reduce the failure rate to about $1\%$, i.e.\ about $20$ failed transactions per month, costing about $4\,000\ \mathrm{FCFA}$. Your total monthly loss would then be about $25\,000+4\,000=29\,000\ \mathrm{FCFA}$, against $31\,200\ \mathrm{FCFA}$ today — a saving of about $2\,200\ \mathrm{FCFA}$ per month, not counting the improved reputation of your kiosk. We therefore recommend that you invest in the backup connection.}
\end{quote}

\begin{attention}[Pitfalls to avoid in this activity]
\begin{itemize}
\item Do not confuse $P(F\mid W)=0.12$ with $P(W\mid F)\approx 0.615$: the order in a conditional probability matters.
\item Do not add $0.05$ and $0.12$ to get a ``failure rate'': the two branches must be weighted by $P(D)$ and $P(W)$ before adding.
\item In the simulation, remember to return the first slip to the bag before the next trial, otherwise the proportions $60\%/40\%$ change.
\end{itemize}
\end{attention}

\begin{savaistu}[Did you know?]
Mobile-money services handle a large share of everyday payments in Cameroonian towns, and network reliability is a real business issue for kiosk agents. The same Bayes' theorem used here is also at the heart of medical test interpretation, spam filters and credit-risk scoring: a small calculation on a flipchart in Buea is the same mathematics used by banks and hospitals worldwide.
\end{savaistu}

\newpage

\coursec{Methods and worked examples}

\courssub{Method 1: Computing a Classical Probability from a Sample Space}

Before applying any formula, recall the framework: a \cle{random experiment} is a process whose outcome cannot be predicted with certainty (throwing a die, drawing a card); the \cle{sample space} $S$ is the set of all possible outcomes; an \cle{event} $E$ is any subset of $S$. When all outcomes are equally likely, the classical (theoretical) probability of $E$ is
\[P(E)=\frac{n(E)}{n(S)}=\frac{\text{number of outcomes favourable to }E}{\text{total number of possible outcomes}}.\]

\begin{methode}[Finding $P(E)$ in an Equally Likely Setting]
\begin{enumerate}
\item \textbf{Identify the experiment} and check that all outcomes are \cle{equally likely} (fair coin, fair die, well-shuffled pack, random selection).
\item \textbf{Describe the sample space} $S$ and count $n(S)$. Use a list, a grid or a table if two objects are involved.
\item \textbf{Define the event} $E$ in words, then list or count its outcomes to get $n(E)$.
\item \textbf{Apply the formula} $P(E)=\dfrac{n(E)}{n(S)}$ and simplify the fraction.
\item \textbf{Check}: $0\le P(E)\le 1$. Use the complement $P(E')=1-P(E)$ when counting $E$ directly is long (especially for ``at least one'').
\end{enumerate}
\textbf{Reflex:} for ``at least one'', always try $1-P(\text{none})$ first.
\end{methode}

\begin{exoresolu}[Two fair dice]
Two fair six-sided dice are thrown together. Find the probability that (a) the scores add up to $8$; (b) the scores add up to at least $10$; (c) at least one die shows a $6$.
\tcblower
The sample space consists of $6\times 6=36$ equally likely ordered pairs $(i,j)$, where $i$ is the score on the first die and $j$ the score on the second.

\textbf{(a)} Sum equal to $8$: the favourable pairs are
$(2,6),(3,5),(4,4),(5,3),(6,2)$, so $n(E)=5$ and
\[P(\text{sum }8)=\frac{5}{36}.\]

\textbf{(b)} Sum at least $10$ means sum $10, 11$ or $12$:
\begin{itemize}
\item sum $10$: $(4,6),(5,5),(6,4)$ — 3 outcomes;
\item sum $11$: $(5,6),(6,5)$ — 2 outcomes;
\item sum $12$: $(6,6)$ — 1 outcome.
\end{itemize}
Hence $P(\text{sum}\ge 10)=\dfrac{3+2+1}{36}=\dfrac{6}{36}=\dfrac{1}{6}$.

\textbf{(c)} ``At least one $6$'': use the complement. The number of outcomes with \emph{no} $6$ is $5\times 5=25$ (each die shows $1$–$5$). Therefore
\[P(\text{at least one }6)=1-\frac{25}{36}=\frac{11}{36}.\]
\textbf{Check:} all answers lie between $0$ and $1$. Note the trap: $(4,4)$ counts \emph{once}, while $(2,6)$ and $(6,2)$ are two different outcomes.
\end{exoresolu}

\begin{popfigure}
\begin{center}
\begin{tikzpicture}[scale=0.62, every node/.style={font=\small}]
\draw[popline] (0,0) grid (6,6);
\foreach \i in {1,...,6} {
  \node at (\i-0.5,-0.4) {\i};
  \node at (-0.4,\i-0.5) {\i};
}
\foreach \p in {(1.5,5.5),(2.5,4.5),(3.5,3.5),(4.5,2.5),(5.5,1.5)} {
  \fill[poporange] \p circle (0.18);
}
\node[popdark] at (3,-1.1) {First die};
\node[popdark, rotate=90] at (-1.1,3) {Second die};
\end{tikzpicture}
\legende{The $36$ equally likely outcomes of two dice; the five dots mark the outcomes with sum $8$.}
\end{center}
\end{popfigure}

\courssub{Method 2: Using the Addition Law and Venn Diagrams}

The \cle{addition law} governs ``or'' questions. For any two events,
\[P(A\cup B)=P(A)+P(B)-P(A\cap B),\]
because adding $P(A)$ and $P(B)$ counts the overlap $A\cap B$ twice, so it must be subtracted once. When $A$ and $B$ are \cle{mutually exclusive} ($A\cap B=\varnothing$), the law reduces to $P(A\cup B)=P(A)+P(B)$.

\begin{methode}[Combining Events with the Addition Law]
\begin{enumerate}
\item \textbf{Translate} the question into set language: $A\cup B$ = ``$A$ or $B$'', $A\cap B$ = ``both'', $A'$ = ``not $A$''.
\item \textbf{Test for mutual exclusivity}: if $A\cap B=\varnothing$ then $P(A\cup B)=P(A)+P(B)$.
\item Otherwise apply the \cle{general addition law}:
\[P(A\cup B)=P(A)+P(B)-P(A\cap B).\]
\item \textbf{Draw a Venn diagram} when data are given as counts: fill the \emph{intersection first}, then the remaining parts, then the outside region.
\item Answer the question asked: ``only $A$'' is $A\cap B'$, ``exactly one'' is $(A\cap B')\cup(A'\cap B)$, ``neither'' is $(A\cup B)'$.
\end{enumerate}
\textbf{Reflex:} never add $P(A)+P(B)$ blindly — check whether the events overlap.
\end{methode}

\begin{exoresolu}[Students in a Form Six class]
In a class of $40$ Upper Sixth students, $22$ study Physics, $18$ study Chemistry and $7$ study both. A student is chosen at random. Find the probability that the student studies (a) Physics or Chemistry; (b) exactly one of the two subjects; (c) neither subject.
\tcblower
Let $P$ = ``studies Physics'', $C$ = ``studies Chemistry''. Then
\[P(P)=\frac{22}{40},\qquad P(C)=\frac{18}{40},\qquad P(P\cap C)=\frac{7}{40}.\]
\textbf{(a)} By the addition law,
\[P(P\cup C)=\frac{22}{40}+\frac{18}{40}-\frac{7}{40}=\frac{33}{40}.\]
\textbf{(b)} ``Exactly one'' excludes the intersection:
\[P(\text{exactly one})=\frac{22-7}{40}+\frac{18-7}{40}=\frac{15+11}{40}=\frac{26}{40}=\frac{13}{20}.\]
\textbf{(c)} ``Neither'' is the complement of the union:
\[P(\text{neither})=1-P(P\cup C)=1-\frac{33}{40}=\frac{7}{40}.\]
\textbf{Venn check:} Physics only $=15$, both $=7$, Chemistry only $=11$, neither $=7$; total $=15+7+11+7=40$. \checkmark
\end{exoresolu}

\begin{popfigure}
\begin{center}
\begin{tikzpicture}[scale=0.9, every node/.style={font=\small}]
\draw[popline] (0,0) rectangle (6.4,4);
\node[popdark] at (6.1,3.7) {$S$};
\draw[popblue, thick] (2.1,2) circle (1.5);
\draw[poporange, thick] (4.1,2) circle (1.5);
\node[popblue] at (1.1,3.2) {$P$};
\node[poporange] at (5.1,3.2) {$C$};
\node at (1.7,2) {$15$};
\node at (3.1,2) {$7$};
\node at (4.5,2) {$11$};
\node at (5.9,0.5) {$7$};
\end{tikzpicture}
\legende{Venn diagram for the class survey: fill the intersection ($7$) first, then the ``only'' regions, then the outside.}
\end{center}
\end{popfigure}

\courssub{Method 3: Conditional Probability and Independence}

The \cle{conditional probability} of $A$ given that $B$ has occurred is written $P(A\mid B)$ and read ``probability of $A$ given $B$''. Knowing that $B$ has occurred shrinks the sample space to $B$, which justifies the definition
\[P(A\mid B)=\frac{P(A\cap B)}{P(B)},\qquad P(B)\ne 0.\]
Rearranging gives the \cle{multiplication law}: $P(A\cap B)=P(B)\,P(A\mid B)$. Two events are \cle{independent} when the occurrence of one does not affect the probability of the other, i.e.\ $P(A\cap B)=P(A)P(B)$, equivalently $P(A\mid B)=P(A)$.

\begin{methode}[Using $P(A\mid B)$ and Testing Independence]
\begin{enumerate}
\item \textbf{Identify the condition}: the event after ``given that'' is the new sample space $B$.
\item Apply the definition:
\[P(A\mid B)=\frac{P(A\cap B)}{P(B)},\qquad P(B)\ne 0.\]
\item Alternatively, \textbf{reduce the sample space}: restrict to outcomes in $B$ and count within it.
\item To \textbf{test independence}, check whether $P(A\cap B)=P(A)\times P(B)$ (equivalently $P(A\mid B)=P(A)$). If equality holds, the events are \cle{independent}; otherwise they are dependent.
\item For independent events, multiply: $P(A\cap B\cap C)=P(A)P(B)P(C)$.
\end{enumerate}
\textbf{Reflex:} ``without replacement'' $\Rightarrow$ dependent; ``with replacement'' or separate experiments $\Rightarrow$ independent.
\end{methode}

\begin{exoresolu}[Drawing balls without replacement]
A bag contains $5$ red and $3$ blue balls. Two balls are drawn at random, one after the other, \textbf{without replacement}. (a) Find the probability that the second ball is blue given that the first is red. (b) Find the probability that both balls are red. (c) Are the events ``first red'' and ``second blue'' independent?
\tcblower
\textbf{(a)} Given the first ball is red, the bag now contains $4$ red and $3$ blue balls, $7$ in total. Hence
\[P(\text{2nd blue}\mid \text{1st red})=\frac{3}{7}.\]
\textbf{(b)} By the multiplication law for dependent events,
\[P(\text{both red})=P(\text{1st red})\times P(\text{2nd red}\mid \text{1st red})=\frac{5}{8}\times\frac{4}{7}=\frac{20}{56}=\frac{5}{14}.\]
\textbf{(c)} Compute each quantity:
\[P(\text{1st red})=\frac{5}{8},\qquad P(\text{2nd blue}\mid \text{1st red})=\frac{3}{7}.\]
Also, $P(\text{2nd blue})$ unconditionally: split according to the first draw (mutually exclusive cases),
\[P(\text{2nd blue})=\frac{5}{8}\cdot\frac{3}{7}+\frac{3}{8}\cdot\frac{2}{7}=\frac{15+6}{56}=\frac{21}{56}=\frac{3}{8}.\]
Since $P(\text{2nd blue}\mid \text{1st red})=\dfrac{3}{7}\ne\dfrac{3}{8}=P(\text{2nd blue})$, the events are \textbf{not independent} — as expected, since the first draw changes the contents of the bag.
\end{exoresolu}

\courssub{Method 4: Tree Diagrams for Multi-Stage Experiments}

A \cle{tree diagram} organises a multi-stage experiment: each branching point is one trial, each branch carries a conditional probability, and a complete path from root to leaf is one outcome of the whole experiment. The two working rules are simply the multiplication law (along a path) and the addition law (across paths).

\begin{methode}[Solving Problems with a Tree Diagram]
\begin{enumerate}
\item \textbf{Draw the stages} in time order; each branching point is one trial.
\item \textbf{Label every branch} with its probability (conditional on the path leading to it). Probabilities at each branching point must sum to $1$.
\item \textbf{Multiply along branches} to get the probability of each complete path.
\item \textbf{Add the paths} that satisfy the event asked.
\item For a ``given that'' question, use the tree to compute $P(A\cap B)$ and $P(B)$, then divide.
\end{enumerate}
\textbf{Reflex:} use the tree to \emph{check} your work — all final path probabilities must add to $1$.
\end{methode}

\begin{exoresolu}[Rain and the school bus]
The probability that it rains on a given morning in Bamenda is $0.4$. If it rains, the probability that the school bus arrives late is $0.7$; if it does not rain, the probability that the bus is late is $0.1$. (a) Draw a tree diagram. (b) Find the probability that the bus is late. (c) Given that the bus is late, find the probability that it rained.
\tcblower
\textbf{(a)} Let $R$ = rain, $L$ = bus late.
\begin{center}
\begin{tikzpicture}[grow=right, level distance=2.6cm,
level 1/.style={sibling distance=2.2cm}, level 2/.style={sibling distance=1.2cm},
every node/.style={font=\small}]
\node {}
child { node {$R'$}
  child { node {$L'$: $0.6\times0.9=0.54$} edge from parent node[below] {$0.9$} }
  child { node {$L$: $0.6\times0.1=0.06$} edge from parent node[above] {$0.1$} }
  edge from parent node[below] {$0.6$} }
child { node {$R$}
  child { node {$L'$: $0.4\times0.3=0.12$} edge from parent node[below] {$0.3$} }
  child { node {$L$: $0.4\times0.7=0.28$} edge from parent node[above] {$0.7$} }
  edge from parent node[above] {$0.4$} };
\end{tikzpicture}
\legende{Tree diagram for rain and bus lateness}
\end{center}
\textbf{(b)} ``Late'' occurs along two paths:
\[P(L)=0.4\times0.7+0.6\times0.1=0.28+0.06=0.34.\]
\textbf{(c)} Conditional probability:
\[P(R\mid L)=\frac{P(R\cap L)}{P(L)}=\frac{0.28}{0.34}=\frac{14}{17}\approx 0.824.\]
\textbf{Check:} path probabilities sum to $0.28+0.12+0.06+0.54=1$. \checkmark
\end{exoresolu}

\courssub{Method 5: Probability with Cards, Dice and Counting Techniques}

When the sample space is large, listing is impossible and we count with \cle{permutations} and \cle{combinations} instead. The golden rule is consistency: the numerator $n(E)$ and the denominator $n(S)$ must be counted under the \emph{same} convention.

\begin{methode}[Using Permutations and Combinations in Probability]
\begin{enumerate}
\item Decide whether \textbf{order matters}: ordered $\Rightarrow$ permutations ${}^{n}P_r=\dfrac{n!}{(n-r)!}$; unordered $\Rightarrow$ combinations $\dbinom{n}{r}=\dfrac{n!}{r!(n-r)!}$.
\item Count $n(S)$ and $n(E)$ using the \emph{same} convention (both ordered or both unordered).
\item For a committee/hand of cards: multiply combinations for each category, e.g.\ choosing $2$ aces from $4$ and $3$ kings from $4$ gives $\binom{4}{2}\binom{4}{3}$.
\item Compute $P(E)=\dfrac{n(E)}{n(S)}$ and simplify by cancelling factorials before evaluating.
\end{enumerate}
\textbf{Key facts:} a standard pack has $52$ cards: $4$ suits, $13$ ranks, $12$ picture cards, $26$ red, $26$ black.
\end{methode}

\begin{exoresolu}[A hand of five cards]
Five cards are drawn at random from a standard pack of $52$ cards (without replacement). Find the probability that the hand contains (a) exactly $3$ hearts; (b) exactly $2$ aces; (c) all five cards of the same colour.
\tcblower
The sample space: $n(S)=\dbinom{52}{5}=2\,598\,960$ unordered hands.

\textbf{(a)} Choose $3$ hearts from $13$ and $2$ non-hearts from $39$:
\[P=\frac{\dbinom{13}{3}\dbinom{39}{2}}{\dbinom{52}{5}}=\frac{286\times 741}{2\,598\,960}=\frac{211\,926}{2\,598\,960}\approx 0.0815.\]

\textbf{(b)} Choose $2$ aces from $4$ and $3$ non-aces from $48$:
\[P=\frac{\dbinom{4}{2}\dbinom{48}{3}}{\dbinom{52}{5}}=\frac{6\times 17\,296}{2\,598\,960}=\frac{103\,776}{2\,598\,960}\approx 0.0399.\]

\textbf{(c)} All red or all black — two mutually exclusive cases, each with $\binom{26}{5}$ hands:
\[P=\frac{2\dbinom{26}{5}}{\dbinom{52}{5}}=\frac{2\times 65\,780}{2\,598\,960}=\frac{131\,560}{2\,598\,960}\approx 0.0506.\]
\textbf{Reflex used:} the same (unordered) convention in numerator and denominator; multiplication principle for mixed categories.
\end{exoresolu}

\courssub{Method 6: Bayes' Theorem (Reversing a Conditional Probability)}

Suppose $B_1,B_2,\dots,B_n$ form a \cle{partition} of the sample space (mutually exclusive and exhaustive). The \cle{law of total probability} expresses any event $A$ through the partition:
\[P(A)=\sum_{i=1}^{n} P(B_i)\,P(A\mid B_i).\]
\cle{Bayes' theorem} then reverses the conditioning — from ``probability of the effect given the cause'' to ``probability of the cause given the effect'':
\[P(B_k\mid A)=\frac{P(B_k)\,P(A\mid B_k)}{\displaystyle\sum_{i=1}^{n} P(B_i)\,P(A\mid B_i)}.\]

\begin{methode}[Applying Bayes' Theorem]
\begin{enumerate}
\item Identify the \textbf{prior partition} $B_1,B_2,\dots,B_n$ (mutually exclusive and exhaustive) and the \textbf{observed event} $A$.
\item Write down the priors $P(B_i)$ and the likelihoods $P(A\mid B_i)$.
\item Compute the total probability:
\[P(A)=\sum_i P(B_i)\,P(A\mid B_i).\]
\item Reverse the conditioning:
\[P(B_k\mid A)=\frac{P(B_k)\,P(A\mid B_k)}{P(A)}.\]
\item Interpret the answer: it is an \emph{updated} probability after observing $A$.
\end{enumerate}
\textbf{Reflex:} Bayes = ``given the effect, find the cause''. A tree diagram does all the bookkeeping.
\end{methode}

\begin{exoresolu}[A medical screening test]
A disease affects $2\%$ of a population. A test detects the disease with probability $0.95$ when it is present, and gives a false positive with probability $0.04$ when the disease is absent. A person tests positive. Find the probability that the person actually has the disease.
\tcblower
Let $D$ = ``has the disease'', $T$ = ``tests positive''. The partition is $\{D,D'\}$ with
\[P(D)=0.02,\quad P(D')=0.98,\quad P(T\mid D)=0.95,\quad P(T\mid D')=0.04.\]
\textbf{Step 1 — total probability of a positive test:}
\[P(T)=P(D)P(T\mid D)+P(D')P(T\mid D')=(0.02)(0.95)+(0.98)(0.04)=0.019+0.0392=0.0582.\]
\textbf{Step 2 — Bayes' theorem:}
\[P(D\mid T)=\frac{P(D)P(T\mid D)}{P(T)}=\frac{0.019}{0.0582}\approx 0.326.\]
\textbf{Interpretation:} even with a positive test, the probability of having the disease is only about $32.6\%$, because the disease is rare and false positives from the large healthy group are numerous. This counter-intuitive result is exactly why Bayes' theorem matters in real screening programmes.
\end{exoresolu}

\courssub{Method 7: Mutually Exclusive, Exhaustive and Complementary Events in Exam Problems}

Three classifications recur constantly in GCE questions. Events $A$ and $B$ are \cle{mutually exclusive} if they cannot occur together ($A\cap B=\varnothing$); they are \cle{exhaustive} if together they cover the whole sample space ($P(A\cup B)=1$); and $A'$ is the \cle{complement} of $A$, with $P(A')=1-P(A)$. Complementary events are both mutually exclusive and exhaustive.

\begin{methode}[Choosing the Right Law Quickly]
\begin{enumerate}
\item \textbf{Classify the events} first: mutually exclusive ($P(A\cap B)=0$)? complementary ($B=A'$)? exhaustive ($P(A\cup B)=1$)? independent?
\item Mutually exclusive $\Rightarrow$ add: $P(A\cup B)=P(A)+P(B)$.
\item Complementary $\Rightarrow$ subtract from $1$: $P(A')=1-P(A)$.
\item Independent $\Rightarrow$ multiply: $P(A\cap B)=P(A)P(B)$.
\item Mixed problems: decompose the event into mutually exclusive \emph{cases}, compute each by multiplication, then add.
\end{enumerate}
\textbf{Warning:} mutually exclusive and independent are \emph{different} ideas — two non-trivial mutually exclusive events can never be independent, since $P(A\cap B)=0\ne P(A)P(B)$.
\end{methode}

\begin{exoresolu}[GCE-style synthesis: hitting a target]
Three marksmen, $X$, $Y$ and $Z$, each fire one shot at a target, independently. Their probabilities of hitting are $\tfrac{1}{2}$, $\tfrac{2}{3}$ and $\tfrac{3}{4}$ respectively. Find the probability that (a) all three hit; (b) none hits; (c) exactly one hits; (d) at least one hits.
\tcblower
\textbf{(a)} Independence:
\[P(\text{all hit})=\frac{1}{2}\times\frac{2}{3}\times\frac{3}{4}=\frac{6}{24}=\frac{1}{4}.\]
\textbf{(b)} Probabilities of missing: $\tfrac12,\tfrac13,\tfrac14$:
\[P(\text{none hits})=\frac{1}{2}\times\frac{1}{3}\times\frac{1}{4}=\frac{1}{24}.\]
\textbf{(c)} ``Exactly one hits'' splits into three mutually exclusive cases:
\begin{align*}
P(X\text{ only})&=\tfrac12\times\tfrac13\times\tfrac14=\tfrac{1}{24},\\
P(Y\text{ only})&=\tfrac12\times\tfrac23\times\tfrac14=\tfrac{2}{24},\\
P(Z\text{ only})&=\tfrac12\times\tfrac13\times\tfrac34=\tfrac{3}{24}.
\end{align*}
Adding (mutually exclusive): $P(\text{exactly one})=\dfrac{1+2+3}{24}=\dfrac{6}{24}=\dfrac14$.
\textbf{(d)} Complement of (b):
\[P(\text{at least one})=1-\frac{1}{24}=\frac{23}{24}.\]
\textbf{Strategy note:} part (d) by the complement took one line; a direct computation would need four cases. Always look for the complement first.
\end{exoresolu}

\begin{attention}[Frequent traps in probability questions]
\begin{itemize}
\item Adding $P(A)+P(B)$ when $A$ and $B$ overlap — subtract $P(A\cap B)$.
\item Treating ``without replacement'' draws as independent.
\item Forgetting that $(3,5)$ and $(5,3)$ are distinct outcomes for two dice, while $(4,4)$ is a single outcome.
\item Confusing $P(A\mid B)$ with $P(B\mid A)$ — they are equal only by accident.
\item Giving a probability greater than $1$ or negative: always sanity-check the final value.
\item Mixing conventions: numerator counted unordered but denominator ordered (or vice versa).
\end{itemize}
\end{attention}

\newpage

\coursec{Exercises}

\courssub{I check my knowledge}

\begin{exercice}[Multiple choice questions]
For each question, choose the single correct answer.
\begin{enumerate}
\item A fair die is thrown once. The probability of obtaining a prime number is:
\begin{enumerate}
\item $\dfrac{1}{6}$ \item $\dfrac{1}{3}$ \item $\dfrac{1}{2}$ \item $\dfrac{2}{3}$
\end{enumerate}
\item If $A$ and $B$ are mutually exclusive events with $P(A)=0.3$ and $P(B)=0.5$, then $P(A\cup B)$ equals:
\begin{enumerate}
\item $0.15$ \item $0.8$ \item $0.65$ \item $0.2$
\end{enumerate}
\item Two events $A$ and $B$ are independent with $P(A)=0.4$ and $P(B)=0.25$. Then $P(A\cap B)$ equals:
\begin{enumerate}
\item $0.65$ \item $0.1$ \item $0.15$ \item $0.55$
\end{enumerate}
\item If $P(A)=0.5$, $P(B)=0.4$ and $P(A\cap B)=0.2$, then $P(A\mid B)$ equals:
\begin{enumerate}
\item $0.4$ \item $0.5$ \item $0.2$ \item $0.8$
\end{enumerate}
\end{enumerate}
\end{exercice}

\begin{exercice}[True or false]
State whether each statement is true or false, justifying your answer with a short proof or a counter-example.
\begin{enumerate}
\item If $A$ and $B$ are mutually exclusive, then they are independent.
\item For any two events $A$ and $B$, $P(A\cup B)=P(A)+P(B)-P(A\cap B)$.
\item If $P(A\mid B)=P(A)$, then $A$ and $B$ are independent.
\item For any event $A$, $P(A)+P(A')=1$.
\end{enumerate}
\end{exercice}

\begin{exercice}[Definitions]
Define each of the following terms, giving one example in each case:
\begin{enumerate}
\item a random experiment and its sample space;
\item mutually exclusive events;
\item exhaustive events;
\item independent events;
\item the conditional probability $P(A\mid B)$.
\end{enumerate}
\end{exercice}

\begin{exercice}[Direct calculations]
\begin{enumerate}
\item A letter is chosen at random from the word \textbf{PROBABILITY}. Find the probability that it is (i) a vowel, (ii) the letter B.
\item A coin is tossed twice. Write down the sample space and find the probability of obtaining exactly one head.
\item Given $P(A)=0.6$, $P(B)=0.5$ and $P(A\cup B)=0.8$, find $P(A\cap B)$ and $P(A\mid B)$.
\end{enumerate}
\end{exercice}

\courssub{I practise}

\begin{exercice}[Two dice]
Two fair dice are thrown together. Find the probability that
\begin{enumerate}
\item the sum of the scores is $7$;
\item the sum is at least $10$;
\item a double is obtained;
\item the sum is $7$ given that at least one die shows a $4$.
\end{enumerate}
\end{exercice}

\begin{exercice}[Balls drawn without replacement]
A bag contains $4$ red and $6$ black balls. Two balls are drawn at random without replacement. Using a tree diagram, find the probability that
\begin{enumerate}
\item both balls are red;
\item the balls are of different colours;
\item the second ball is red given that the first was black.
\end{enumerate}
\end{exercice}

\begin{exercice}[Choosing a committee]
A committee of $4$ people is chosen at random from $5$ men and $6$ women. Using combinations, find the probability that the committee contains
\begin{enumerate}
\item exactly two women;
\item at least one woman;
\item no man.
\end{enumerate}
\end{exercice}

\begin{exercice}[Cards with and without replacement]
A card is drawn from a standard pack of $52$ cards, replaced, and then a second card is drawn. Find the probability of obtaining
\begin{enumerate}
\item two aces;
\item an ace followed by a king;
\item at least one heart.
\end{enumerate}
Repeat parts (a) and (c) assuming the first card is \emph{not} replaced, and comment briefly on the difference.
\end{exercice}

\courssub{I prepare for the GCE}

\begin{exercice}[Subjects chosen by candidates]
In a school, $55\%$ of candidates take Chemistry, $40\%$ take Biology and $25\%$ take both subjects. A candidate is chosen at random.
\begin{enumerate}
\item Illustrate the information on a Venn diagram. \hfill (2 marks)
\item Find the probability that the candidate takes
\begin{enumerate}
\item at least one of the two subjects; \hfill (2 marks)
\item exactly one of the two subjects; \hfill (2 marks)
\item neither subject; \hfill (1 mark)
\item Biology, given that she takes Chemistry. \hfill (2 marks)
\end{enumerate}
\item Determine, with justification, whether taking Chemistry and taking Biology are independent events. \hfill (3 marks)
\end{enumerate}
\end{exercice}

\begin{exercice}[Defective items and Bayes' theorem]
Two machines $A$ and $B$ produce $70\%$ and $30\%$ respectively of a factory's output. It is known that $2\%$ of the items produced by $A$ and $5\%$ of the items produced by $B$ are defective. An item is chosen at random from the total output.
\begin{enumerate}
\item Draw a fully labelled tree diagram to represent this information. \hfill (3 marks)
\item Find the probability that the chosen item is defective. \hfill (3 marks)
\item Given that the item is defective, use Bayes' theorem to find the probability that it came from machine $B$. \hfill (4 marks)
\item Three items are chosen at random from the total output. Find the probability that exactly one of them is defective. \hfill (3 marks)
\end{enumerate}
\end{exercice}

\begin{exercice}[Biased coin and GCE mock examination]
\textbf{Part A.} A biased coin shows heads with probability $\dfrac{2}{3}$. It is tossed three times.
\begin{enumerate}
\item Draw a tree diagram showing all possible outcomes. \hfill (3 marks)
\item Find the probability of obtaining
\begin{enumerate}
\item exactly two heads; \hfill (2 marks)
\item at least one head. \hfill (2 marks)
\end{enumerate}
\item Given that exactly two heads occurred, find the probability that the first toss was a head. \hfill (3 marks)
\end{enumerate}
\textbf{Part B.} In a GCE mock examination, the probability that a candidate passes Mathematics is $0.6$, the probability that he passes Physics is $0.5$, and the probability that he passes both is $0.3$.
\begin{enumerate}
\item Find the probability that a candidate passes
\begin{enumerate}
\item at least one of the two subjects; \hfill (2 marks)
\item exactly one of the two subjects; \hfill (2 marks)
\item Physics, given that he passed Mathematics. \hfill (2 marks)
\end{enumerate}
\item Are the events ``passes Mathematics'' and ``passes Physics'' independent? Justify your answer. \hfill (3 marks)
\end{enumerate}
\end{exercice}

\newpage

\coursec{Solutions to the exercises}

\courssub{I check my knowledge}

\begin{corrige}[Exercise 1 --- Multiple choice questions]
\begin{enumerate}
\item The prime numbers on a standard die are $\{2,3,5\}$, i.e. $3$ favourable outcomes out of $6$ equally likely outcomes: $P=\dfrac{3}{6}=\dfrac{1}{2}$. \textbf{Answer: (c).}
\item For mutually exclusive events $A$ and $B$, $P(A\cup B)=P(A)+P(B)=0.3+0.5=0.8$. \textbf{Answer: (b).}
\item For independent events $A$ and $B$, $P(A\cap B)=P(A)\times P(B)=0.4\times 0.25=0.1$. \textbf{Answer: (b).}
\item Conditional probability: $P(A\mid B)=\dfrac{P(A\cap B)}{P(B)}=\dfrac{0.2}{0.4}=0.5$. \textbf{Answer: (b).}
\end{enumerate}
\end{corrige}

\begin{corrige}[Exercise 2 --- True or false]
\begin{enumerate}
\item \textbf{False.} If $A$ and $B$ are mutually exclusive with $P(A)>0$ and $P(B)>0$, then $P(A\cap B)=0$ whereas $P(A)P(B)>0$, so $P(A\cap B)\neq P(A)P(B)$: they are \emph{not} independent. Counter-example: throwing a fair die, $A=$ ``obtain $1$'', $B=$ ``obtain $2$'' are mutually exclusive but $P(A\cap B)=0\neq \tfrac{1}{36}=P(A)P(B)$.
\item \textbf{True.} This is the addition law. Proof: $A\cup B$ is the disjoint union of $A\setminus B$, $B\setminus A$ and $A\cap B$. Hence
\[P(A\cup B)=\bigl[P(A)-P(A\cap B)\bigr]+\bigl[P(B)-P(A\cap B)\bigr]+P(A\cap B)=P(A)+P(B)-P(A\cap B).\]
\item \textbf{True.} By definition $P(A\mid B)=\dfrac{P(A\cap B)}{P(B)}$ (with $P(B)>0$). If $P(A\mid B)=P(A)$, then $P(A\cap B)=P(A)P(B)$, which is exactly the definition of independence.
\item \textbf{True.} $A$ and $A'$ are mutually exclusive and exhaustive, so $P(A)+P(A')=P(A\cup A')=P(S)=1$.
\end{enumerate}
\end{corrige}

\begin{corrige}[Exercise 3 --- Definitions]
\begin{enumerate}
\item A \cle{random experiment} is a process whose outcome cannot be predicted with certainty in advance; its \cle{sample space} $S$ is the set of all possible outcomes. \emph{Example:} tossing a fair coin once; $S=\{H,T\}$.
\item Two events are \cle{mutually exclusive} if they cannot occur simultaneously, i.e. $A\cap B=\varnothing$. \emph{Example:} on a die, $A=$ ``even score'' and $B=$ ``score $5$''.
\item Events $A_1,\dots,A_n$ are \cle{exhaustive} if their union is the whole sample space: $A_1\cup\cdots\cup A_n=S$. \emph{Example:} on a die, ``even'' and ``odd'' are exhaustive.
\item Two events are \cle{independent} if the occurrence of one does not affect the probability of the other: $P(A\cap B)=P(A)P(B)$. \emph{Example:} the outcomes of two successive tosses of a fair coin.
\item The \cle{conditional probability} of $A$ given $B$ (with $P(B)>0$) is $P(A\mid B)=\dfrac{P(A\cap B)}{P(B)}$. \emph{Example:} for a fair die, $P(\text{score }2\mid\text{even})=\dfrac{1/6}{1/2}=\dfrac13$.
\end{enumerate}
\end{corrige}

\begin{corrige}[Exercise 4 --- Direct calculations]
\begin{enumerate}
\item The word PROBABILITY has $11$ letters: P, R, O, B, A, B, I, L, I, T, Y.
\begin{enumerate}
\item Vowels: O, A, I, I, i.e. $4$ letters. $P(\text{vowel})=\dfrac{4}{11}$.
\item The letter B occurs twice: $P(B)=\dfrac{2}{11}$.
\end{enumerate}
\item Sample space for two fair coins: $S=\{HH,\ HT,\ TH,\ TT\}$. Exactly one head corresponds to $\{HT,TH\}$, so $P=\dfrac{2}{4}=\dfrac{1}{2}$.
\item By the addition law:
\[P(A\cap B)=P(A)+P(B)-P(A\cup B)=0.6+0.5-0.8=0.3.\]
Then $P(A\mid B)=\dfrac{P(A\cap B)}{P(B)}=\dfrac{0.3}{0.5}=0.6$.
\end{enumerate}
\end{corrige}

\courssub{I practise}

\begin{corrige}[Exercise 5 --- Two dice]
The sample space contains $6\times 6=36$ equally likely ordered pairs.
\begin{enumerate}
\item Sum $7$: $(1,6),(2,5),(3,4),(4,3),(5,2),(6,1)$ --- $6$ outcomes. $P=\dfrac{6}{36}=\dfrac{1}{6}$.
\item Sum at least $10$: sums $10,11,12$. Sum $10$: $(4,6),(5,5),(6,4)$; sum $11$: $(5,6),(6,5)$; sum $12$: $(6,6)$. Total $6$ outcomes: $P=\dfrac{6}{36}=\dfrac{1}{6}$.
\item Doubles: $(1,1),\dots,(6,6)$, i.e. $6$ outcomes: $P=\dfrac{6}{36}=\dfrac{1}{6}$.
\item Let $B=$ ``at least one die shows $4$'': $6+6-1=11$ outcomes (subtracting the double $(4,4)$ counted twice). Among these, the sum is $7$ only for $(3,4)$ and $(4,3)$: $2$ outcomes. Hence
\[P(\text{sum }7\mid B)=\dfrac{2/36}{11/36}=\dfrac{2}{11}.\]
\end{enumerate}
\end{corrige}

\begin{corrige}[Exercise 6 --- Balls drawn without replacement]
A bag contains $4$ red and $6$ black balls. Two balls are drawn successively without replacement.
First draw: $P(R)=\dfrac{4}{10}=\dfrac25$, $P(B)=\dfrac{6}{10}=\dfrac35$.
Second draw (without replacement): after $R$, $3$ red and $6$ black remain $\Rightarrow P(R)=\dfrac39=\dfrac13$, $P(B)=\dfrac69=\dfrac23$; after $B$, $4$ red and $5$ black remain $\Rightarrow P(R)=\dfrac49$, $P(B)=\dfrac59$.
\begin{enumerate}
\item $P(RR)=\dfrac{4}{10}\times\dfrac{3}{9}=\dfrac{12}{90}=\dfrac{2}{15}$.
\item $P(\text{different colours})=P(RB)+P(BR)=\dfrac{4}{10}\times\dfrac{6}{9}+\dfrac{6}{10}\times\dfrac{4}{9}=\dfrac{24}{90}+\dfrac{24}{90}=\dfrac{48}{90}=\dfrac{8}{15}$.
\item Given the first ball was black, $9$ balls remain, of which $4$ are red:
\[P(\text{2nd red}\mid\text{1st black})=\dfrac{4}{9}.\]
\end{enumerate}
\end{corrige}

\begin{corrige}[Exercise 7 --- Choosing a committee]
A committee of $4$ is chosen from $6$ women and $5$ men.
Total number of committees: $\dbinom{11}{4}=\dfrac{11\cdot10\cdot9\cdot8}{24}=330$.
\begin{enumerate}
\item Exactly $2$ women (and $2$ men): $\dbinom{6}{2}\dbinom{5}{2}=15\times 10=150$ committees.
\[P=\dfrac{150}{330}=\dfrac{5}{11}.\]
\item At least one woman $=1-P(\text{no woman})$. No woman means $4$ men: $\dbinom{5}{4}=5$ committees.
\[P=1-\dfrac{5}{330}=1-\dfrac{1}{66}=\dfrac{65}{66}.\]
\item No man means $4$ women: $\dbinom{6}{4}=15$ committees.
\[P=\dfrac{15}{330}=\dfrac{1}{22}.\]
\end{enumerate}
\end{corrige}

\begin{corrige}[Exercise 8 --- Cards with and without replacement]
A standard deck of $52$ cards ($4$ suits, $13$ ranks, $4$ aces, $4$ kings, $13$ hearts).
\textbf{With replacement} (independent draws, $52$ cards each time):
\begin{enumerate}
\item $P(\text{two aces})=\dfrac{4}{52}\times\dfrac{4}{52}=\dfrac{1}{13}\times\dfrac{1}{13}=\dfrac{1}{169}$.
\item $P(\text{ace then king})=\dfrac{4}{52}\times\dfrac{4}{52}=\dfrac{1}{169}$.
\item $P(\text{at least one heart})=1-P(\text{no heart})=1-\left(\dfrac{39}{52}\right)^{2}=1-\dfrac{9}{16}=\dfrac{7}{16}$.
\end{enumerate}
\textbf{Without replacement:}
\begin{enumerate}
\item $P(\text{two aces})=\dfrac{4}{52}\times\dfrac{3}{51}=\dfrac{12}{2652}=\dfrac{1}{221}$.
\item[(c)] $P(\text{at least one heart})=1-\dfrac{39}{52}\times\dfrac{38}{51}=1-\dfrac{1482}{2652}=1-\dfrac{19}{34}=\dfrac{15}{34}$.
\end{enumerate}
\textbf{Comment:} without replacement the second draw is no longer independent of the first; the probability of two aces decreases (fewer aces remain after an ace), while the probability of at least one heart increases slightly (fewer non-hearts remain after a non-heart).
\end{corrige}

\courssub{I prepare for the GCE}

\begin{corrige}[Exercise 9 --- Subjects chosen by candidates]
Let $C=$ ``takes Chemistry'', $B=$ ``takes Biology''. Given: $P(C)=0.55$, $P(B)=0.40$, $P(C\cap B)=0.25$.
\begin{enumerate}
\item Venn diagram regions:
\begin{itemize}
\item $C$ only $=P(C)-P(C\cap B)=0.55-0.25=0.30$;
\item $B$ only $=P(B)-P(C\cap B)=0.40-0.25=0.15$;
\item Intersection $=0.25$;
\item Outside $=1-(0.30+0.25+0.15)=0.30$.
\end{itemize}
\begin{popfigure}
\begin{tikzpicture}
  \draw[popblue, thick] (0,0) circle (1.5);
  \draw[poporange, thick] (2.5,0) circle (1.5);
  \node[popblue] at (-1.2,0) {$C$};
  \node[poporange] at (3.7,0) {$B$};
  \node at (0,0) {$0.30$};
  \node at (2.5,0) {$0.15$};
  \node at (1.25,0) {$0.25$};
  \node at (1.25,-2) {$0.30$};
  \draw[popdark] (-3,-2.5) rectangle (5.5,2);
  \node[popdark, anchor=north west] at (-2.8,1.8) {$S$};
\end{tikzpicture}
\legende{Venn diagram for Chemistry and Biology}
\end{popfigure}
\emph{(2 marks: correct regions, correct outside value)}
\item
\begin{enumerate}
\item $P(C\cup B)=0.55+0.40-0.25=0.70$. \hfill \emph{(2 marks: formula, answer)}
\item Exactly one subject: $0.30+0.15=0.45$. \hfill \emph{(2 marks)}
\item Neither: $1-0.70=0.30$. \hfill \emph{(1 mark)}
\item $P(B\mid C)=\dfrac{P(B\cap C)}{P(C)}=\dfrac{0.25}{0.55}=\dfrac{5}{11}\approx 0.455$. \hfill \emph{(2 marks)}
\end{enumerate}
\item Independence test: $P(C)P(B)=0.55\times 0.40=0.22\neq 0.25=P(C\cap B)$. Hence $C$ and $B$ are \textbf{not} independent. \hfill \emph{(3 marks: product computed, comparison, conclusion)}
\end{enumerate}
\textbf{Examiner's expectations:} a clearly labelled two-circle Venn diagram with the four numerical regions; the addition law quoted before use; a numerical test (not a vague argument) for independence.
\end{corrige}

\begin{corrige}[Exercise 10 --- Defective items and Bayes' theorem]
Let $D=$ ``item is defective''. Machines $A$ and $B$ produce $70\%$ and $30\%$ of output respectively. Defective rates: $P(D\mid A)=0.02$, $P(D\mid B)=0.05$.
\begin{enumerate}
\item Tree diagram:
\begin{popfigure}
\begin{tikzpicture}[level distance=2.5cm, sibling distance=3cm]
  \node[popdark] {}
    child {node[popblue] {$A$ ($0.7$)}
      child {node[popgreen] {$D$ ($0.02$)} edge from parent node[left] {$0.02$}}
      child {node[popred] {$D'$ ($0.98$)} edge from parent node[right] {$0.98$}}
      edge from parent node[left] {$0.7$}
    }
    child {node[poporange] {$B$ ($0.3$)}
      child {node[popgreen] {$D$ ($0.05$)} edge from parent node[left] {$0.05$}}
      child {node[popred] {$D'$ ($0.95$)} edge from parent node[right] {$0.95$}}
      edge from parent node[right] {$0.3$}
    };
  % Terminal probabilities
  \node[popdark, anchor=north] at (-3.5,-5) {$P(A\cap D)=0.014$};
  \node[popdark, anchor=north] at (-0.5,-5) {$P(A\cap D')=0.686$};
  \node[popdark, anchor=north] at (2.5,-5) {$P(B\cap D)=0.015$};
  \node[popdark, anchor=north] at (5.5,-5) {$P(B\cap D')=0.285$};
\end{tikzpicture}
\legende{Tree diagram for machine and defect status}
\end{popfigure}
\emph{(3 marks: structure, branch probabilities, terminal probabilities)}
\item Law of total probability:
\[P(D)=P(A)P(D\mid A)+P(B)P(D\mid B)=0.7\times 0.02+0.3\times 0.05=0.014+0.015=0.029.\]
\emph{(3 marks: total probability formula, substitution, answer)}
\item Bayes' theorem:
\[P(B\mid D)=\dfrac{P(B)P(D\mid B)}{P(D)}=\dfrac{0.015}{0.029}=\dfrac{15}{29}\approx 0.517.\]
\emph{(4 marks: Bayes' formula stated, numerator, denominator from (b), simplified answer)}
\item Each item is defective with probability $p=0.029$, independently. For $3$ items, exactly one defective:
\[P=\binom{3}{1}(0.029)(0.971)^{2}=3\times 0.029\times 0.942841\approx 0.0820.\]
\emph{(3 marks: binomial structure with $\binom{3}{1}$, correct powers, numerical answer)}
\end{enumerate}
\textbf{Examiner's expectations:} the tree must carry probabilities on \emph{every} branch; Bayes' theorem must be written explicitly before substitution; answers left as exact fractions or given to at least $3$ significant figures.
\end{corrige}

\begin{corrige}[Exercise 11 --- Biased coin and GCE mock examination]
\textbf{Part A.} A biased coin has $P(H)=\dfrac23$, $P(T)=\dfrac13$. Three independent tosses.
\begin{enumerate}
\item Tree diagram with three levels:
\begin{popfigure}
\begin{tikzpicture}[level distance=2cm, sibling distance=2.5cm]
  \node[popdark] {}
    child {node[popblue] {$H$ ($\frac23$)}
      child {node[popblue] {$H$ ($\frac23$)}
        child {node[popgreen] {$H$ ($\frac23$)} edge from parent node[left] {$\frac23$}}
        child {node[popred] {$T$ ($\frac13$)} edge from parent node[right] {$\frac13$}}
        edge from parent node[left] {$\frac23$}
      }
      child {node[poporange] {$T$ ($\frac13$)}
        child {node[popgreen] {$H$ ($\frac23$)} edge from parent node[left] {$\frac23$}}
        child {node[popred] {$T$ ($\frac13$)} edge from parent node[right] {$\frac13$}}
        edge from parent node[right] {$\frac13$}
      }
      edge from parent node[left] {$\frac23$}
    }
    child {node[poporange] {$T$ ($\frac13$)}
      child {node[popblue] {$H$ ($\frac23$)}
        child {node[popgreen] {$H$ ($\frac23$)} edge from parent node[left] {$\frac23$}}
        child {node[popred] {$T$ ($\frac13$)} edge from parent node[right] {$\frac13$}}
        edge from parent node[left] {$\frac23$}
      }
      child {node[poporange] {$T$ ($\frac13$)}
        child {node[popgreen] {$H$ ($\frac23$)} edge from parent node[left] {$\frac23$}}
        child {node[popred] {$T$ ($\frac13$)} edge from parent node[right] {$\frac13$}}
        edge from parent node[right] {$\frac13$}
      }
      edge from parent node[right] {$\frac13$}
    };
\end{tikzpicture}
\legende{Tree diagram for three tosses of a biased coin}
\end{popfigure}
\emph{(3 marks)}
\item
\begin{enumerate}
\item Exactly two heads: paths $HHT$, $HTH$, $THH$, each of probability $\left(\tfrac23\right)^{2}\left(\tfrac13\right)=\tfrac{4}{27}$.
\[P=3\times\dfrac{4}{27}=\dfrac{12}{27}=\dfrac{4}{9}.\]
\emph{(2 marks)}
\item At least one head: $P=1-P(TTT)=1-\left(\tfrac13\right)^{3}=1-\dfrac{1}{27}=\dfrac{26}{27}$. \hfill \emph{(2 marks)}
\end{enumerate}
\item Let $E=$ ``exactly two heads'' and $F=$ ``first toss a head''. $E\cap F=\{HHT,HTH\}$ with probability $2\times\tfrac{4}{27}=\tfrac{8}{27}$.
\[P(F\mid E)=\dfrac{P(E\cap F)}{P(E)}=\dfrac{8/27}{12/27}=\dfrac{8}{12}=\dfrac{2}{3}.\]
\emph{(3 marks: conditional probability formula, correct events, answer)}
\end{enumerate}
\textbf{Part B.} Let $M=$ ``passes Mathematics'', $P=$ ``passes Physics'': $P(M)=0.6$, $P(P)=0.5$, $P(M\cap P)=0.3$.
\begin{enumerate}
\item
\begin{enumerate}
\item $P(M\cup P)=0.6+0.5-0.3=0.8$. \hfill \emph{(2 marks)}
\item Exactly one subject: $P(M\cup P)-P(M\cap P)=0.8-0.3=0.5$ (equivalently $0.3+0.2$). \hfill \emph{(2 marks)}
\item $P(P\mid M)=\dfrac{P(P\cap M)}{P(M)}=\dfrac{0.3}{0.6}=0.5$. \hfill \emph{(2 marks)}
\end{enumerate}
\item Test: $P(M)P(P)=0.6\times 0.5=0.3=P(M\cap P)$. The condition for independence is satisfied, so the two events \textbf{are} independent. (Equivalently, $P(P\mid M)=0.5=P(P)$.) \hfill \emph{(3 marks)}
\end{enumerate}
\textbf{Examiner's expectations:} in Part A, the complement rule for ``at least one'' earns full credit only if the complement event is named; in Part B, the independence conclusion must follow from an explicit numerical comparison, not from intuition.
\end{corrige}

\newpage

\coursec{Summary sheet}

\begin{popfigure}
\begin{tikzpicture}[mindmap, grow cyclic, every node/.style={concept}, concept color=popblue, text=white,
level 1/.append style={level distance=3.5cm, sibling angle=51.4},
level 2/.append style={level distance=2.2cm, sibling angle=40},
root concept/.append style={font=\large\bfseries}]
\node[root concept]{Elementary Probability}
child[concept color=poporange]{ node{Foundations}
  child{ node{Experiment, outcome, sample space $S$} }
  child{ node{Events: simple, compound, impossible, certain} }
}
child[concept color=popgreen]{ node{Two Approaches}
  child{ node{Classical: $P(A)=\frac{n(A)}{n(S)}$} }
  child{ node{Empirical: relative frequency} }
}
child[concept color=poppurple]{ node{Addition Law}
  child{ node{$P(A\cup B)=P(A)+P(B)-P(A\cap B)$} }
  child{ node{Exclusive: $P(A\cap B)=0$} }
}
child[concept color=popdark]{ node{Multiplication Law}
  child{ node{$P(A\cap B)=P(A)\,P(B\mid A)$} }
  child{ node{Independent: $P(B\mid A)=P(B)$} }
}
child[concept color=popgold]{ node{Tools}
  child{ node{Venn diagrams} }
  child{ node{Tree diagrams} }
  child{ node{Permutations \& combinations} }
}
child[concept color=poppink]{ node{Bayes' Theorem}
  child{ node{$P(A\mid B)=\dfrac{P(B\mid A)P(A)}{P(B)}$} }
}
child[concept color=popmuted]{ node{Applications}
  child{ node{Cards, dice, coins; real-life risk} }
};
\end{tikzpicture}
\legende{Mind map of the chapter Elementary Probability: the seven families of ideas and how they connect.}
\end{popfigure}

\begin{aretenir}[Chapter Summary --- Elementary Probability]
\textbf{Key definitions.}
\begin{itemize}
\item A \cle{random experiment} is a process whose outcome cannot be predicted with certainty before it is carried out (tossing a coin, rolling a die, drawing a card).
\item Each possible result of the experiment is an \cle{outcome}; the \cle{sample space} $S$ (also written $\Omega$ or $E$) is the set of \emph{all} possible outcomes.
\item An \cle{event} is a subset of $S$. It is a \emph{simple} (elementary) event if it contains exactly one outcome, \emph{compound} if it contains several, \emph{impossible} if it is $\varnothing$, and \emph{certain} if it equals $S$.
\item \cle{Classical (theoretical) probability}: when the outcomes of $S$ are equally likely,
\[ P(A)=\frac{n(A)}{n(S)}=\frac{\text{number of favourable outcomes}}{\text{total number of outcomes}}. \]
\item \cle{Empirical (experimental) probability}: $P(A)\approx \dfrac{\text{number of times }A\text{ occurs}}{\text{number of trials}}$; the estimate becomes more reliable as the number of trials grows large.
\item $A$ and $B$ are \cle{mutually exclusive} if $A\cap B=\varnothing$ (they cannot occur together); they are \cle{exhaustive} if $A\cup B=S$ (one of them must occur).
\item The \cle{complement} of $A$ is $A'$ (also written $\bar{A}$ or $A^c$), the event ``$A$ does not occur'', with $P(A')=1-P(A)$.
\item $A$ and $B$ are \cle{independent} if the occurrence of one does not affect the probability of the other: $P(A\cap B)=P(A)\,P(B)$, equivalently $P(B\mid A)=P(B)$ when $P(A)>0$.
\end{itemize}

\textbf{Laws and formulas to know.}
\begin{itemize}
\item Axioms: $0\le P(A)\le 1$ for every event $A$; $P(S)=1$; $P(\varnothing)=0$.
\item Complement rule: $P(A')=1-P(A)$.
\item Addition law: $P(A\cup B)=P(A)+P(B)-P(A\cap B)$; if $A$ and $B$ are mutually exclusive, $P(A\cup B)=P(A)+P(B)$.
\item Conditional probability: $P(B\mid A)=\dfrac{P(A\cap B)}{P(A)}$, provided $P(A)>0$.
\item Multiplication law: $P(A\cap B)=P(A)\,P(B\mid A)$; if $A$ and $B$ are independent, $P(A\cap B)=P(A)\,P(B)$.
\item Law of total probability: if $A_1,A_2,\dots,A_n$ form a partition of $S$ (mutually exclusive and exhaustive, each with non-zero probability), then for any event $B$,
\[ P(B)=\sum_{i=1}^{n} P(A_i)\,P(B\mid A_i). \]
\item Bayes' theorem: under the same conditions,
\[ P(A_i\mid B)=\frac{P(A_i)\,P(B\mid A_i)}{\displaystyle\sum_{j=1}^{n} P(A_j)\,P(B\mid A_j)}. \]
\item Counting techniques: use permutations ${}^nP_r=\dfrac{n!}{(n-r)!}$ when order matters, and combinations $\dbinom{n}{r}=\dfrac{n!}{r!\,(n-r)!}$ when order does not matter, to evaluate $n(A)$ and $n(S)$.
\end{itemize}

\textbf{Methods.}
\begin{itemize}
\item \emph{Venn diagrams}: represent $S$ by a rectangle and events by circles; shade regions for $A\cup B$, $A\cap B$ and complements; use $n(A\cup B)=n(A)+n(B)-n(A\cap B)$ to pass between counts and probabilities.
\item \emph{Tree diagrams}: multiply probabilities along the branches of a path, then add the probabilities of all paths leading to the desired event. The probabilities on the branches leaving any one node sum to $1$ --- use this to check your tree.
\item \emph{Cards, dice, coins}: an ordinary pack has $52$ cards ($4$ suits of $13$ ranks; $26$ red, $26$ black; $12$ picture cards, $4$ aces); two fair dice give $36$ equally likely ordered pairs --- display them in a $6\times 6$ table; coins give powers of $2$ outcomes ($2^n$ for $n$ tosses).
\item \emph{Without replacement}: the composition changes, so the second draw is conditional on the first; \emph{with replacement}: the stages are independent and probabilities stay constant.
\item \emph{``At least one'' problems}: pass to the complement, $P(\text{at least one})=1-P(\text{none})$.
\end{itemize}

\textbf{Common mistakes.}
\begin{itemize}
\item Adding $P(A)+P(B)$ when $A$ and $B$ are \emph{not} mutually exclusive --- this double-counts $A\cap B$; subtract $P(A\cap B)$.
\item Multiplying $P(A)\,P(B)$ when the events are \emph{not} independent --- use $P(A)\,P(B\mid A)$ instead.
\item Confusing \emph{mutually exclusive} with \emph{independent}: two mutually exclusive events with non-zero probabilities are \emph{never} independent, since $P(A\cap B)=0\ne P(A)P(B)$.
\item Giving a probability outside $[0,1]$: a negative answer or one greater than $1$ signals an error --- check your working.
\item Ignoring the words ``without replacement'': after the first draw there is one item fewer, and the probabilities change.
\item Mixing up $\dbinom{n}{r}$ and ${}^nP_r$: always ask first whether the order of selection matters.
\item In conditional probability, inverting the roles: $P(A\mid B)\ne P(B\mid A)$ in general.
\end{itemize}

\textbf{GCE tips.}
\begin{itemize}
\item Always state the sample space clearly and count outcomes carefully; for two dice, draw the $6\times 6$ table of outcomes to avoid omissions.
\item Write down the formula \emph{before} substituting numerical values --- method marks are awarded for correct statements of the laws.
\item Give final answers as simplified fractions unless the question asks for decimals or percentages.
\item For ``at least one'' questions, use the complement: it is almost always shorter.
\item In Bayes-type questions, draw the tree diagram first, compute $P(A\cap B)$ and $P(B)$ from it, then apply $P(A\mid B)=\dfrac{P(A\cap B)}{P(B)}$.
\item Finish by checking that every probability you quote lies between $0$ and $1$, and that the probabilities of a complete set of exhaustive events add up to $1$.
\end{itemize}
\end{aretenir}

\findecours

\end{document}
